\chapter{Conduction, convection, radiation}
\label{vol04:ch05}

\epigraph{Heat, like gravity, penetrates every substance of the
universe, its rays occupy all parts of space. The object of our work
is to set forth the mathematical laws which this element obeys.}{
Joseph Fourier, \emph{The Analytical Theory of Heat} (1822), in the
Freeman translation \cite[Preliminary Discourse, p.~1]{hist:fourier1822}.}

\chapmeta{Half-life: foundational. Archetypes: transport, scaling, failure. Exercise target: 32. Project track: household wall U-value measurement. Hazard class: low (winter outdoor measurements; standard precautions). Reader path default: standard, with sections explicitly tagged. Named cases: Three Mile Island (1979, cross-referenced for thermal-runaway discussion); Chernobyl (1986, cross-referenced for positive-void-coefficient instability); Bhopal (1984, cross-referenced for chemical-reactor runaway).}

\noindent
Heat transfer is the engineering of energy flow driven by temperature
difference. The three modes (conduction through matter at rest,
convection in moving fluids, radiation through electromagnetic
fields) are governed by different physics but share a common
engineering form: a flux proportional to a driving force, with a
geometric and material conductance setting the proportionality. The
chapter develops each mode, the combined-resistance machinery that
working engineers use, and the failure section closes with the
thermal-runaway pattern that destroys electronics, batteries,
reactors, and reactive chemical systems.

\section{Fourier's law of conduction}\pathtag{core}

A solid (or stationary fluid) with a temperature gradient conducts
heat at a rate set by \engterm{Fourier's law},
\[
\vect{q} = -k \grad T,
\]
where $\vect{q}$ is the heat flux (W/m$^2$), $k$ is the thermal
conductivity of the material, and $\grad T$ is the temperature
gradient \cite{hist:fourier1822}. The minus sign reflects that heat
flows from hot to cold. In one dimension,
\[
q = -k \frac{dT}{dx}.
\]
Thermal conductivity has units of W/(m$\cdot$K). Working values:

\begin{itemize}[leftmargin=*]
  \item Copper: $400\,\text{W}/(\text{m}\cdot\si{\kelvin})$.
  \item Aluminium: $237\,\text{W}/(\text{m}\cdot\si{\kelvin})$.
  \item Steel: $50\,\text{W}/(\text{m}\cdot\si{\kelvin})$.
  \item Stainless steel: $16\,\text{W}/(\text{m}\cdot\si{\kelvin})$.
  \item Glass: $1.0\,\text{W}/(\text{m}\cdot\si{\kelvin})$.
  \item Concrete: $1.7\,\text{W}/(\text{m}\cdot\si{\kelvin})$.
  \item Brick: $0.7\,\text{W}/(\text{m}\cdot\si{\kelvin})$.
  \item Wood (across grain): $0.15\,\text{W}/(\text{m}\cdot\si{\kelvin})$.
  \item Water: $0.6\,\text{W}/(\text{m}\cdot\si{\kelvin})$.
  \item Air (still): $0.025\,\text{W}/(\text{m}\cdot\si{\kelvin})$.
  \item Polyurethane foam: $0.025\,\text{W}/(\text{m}\cdot\si{\kelvin})$.
  \item Aerogel: $0.013\,\text{W}/(\text{m}\cdot\si{\kelvin})$.
\end{itemize}

The four orders of magnitude separating copper from polyurethane
foam are the engineering working range. Metals conduct heat through
their electron gas (the same charge carriers that conduct electricity,
giving the Wiedemann-Franz law $k/\sigma T \approx 2.45 \times
10^{-8}\,\text{W}\cdot\Omega/\si{\kelvin}^2$ at room temperature).
Insulators conduct through lattice phonons and gas-pore convection;
the lowest conductivities are achieved by minimising both contributions
(closed-cell foams, evacuated multilayer insulation, aerogels with
sub-mean-free-path pore size).

The four-order-of-magnitude span is not the limit of the table. A
diamond crystal conducts heat at $2000\,\text{W}/(\text{m}\cdot\si{
\kelvin})$, five times copper, because its stiff light lattice carries
phonons with little scattering, and that is why diamond and its
relatives appear as heat spreaders under high-power semiconductors. At
the other extreme, evacuated multilayer insulation reaches an effective
conductivity below $10^{-3}\,\text{W}/(\text{m}\cdot\si{\kelvin})$ by
removing the gas that carries pore conduction and stacking reflective
films that block radiation, a six-order span from diamond to the best
insulation. Conductivity is also temperature-dependent in ways the
working tables hide: a metal's conductivity falls slowly as temperature
rises and its electrons scatter more, while a gas's conductivity climbs
as the square root of temperature, so the still-air figure of $0.025$ is
a room-temperature value that an engineer sizing a high-temperature
furnace lining must correct upward. The single tabulated number is a
starting point; the property an engineer actually needs is the
conductivity at the operating temperature, read from a property table or
correlation rather than assumed constant across a wide range.

% Figure: the working span of thermal conductivity on a logarithmic axis,
% from the best insulators through nonmetals to metals and diamond. A
% horizontal log ruler with labelled markers, drawn from tabulated values,
% not measured. Shows the six-order span the chapter opens with.
\begin{figure}[htbp]
\centering
\begin{tikzpicture}[every node/.style={font=\small}]
\begin{axis}[
  width=13cm, height=4.2cm,
  xmode=log,
  xmin=0.008, xmax=3000,
  ymin=0, ymax=1,
  axis lines=left,
  ytick=\empty,
  xlabel={thermal conductivity $k$ (W/(m$\cdot$K)), log scale},
  tick label style={font=\scriptsize},
  label style={font=\small},
  clip=false,
]
% markers along the axis
\addplot[only marks, mark=*, engred, mark size=1.6pt] coordinates {
  (0.013,0.5) (0.025,0.5) (0.15,0.5) (0.6,0.5) (1.0,0.5)
  (16,0.5) (50,0.5) (237,0.5) (400,0.5) (2000,0.5)};
\node[font=\scriptsize, engred, rotate=55, anchor=west] at (axis cs:0.013,0.55) {aerogel};
\node[font=\scriptsize, enggray, rotate=55, anchor=west] at (axis cs:0.025,0.55) {foam, still air};
\node[font=\scriptsize, enggray, rotate=55, anchor=west] at (axis cs:0.15,0.55) {wood};
\node[font=\scriptsize, enggray, rotate=55, anchor=west] at (axis cs:0.6,0.55) {water};
\node[font=\scriptsize, enggray, rotate=55, anchor=west] at (axis cs:1.0,0.55) {glass};
\node[font=\scriptsize, enggray, rotate=55, anchor=west] at (axis cs:16,0.55) {stainless};
\node[font=\scriptsize, enggray, rotate=55, anchor=west] at (axis cs:50,0.55) {steel};
\node[font=\scriptsize, engblue, rotate=55, anchor=west] at (axis cs:237,0.55) {aluminium};
\node[font=\scriptsize, engblue, rotate=55, anchor=west] at (axis cs:400,0.55) {copper};
\node[font=\scriptsize, engblue, rotate=55, anchor=west] at (axis cs:2000,0.55) {diamond};
% band annotations below the axis
\node[font=\scriptsize, engred, anchor=north] at (axis cs:0.03,0.30) {insulators};
\node[font=\scriptsize, enggray, anchor=north] at (axis cs:1.0,0.30) {nonmetals};
\node[font=\scriptsize, engblue, anchor=north] at (axis cs:300,0.30) {metals};
\end{axis}
\end{tikzpicture}
\caption[The working span of thermal conductivity]{The thermal
conductivities an engineer meets, laid on a logarithmic ruler. The best
insulators sit near $0.01\,$W/(m$\cdot$K), the nonmetals and stagnant
fluids cluster near unity, the structural metals run from tens to a few
hundred, and diamond tops the scale above copper. Nearly six decades
separate the ends, which is why the choice of material, far more than the
choice of thickness, sets whether a path conducts or insulates. The
grouping is physical: insulators carry heat through slow lattice phonons
and trapped gas, metals carry it through their fast electron gas, and
diamond carries it through a stiff light lattice that scatters phonons
little.}
\label{fig:vol04ch05:conductivity-spectrum}
\end{figure}


Figure~\ref{fig:vol04ch05:conductivity-spectrum} lays the working values on
a single logarithmic ruler, and the shape of that ruler is the first thing
to read from it. The metals crowd into the top decade, all within a factor
of ten of one another, because they share the same electron-gas mechanism
and differ only in how cleanly that gas moves. The insulators crowd into
the bottom decade for the opposite reason: each has driven both its solid
conduction and its pore conduction as low as the material allows, so they
converge from below toward the still-air floor near $0.025$. The wide
empty middle is the nonmetal band, where a single mechanism, lattice
conduction through a disordered solid or a liquid, gives values near unity.
The practical reading is that the material choice, not the thickness, sets
the decade a path lives in: no achievable thickness of copper insulates,
and no achievable thinness of foam conducts, because four to six decades of
conductivity swamp the one or two decades of thickness an engineer can
reach. A design that needs to move heat picks from the top of the ruler and
then works on area and film; a design that needs to block heat picks from
the bottom and then works on thickness. The ruler also warns against the
common category error of treating a good structural metal as a thermal
insulator or a foam as a structural member: the properties that make a
material good at one job place it at the wrong end of this scale for the
other.

For a flat wall of thickness $L$ and area $A$ with temperatures $T_1$
and $T_2$ on the two faces and constant $k$, the steady-state heat
transfer rate is
\[
\dot{Q} = \frac{k A (T_1 - T_2)}{L} = \frac{T_1 - T_2}{R_{\text{cond}}},
\quad R_{\text{cond}} = \frac{L}{k A}.
\]
The form is analogous to Ohm's law for electrical conductors with
$\dot{Q}$ playing the role of current, $\Delta T$ the role of
voltage, and $R_{\text{cond}}$ the thermal resistance. Series and
parallel combinations of thermal resistances follow the same rules
as electrical resistors:
\[
R_{\text{series}} = R_1 + R_2 + \ldots, \quad R_{\text{parallel}} =
\left(\sum_i \frac{1}{R_i}\right)^{-1}.
\]
For cylindrical conduction (a pipe wall of inner radius $r_1$, outer
radius $r_2$, length $L$),
\[
R_{\text{cyl}} = \frac{\ln(r_2/r_1)}{2\pi k L},
\]
and for spherical conduction (a sphere of inner radius $r_1$, outer
radius $r_2$),
\[
R_{\text{sph}} = \frac{1}{4\pi k}\left(\frac{1}{r_1} - \frac{1}{r_2}
\right).
\]
The three resistance formulas are the working tool for steady
conduction in the canonical geometries; for more complex shapes,
shape factors or numerical methods are used.

% Figure: steady temperature profile through a composite wall.
% Three solid layers (brick, insulation, gypsum) plus two air films;
% the temperature drops linearly within each conductive layer and the
% slopes scale inversely with conductivity, so the insulation carries
% the largest drop. Bulk fluid temperatures shown at the two ends.
\begin{figure}[htbp]
\centering
\begin{tikzpicture}[scale=1.0, >=latex]
% layer boundaries on x, temperature on y (deg C)
% outdoor film | brick | insulation | gypsum | indoor film
% x positions
\def\xa{0}    % outdoor air
\def\xb{1.2}  % outer brick face
\def\xc{3.0}  % brick/insulation
\def\xd{6.4}  % insulation/gypsum
\def\xe{7.2}  % gypsum/indoor face
\def\xf{8.4}  % indoor air
% axes
\draw[->] (-0.3,0) -- (9.0,0) node[right, font=\small] {position};
\draw[->] (-0.3,0) -- (-0.3,5.2) node[above, font=\small] {$T$ ($^{\circ}\mathrm{C}$)};
% layer shading bands
\fill[engblue!8]  (\xb,0) rectangle (\xc,5.0);
\fill[englime!12] (\xc,0) rectangle (\xd,5.0);
\fill[engamber!12] (\xd,0) rectangle (\xe,5.0);
% layer separators
\foreach \x in {\xb,\xc,\xd,\xe}{ \draw[enggray, thin] (\x,0) -- (\x,5.0); }
% temperature node values (deg C), illustrative for 22 indoor / -15 outdoor
% Tout_air=-15, Tout_face=-13.5, T_brick/ins=-12, T_ins/gyp=18.5, T_gyp_face=20.5, Tin_air=22
% map T in [-15,22] to y in [0,5]: y = (T+15)/37*5
\def\Tya{0.0}   % -15
\def\Tyb{0.203} % -13.5
\def\Tyc{0.405} % -12
\def\Tyd{4.527} % 18.5
\def\Tye{4.797} % 20.5
\def\Tyf{5.0}   % 22
% profile line
\draw[engred, very thick]
  (\xa,\Tya) -- (\xb,\Tyb)   % outdoor film
  -- (\xc,\Tyc)              % brick (gentle slope)
  -- (\xd,\Tyd)              % insulation (steep slope, large drop)
  -- (\xe,\Tye)              % gypsum
  -- (\xf,\Tyf);             % indoor film
% endpoint markers
\fill[engred] (\xa,\Tya) circle (1.6pt);
\fill[engred] (\xf,\Tyf) circle (1.6pt);
% labels
\node[font=\scriptsize, rotate=90] at ({(\xb+\xc)/2},2.5) {brick};
\node[font=\scriptsize, rotate=90] at ({(\xc+\xd)/2},2.5) {insulation};
\node[font=\scriptsize, rotate=90] at ({(\xd+\xe)/2},2.5) {gypsum};
\node[font=\scriptsize, enggray] at (0.5,3.4) {outdoor};
\node[font=\scriptsize, enggray] at (0.5,3.0) {film};
\node[font=\scriptsize, enggray] at (7.8,1.6) {indoor};
\node[font=\scriptsize, enggray] at (7.8,1.2) {film};
\node[font=\scriptsize, anchor=east] at (-0.32,\Tya) {$-15$};
\node[font=\scriptsize, anchor=east] at (-0.32,\Tyf) {$22$};
\node[font=\scriptsize, anchor=south] at (\xa,\Tya) {$T_{\infty,o}$};
\node[font=\scriptsize, anchor=south] at (\xf,\Tyf) {$T_{\infty,i}$};
\end{tikzpicture}
\caption[Steady temperature profile through a composite wall]{Steady
temperature profile through a composite wall separating outdoor air at
$-15\,^{\circ}\mathrm{C}$ from indoor air at $22\,^{\circ}\mathrm{C}$.
The temperature falls linearly within each conductive layer, and the
slope in a layer is inversely proportional to its conductivity, so the
low-conductivity insulation carries the largest temperature drop while
the high-conductivity brick carries almost none. The two air films at
the surfaces add convective resistance and account for the temperature
jumps between the bulk air and the wall faces. The total heat flux is
the same through every layer in series, which is what makes the
resistance-network treatment of section 5.1 work.}
\label{fig:vol04ch05:wall-temperature-profile}
\end{figure}


% Figure: thermal-resistance network for the composite wall, drawn as an
% electrical circuit. Series chain of five resistors between the two bulk
% temperatures, with one parallel branch showing a thermal bridge (framing
% stud) across the insulation layer.
\begin{figure}[htbp]
\centering
\begin{tikzpicture}[scale=1.0, >=latex,
  res/.style={draw, thick, minimum width=1.1cm, minimum height=0.42cm, font=\scriptsize, fill=white}]
% baseline y
\def\y{0}
% nodes (temperatures) along the chain
\node[font=\small] (Ti) at (0,\y) {$T_{\infty,i}$};
\coordinate (n1) at (1.4,\y);
\coordinate (n2) at (3.4,\y);
\coordinate (n3) at (5.4,\y);
\coordinate (n4) at (7.4,\y);
\coordinate (n5) at (9.4,\y);
\node[font=\small] (To) at (10.8,\y) {$T_{\infty,o}$};
% wires
\draw[thick] (Ti) -- (n1);
\node[res] at ($(n1)!0.5!(n2)$) {$R_{h,i}$};
\draw[thick] (n1) -- ($(n1)!0.18!(n2)$);
\draw[thick] ($(n1)!0.82!(n2)$) -- (n2);
\node[res] at ($(n2)!0.5!(n3)$) {$R_{\text{gyp}}$};
\draw[thick] (n2) -- ($(n2)!0.18!(n3)$);
\draw[thick] ($(n2)!0.82!(n3)$) -- (n3);
% parallel section: insulation field vs framing bridge
\node[res] (Rins) at ($(n3)!0.5!(n4)$) {$R_{\text{ins}}$};
\draw[thick] (n3) -- ($(n3)!0.18!(n4)$);
\draw[thick] ($(n3)!0.82!(n4)$) -- (n4);
% upper parallel branch (thermal bridge)
\coordinate (b3) at (5.4,1.2);
\coordinate (b4) at (7.4,1.2);
\draw[thick] (n3) -- (b3);
\draw[thick] (n4) -- (b4);
\node[res, engred, draw=engred] at ($(b3)!0.5!(b4)$) {$R_{\text{stud}}$};
\draw[thick] (b3) -- ($(b3)!0.18!(b4)$);
\draw[thick] ($(b3)!0.82!(b4)$) -- (b4);
\node[font=\scriptsize, engred] at (6.4,1.7) {thermal bridge (parallel)};
% continue series
\node[res] at ($(n4)!0.5!(n5)$) {$R_{\text{brick}}$};
\draw[thick] (n4) -- ($(n4)!0.18!(n5)$);
\draw[thick] ($(n4)!0.82!(n5)$) -- (n5);
\draw[thick] (n5) -- (To);
% node temperature dots
\foreach \p in {n1,n2,n3,n4,n5}{ \fill (\p) circle (1.3pt); }
% heat flow arrow
\draw[->, engblue, very thick] (1.4,-0.9) -- (9.4,-0.9)
  node[midway, below, font=\scriptsize] {$\dot{Q} = (T_{\infty,i}-T_{\infty,o})/R_{\text{tot}}$};
\end{tikzpicture}
\caption[Thermal-resistance network for a composite wall]{The composite
wall of Figure~\ref{fig:vol04ch05:wall-temperature-profile} drawn as a
thermal-resistance network. Heat flows from the warm indoor air through
the indoor film resistance $R_{h,i}$, the gypsum, the insulation, the
brick, and out through the outdoor film, with the resistances in series.
The driving temperature difference divided by the total series resistance
gives the heat-flow rate, exactly as voltage divided by resistance gives
current. A framing stud crossing the insulation provides a parallel
low-resistance path $R_{\text{stud}}$, the thermal bridge that the
infrared survey in the diagnosis exercises detects as a cold stripe on
the interior surface.}
\label{fig:vol04ch05:thermal-resistance-circuit}
\end{figure}


Figure~\ref{fig:vol04ch05:wall-temperature-profile} shows the steady
temperature profile through a composite wall, and
Figure~\ref{fig:vol04ch05:thermal-resistance-circuit} draws the same
wall as the series-resistance network the calculation rests on. The
\texttt{thermal\_resistance.py} module accompanying this chapter builds
the network for a plane wall or an insulated tube and returns the
U-value and the interface temperatures; representative conductivities
for the calculation are tabulated in
\texttt{material\_thermal\_properties.csv}.

\engterm{A parallel path through a framed wall}. The series formula
imagines heat flowing straight through stacked layers, but a real wall
also carries parallel paths that let heat detour around the insulation, and
combining the two requires both rules at once. We take one square metre of
a timber-framed wall whose studs occupy a fraction $f = 0.15$ of the face
area (the framing fraction of a wood-stud wall) and whose cavity carries
the remaining $0.85$. The cavity is $140\,\si{\milli\meter}$ of mineral
wool ($k = 0.038\,\text{W}/(\text{m}\cdot\si{\kelvin})$); the studs are
$140\,\si{\milli\meter}$ of softwood ($k = 0.13\,\text{W}/(\text{m}\cdot
\si{\kelvin})$) spanning the same depth. Ignoring the shared linings for
the moment, the two paths run side by side from the inner face to the
outer, so their conductances add. The cavity path has area $0.85\,\text{m}^2$
and resistance $L/(kA) = 0.140/(0.038\times0.85) = 4.33\,\si{\kelvin}/
\text{W}$; the stud path has area $0.15\,\text{m}^2$ and resistance
$0.140/(0.13\times0.15) = 7.18\,\si{\kelvin}/\text{W}$. The conductances
are the reciprocals, $0.231$ and $0.139\,\text{W}/\si{\kelvin}$, so the
combined conductance is $0.370\,\text{W}/\si{\kelvin}$ and the effective
resistance of the parallel pair is $1/0.370 = 2.70\,\si{\kelvin}/\text{W}$.
The instructive comparison is against the cavity alone: a clear-field
calculation that saw only the insulation would have taken the whole
$1\,\text{m}^2$ at $0.140/0.038 = 3.68\,\si{\kelvin}\cdot\text{m}^2/
\text{W}$, a conductance of $0.272\,\text{W}/\si{\kelvin}$, understating the
loss by the $0.098\,\text{W}/\si{\kelvin}$ the studs carry, a $36\,\%$
error on this stud-heavy wall. The studs occupy $15\,\%$ of the area but
carry $38\,\%$ of the conductance, because their conductivity is more than
three times the wool's, and this is the arithmetic behind the
thermal-bridge uplift the house-retrofit case applies as a lump. The two
rules combine cleanly only when the paths are truly independent from face
to face; where the paths merge partway (a stud that stops short of the
outer sheathing), the honest model is a two-dimensional field solution, and
the parallel-resistance estimate becomes a bound rather than an answer.

\section{Steady and transient conduction}\pathtag{standard}

\engterm{The heat equation}. In a region with no internal heat
generation, conservation of energy combined with Fourier's law gives
the \engterm{heat equation}
\[
\frac{\partial T}{\partial t} = \alpha \laplacian T,
\]
where $\alpha = k/(\rho c_p)$ is the \engterm{thermal diffusivity},
with $\rho$ density and $c_p$ specific heat \cite{text:incropera2007}.
Thermal diffusivity has units of m$^2$/s. Working values: aluminium
$9.7 \times 10^{-5}\,\text{m}^2/\text{s}$, steel $1.2 \times 10^{-5}\,
\text{m}^2/\text{s}$, glass $5 \times 10^{-7}\,\text{m}^2/\text{s}$,
water $1.5 \times 10^{-7}\,\text{m}^2/\text{s}$, polyurethane foam
$3 \times 10^{-7}\,\text{m}^2/\text{s}$.

When internal heat generation is present at volumetric rate $\dot{q}_v$
(W/m$^3$), the heat equation reads
\[
\frac{\partial T}{\partial t} = \alpha \laplacian T + \frac{\dot{q}_v}{
\rho c_p}.
\]
The generation term enters for resistive electrical heating, chemical
reaction, nuclear heat in reactor fuel, and viscous dissipation in
flowing fluids.

\engterm{The heat equation from an energy balance}. The equation is not a
postulate; it is conservation of energy applied to a differential element with
Fourier's law substituted for the flux. Take the one-dimensional element of
cross-section $A$ and thickness $dx$ in
Figure~\ref{fig:vol04ch05:control-volume}. Energy conservation on the element
balances four terms over a time $dt$: the heat conducted in at the left face,
the heat conducted out at the right face, the heat generated inside, and the
change in stored internal energy,
\[
\dot{Q}_x - \dot{Q}_{x+dx} + \dot{q}_v\,A\,dx = \rho\,A\,dx\,c_p
\frac{\partial T}{\partial t}.
\]
The outgoing conducted rate differs from the incoming by a Taylor step,
$\dot{Q}_{x+dx} = \dot{Q}_x + (\partial \dot{Q}/\partial x)\,dx$, so the first
two terms collapse to $-(\partial \dot{Q}/\partial x)\,dx$. Fourier's law sets
the conducted rate, $\dot{Q} = -kA\,\partial T/\partial x$, and for constant
conductivity $\partial \dot{Q}/\partial x = -kA\,\partial^2 T/\partial x^2$.
Substituting and dividing through by $\rho A c_p\,dx$ leaves
\[
\frac{\partial T}{\partial t} = \frac{k}{\rho c_p}
\frac{\partial^2 T}{\partial x^2} + \frac{\dot{q}_v}{\rho c_p} = \alpha
\frac{\partial^2 T}{\partial x^2} + \frac{\dot{q}_v}{\rho c_p},
\]
the one-dimensional heat equation, with the three-dimensional form following by
applying the same balance to each coordinate and summing the second derivatives
into the Laplacian. The second spatial derivative is the signature of the
derivation: stored energy responds to the \emph{difference} between the flux
entering and the flux leaving, which is the gradient of a gradient, not the
gradient itself. A region where the temperature profile is straight
($\partial^2 T/\partial x^2 = 0$) conducts the same rate in as out and neither
heats nor cools; curvature in the profile is what drives the transient. The
thermal diffusivity $\alpha = k/(\rho c_p)$ that emerges is the ratio of how
fast a material conducts heat to how much heat it must store to change
temperature, and it is the property that sets the speed of every conduction
transient: a high-$\alpha$ metal equilibrates fast, a low-$\alpha$ insulator or
masonry slowly, even when their steady conductances differ in the opposite
direction.

% Figure: the differential control volume from which the heat equation is
% derived. A slab element of thickness dx carries a conducted flux q(x) in at
% the left face and q(x+dx) out at the right, with internal generation and a
% stored-energy term accounting for the imbalance. The Taylor expansion of the
% flux is the source of the second-derivative term in the heat equation.
\begin{figure}[htbp]
\centering
\begin{tikzpicture}[every node/.style={font=\small}]
% the control volume
\draw[fill=engblue!8, draw=engblue, thick] (0,0) rectangle (3.2,2.4);
% left and right face markers
\draw[engred, thick] (0,-0.05) -- (0,2.45);
\draw[engred, thick] (3.2,-0.05) -- (3.2,2.45);
% incoming flux arrow
\draw[engred, -{Stealth}, very thick] (-1.8,1.2) -- (-0.05,1.2);
\node[engred, font=\scriptsize, anchor=south] at (-0.9,1.25) {$\dot{Q}_x$};
% outgoing flux arrow
\draw[engred, -{Stealth}, very thick] (3.25,1.2) -- (5.0,1.2);
\node[engred, font=\scriptsize, anchor=south] at (4.1,1.25) {$\dot{Q}_{x+dx}$};
% generation
\node[engorange, font=\scriptsize] at (1.6,1.7) {$\dot{q}_v\,A\,dx$};
\draw[engorange, -{Stealth}] (1.6,1.5) -- (1.6,1.05);
% storage
\node[enggray, font=\scriptsize] at (1.6,0.6) {$\rho c_p A\,dx\,\dfrac{\partial T}{\partial t}$};
% dx dimension
\draw[<->, enggray] (0,-0.45) -- (3.2,-0.45);
\node[enggray, font=\scriptsize, anchor=north] at (1.6,-0.5) {$dx$};
% area label
\node[engblue, font=\scriptsize, anchor=south] at (1.6,2.5) {cross-section $A$};
% x axis
\draw[-{Stealth}] (-1.8,-0.9) -- (5.0,-0.9);
\node[anchor=north] at (5.0,-0.95) {$x$};
\end{tikzpicture}
\caption[The differential control volume for the heat equation]{The
differential control volume from which the heat equation follows. Heat is
conducted in at the left face at rate $\dot{Q}_x$ and out at the right face at
rate $\dot{Q}_{x+dx}$, generation adds $\dot{q}_v A\,dx$ inside, and the
imbalance is stored as the rate of change of internal energy $\rho c_p A\,dx\,
\partial T/\partial t$. Expanding the outgoing flux as $\dot{Q}_{x+dx} =
\dot{Q}_x + (\partial \dot{Q}/\partial x)\,dx$ and substituting Fourier's law
$\dot{Q} = -kA\,\partial T/\partial x$ turns the balance into the
one-dimensional heat equation. The second spatial derivative enters because the
stored energy is set by the net flux difference across the element, not by the
flux itself.}
\label{fig:vol04ch05:control-volume}
\end{figure}


\engterm{Steady-state solutions}. With $\partial T/\partial t = 0$,
the heat equation becomes Laplace's equation $\laplacian T = 0$ (no
generation) or Poisson's equation $\laplacian T = -\dot{q}_v/k$ (with
generation). The solutions in canonical geometries are standard and
give the resistance formulas of section 5.1. For arbitrary geometries,
the standard tools are conformal mapping (in two dimensions for
specific shapes), the shape-factor catalogue (for problems where two
isothermal surfaces are connected by a conducting medium), and
finite-element numerical solution.

Two structural properties of the steady field make it far more tractable
than the transient. The first is linearity: because Laplace's equation is
linear in $T$, any sum of solutions is again a solution, so a complicated
boundary condition can be split into pieces whose responses add. A room
corner heated on two faces is the superposition of the same corner heated on
each face alone, and a wall carrying both a temperature difference and an
embedded heated pipe is the sum of the plain-wall field and the pipe's field.
The superposition principle is what lets the shape-factor catalogue exist at
all: each entry is a single isothermal-pair field, and a real problem is
assembled by adding the entries that apply. The shape factor $S$ itself is
defined so that $\dot{Q} = S k\,\Delta T$ carries the whole geometry in one
number with the dimensions of length, and it depends only on the shape, never
on the conductivity or the temperature, which is why the same $S$ serves a
copper part and a concrete one. The second property is the maximum principle:
a solution of Laplace's equation takes its largest and smallest values on the
boundary, never at an interior point, so a body with no internal generation
can hold no interior hot spot hotter than its hottest surface. The moment an
interior peak appears, generation is present and the field obeys Poisson's
equation instead, whose interior maximum is exactly the centreline rise the
busbar and fuel-pin examples compute. The two principles together tell an
engineer where to look: for a passive conducting shape the extreme
temperatures live on the boundary and the shape factor sizes the flow, while
an interior peak is the signature of a source that Poisson's equation, not
Laplace's, must carry.

\engterm{Transient conduction: lumped capacitance}. When the
conduction within a body is fast compared to convection from its
surface, the body's temperature can be treated as spatially uniform
and only time-dependent. The criterion is the \engterm{Biot number}
\[
\text{Bi} = \frac{h L_c}{k_s} \ll 1,
\]
where $h$ is the surface convection coefficient, $L_c$ is the body's
characteristic length (volume divided by surface area), and $k_s$
is the body's thermal conductivity. For $\text{Bi} < 0.1$, the
lumped-capacitance approximation gives the body's temperature
$T(t)$ from
\[
\rho V c_p \frac{dT}{dt} = -h A (T - T_\infty),
\]
which integrates to
\[
T(t) - T_\infty = (T_0 - T_\infty) e^{-t/\tau}, \quad \tau =
\frac{\rho V c_p}{h A}.
\]
The time constant $\tau$ is the body's thermal response time. A
small piece of metal ($L_c \sim 1\,\si{\centi\meter}$) in air ($h
\sim 25\,\text{W}/(\text{m}^2\cdot\si{\kelvin})$) has $\tau$ of
order minutes. A large concrete wall ($L_c \sim 0.2\,\text{m}$) in
the same air has $\text{Bi} \gg 1$ and the lumped approximation
fails; the wall has spatial temperature gradients that take hours
to relax.

% Figure: lumped-capacitance transient cooling. Excess temperature decays
% exponentially toward ambient; two bodies with different time constants
% shown, plus the one-time-constant and the 63 percent markers.
\begin{figure}[htbp]
\centering
\begin{tikzpicture}
\begin{axis}[
  width=12cm, height=7.5cm,
  xlabel={time $t/\tau$ (in units of the time constant)},
  ylabel={excess temperature $(T-T_\infty)/(T_0-T_\infty)$},
  xmin=0, xmax=5, ymin=0, ymax=1.05,
  axis lines=left,
  legend pos=north east,
  legend cell align=left,
  tick label style={font=\scriptsize},
  label style={font=\small},
  legend style={font=\scriptsize},
  samples=120, domain=0:5,
  clip=false,
]
% the universal decay curve
\addplot[engblue, very thick] {exp(-x)};
\addlegendentry{$e^{-t/\tau}$ (lumped body)}
% 1/e marker at t=tau
\draw[enggray, dashed] (axis cs:1,0) -- (axis cs:1,0.3679);
\draw[enggray, dashed] (axis cs:0,0.3679) -- (axis cs:1,0.3679);
\node[font=\scriptsize, enggray, anchor=west] at (axis cs:1.05,0.42)
  {at $t=\tau$: $37\%$ remains};
\fill[engblue] (axis cs:1,0.3679) circle (1.6pt);
% 5 percent (3 time constants)
\draw[enggray, dashed] (axis cs:3,0) -- (axis cs:3,0.0498);
\node[font=\scriptsize, enggray, anchor=west] at (axis cs:3.05,0.13)
  {$t=3\tau$: $5\%$};
\fill[engblue] (axis cs:3,0.0498) circle (1.4pt);
\end{axis}
\end{tikzpicture}
\caption[Lumped-capacitance transient cooling]{Transient cooling of a
body small enough that its internal conduction outruns its surface
convection (Biot number below $0.1$), so the temperature is uniform and
decays exponentially toward the ambient value. Plotted against time in
units of the time constant $\tau=\rho V c_p/(hA)$, every such body
follows the same curve: after one time constant $37\%$ of the initial
temperature excess remains, after three time constants only $5\%$. The
time constant scales with the body's thermal mass $\rho V c_p$ and
inversely with the surface conductance $hA$, so a small well-ventilated
part settles in seconds while a large still-air part takes hours.}
\label{fig:vol04ch05:transient-cooling}
\end{figure}


Figure~\ref{fig:vol04ch05:transient-cooling} plots the universal
exponential decay that every lumped body follows when scaled by its own
time constant.

\engterm{Transient conduction: semi-infinite solid}. For a body
whose thickness is much larger than the thermal penetration depth on
the timescale of interest, the surface temperature determines a
diffusion-like temperature profile that penetrates with depth $\delta
\sim \sqrt{\alpha t}$. For a step change in surface temperature from
$T_0$ to $T_s$ at $t = 0$,
\[
\frac{T(x, t) - T_s}{T_0 - T_s} = \text{erf}\left(\frac{x}{2\sqrt{
\alpha t}}\right),
\]
with $\text{erf}$ the error function. The heat flux at the surface
is $q_s(t) = (T_s - T_0) \sqrt{k \rho c_p/(\pi t)}$, decreasing as
$t^{-1/2}$. The semi-infinite-solid result governs the response of
buildings to outdoor temperature swings (the daily and seasonal
penetration depths in concrete are $\sim 0.1$ and $\sim 1\,\text{m}$
respectively), the response of soil to surface heating, and the
temperature transient in any large body after a sudden surface
condition change.

\engterm{Transient conduction: shape factors and Heisler charts}.
For intermediate cases (a finite body with intermediate Biot number),
the heat equation must be solved in full. Standard solutions for the
slab, the long cylinder, and the sphere, parameterised by Biot
number and dimensionless time (Fourier number $\text{Fo} = \alpha
t/L_c^2$), are tabulated in Heisler charts. For engineering design,
the charts have largely been replaced by direct numerical solution
in finite-element packages, but the underlying scaling (Biot and
Fourier numbers) is the working language.

\begin{mastery}
The lumped-capacitance model and the $\text{Bi} < 0.1$ rule are not
separate facts; the rule is the lumped model's own statement of when it
may be trusted, and the one-term series solution shows where the bound
comes from. Take a plane wall of half-thickness $L$, convecting from
both faces into a fluid at $T_\infty$. The transient heat equation
$\partial T/\partial t = \alpha\,\partial^2 T/\partial x^2$ with the
convective boundary condition $-k\,\partial T/\partial x = h(T - T_\infty)$
at $x = L$ separates into modes $\theta = \sum_n C_n \cos(\lambda_n x/L)
\exp(-\lambda_n^2 \text{Fo})$, where $\theta = T - T_\infty$, the Fourier
number is $\text{Fo} = \alpha t/L^2$, and the eigenvalues solve the
transcendental relation $\lambda_n \tan\lambda_n = \text{Bi}$. After the
first few diffusion times the higher modes die and the leading term
governs, $\theta/\theta_0 \approx C_1 \cos(\lambda_1 x/L)
\exp(-\lambda_1^2 \text{Fo})$. For small Biot number the smallest root
expands as $\lambda_1^2 \approx \text{Bi}(1 - \text{Bi}/3 + \cdots)$, so
the centre cooling exponent $\lambda_1^2 \text{Fo} = \lambda_1^2 \alpha
t/L^2$ becomes $\text{Bi}\,\alpha t/L^2 = h t/(\rho c_p L)$, which is
exactly $t/\tau$ for the lumped time constant $\tau = \rho c_p L/h$ with
$L_c = L$. The lumped model is therefore the leading term of the exact
series in the limit of small Biot number, and the spatial
non-uniformity it discards is the cosine factor: the surface-to-centre
temperature difference is of order $\lambda_1^2/2 \approx \text{Bi}/2$
times the centre excess. Holding that error under a few percent is what
fixes the threshold at $\text{Bi} < 0.1$. The same expansion explains why
the steel-shaft quench at $\text{Bi} = 0.14$ ran slightly fast: the
$\text{Bi}/3$ correction to $\lambda_1^2$ lengthens the true cooling
exponent below the lumped estimate, so the real core reaches the target a
little later than the single-exponential prediction.
\end{mastery}

\begin{mastery}
The one-term result of the previous box is the tail of a full eigenfunction
expansion, and the machinery that produces it is the separation of variables
that recurs across the linear partial differential equations of engineering.
Take the symmetric plane wall again, half-thickness $L$, initially at uniform
$T_i$, suddenly exposed on both faces to a fluid at $T_\infty$ with coefficient
$h$. Writing $\theta = T - T_\infty$ and the dimensionless position
$\xi = x/L$, the problem is
\[
\frac{\partial \theta}{\partial t} = \alpha \frac{\partial^2 \theta}{\partial
x^2}, \quad \theta(x,0) = \theta_i, \quad \left.\frac{\partial \theta}{\partial
x}\right|_{0} = 0, \quad -k\left.\frac{\partial \theta}{\partial x}\right|_{L}
= h\,\theta(L,t).
\]
Seeking a product solution $\theta = X(x)\,\Gamma(t)$ and dividing by
$X\Gamma$ forces each side to a separation constant $-\lambda^2/L^2$, giving a
decaying time factor $\Gamma \propto \exp(-\lambda^2 \alpha t/L^2) =
\exp(-\lambda^2 \text{Fo})$ and a spatial factor $X = \cos(\lambda \xi)$ once
the symmetry condition at the centre kills the sine. The convective face
condition turns into the transcendental eigenvalue equation
$\lambda_n \tan\lambda_n = \text{Bi}$, an infinite ladder of roots
$\lambda_1 < \lambda_2 < \cdots$ that crowd toward $(n-\tfrac12)\pi$ for large
$\text{Bi}$ and toward $\sqrt{\text{Bi}}\,$ at the bottom for small $\text{Bi}$.
The general solution superposes the modes,
\[
\frac{\theta(\xi,\text{Fo})}{\theta_i} = \sum_{n=1}^{\infty} C_n
\cos(\lambda_n \xi)\,e^{-\lambda_n^2 \text{Fo}},
\]
and the coefficients $C_n$ are fixed by the initial condition through the
orthogonality of the eigenfunctions. The cosines $\cos(\lambda_n \xi)$ are
mutually orthogonal on $[0,1]$ under these boundary conditions, so projecting
the flat initial profile onto each mode gives
\[
C_n = \frac{\displaystyle\int_0^1 \cos(\lambda_n \xi)\,d\xi}
{\displaystyle\int_0^1 \cos^2(\lambda_n \xi)\,d\xi} =
\frac{4\sin\lambda_n}{2\lambda_n + \sin 2\lambda_n}.
\]
Each term decays at its own rate $\exp(-\lambda_n^2 \text{Fo})$, and because the
eigenvalues grow, the higher modes vanish first: after $\text{Fo} \gtrsim 0.2$
the series is the single leading term to better than one percent, which is the
one-term approximation the Heisler charts tabulate and the previous box
expanded. The whole apparatus, separate, solve an eigenvalue problem, project
the data by orthogonality, sum the modes, carries over without change to
transient diffusion in cylinders and spheres (where the eigenfunctions become
Bessel functions and spherical sincs), to the vibrating beam, and to the
diffusion of mass in chapter~7. The conduction transient is the cleanest place
to meet the method because the eigenvalues are real, the decay is monotone, and
the physics of why the short-wavelength modes die first, smaller features carry
steeper curvature and so larger $\partial^2 T/\partial x^2$, reads directly off
the heat equation.
\end{mastery}

\section{Convection: free and forced}\pathtag{core}

Convection is heat transfer between a solid surface and a moving
fluid in contact with it. The heat flux at the surface is given by
\engterm{Newton's law of cooling},
\[
q = h (T_s - T_\infty),
\]
where $h$ is the \engterm{convection heat-transfer coefficient}
(W/(m$^2\cdot$K)), $T_s$ is the surface temperature, and $T_\infty$
is the bulk fluid temperature. The coefficient $h$ depends on the
fluid properties, the flow geometry, the flow regime (laminar or
turbulent), and the velocity. Newton's law is a definition of $h$
rather than a physical law; the engineering work is to compute or
correlate $h$ for the operating conditions.

% Figure: thermal and velocity boundary layers growing over a flat plate.
% Free stream U_inf, T_inf approaches the leading edge; the boundary layer
% thickness grows like sqrt(x) in the laminar region, then thickens and
% transitions to turbulent. Plate held at T_s.
\begin{figure}[htbp]
\centering
\begin{tikzpicture}[scale=1.0, >=latex]
% plate
\fill[enggray!25] (0,-0.35) rectangle (10,0);
\draw[thick] (0,0) -- (10,0);
\node[font=\scriptsize] at (5,-0.62) {heated plate at $T_s$};
% free-stream arrows
\foreach \y in {1.4,1.9,2.4}{
  \draw[->, engblue, thick] (-1.1,\y) -- (-0.2,\y);
}
\node[font=\scriptsize, engblue, anchor=east] at (-1.15,1.9) {$U_\infty,\,T_\infty$};
% laminar boundary-layer edge: delta ~ sqrt(x), from 0 to x=5
\draw[engred, thick, samples=60, domain=0.02:5, variable=\x]
  plot ({\x}, {0.62*sqrt(\x)});
% transition marker
\draw[enggray, dashed] (5,0) -- (5,2.2);
\node[font=\scriptsize, enggray] at (5,2.45) {transition};
% turbulent region: thicker, grows faster ~ x^{0.8}, continue to x=10
\draw[engred, thick, samples=60, domain=5:10, variable=\x]
  plot ({\x}, {1.275*0.30*(\x^0.8)});
\node[font=\scriptsize, engred, anchor=west] at (9.0,2.7) {$\delta(x)$};
% velocity profile glyphs at two stations
\foreach \xs/\dl in {2/0.88,8/2.0}{
  \draw[engblue, thin] (\xs,0) -- (\xs,\dl);
  \foreach \f in {0.25,0.5,0.75,1.0}{
    \draw[->, engblue, thin] (\xs,{\f*\dl}) -- ({\xs + 0.7*sqrt(\f)},{\f*\dl});
  }
}
\node[font=\scriptsize, engblue] at (2.0,1.15) {laminar};
\node[font=\scriptsize, engblue] at (8.0,2.3) {turbulent};
% x axis
\draw[->] (0,-0.95) -- (10.2,-0.95) node[right, font=\scriptsize] {$x$};
\node[font=\scriptsize] at (0,-1.2) {leading edge};
\end{tikzpicture}
\caption[Boundary layer growth over a heated flat plate]{The boundary
layer growing over a heated flat plate. The free stream arrives at
velocity $U_\infty$ and temperature $T_\infty$; the layer of fluid
slowed and warmed by the plate thickens with distance, as
$\delta \sim \sqrt{x}$ in the laminar region near the leading edge. Past
a transition Reynolds number near $5\times 10^5$ the layer becomes
turbulent, thickens faster, and mixes more vigorously, which raises the
local convection coefficient. The velocity profiles sketched at two
stations show the no-slip condition at the wall and the approach to the
free stream at the boundary-layer edge. The local Nusselt-number
correlations of section 5.3 follow this same $\sqrt{x}$ scaling in the
laminar region and the steeper $x^{4/5}$ scaling once turbulent.}
\label{fig:vol04ch05:boundary-layer}
\end{figure}


The coefficient $h$ is set by the boundary layer, the thin region of
slowed and thermally affected fluid next to the surface
(Figure~\ref{fig:vol04ch05:boundary-layer}). A thin boundary layer
puts a steep temperature gradient at the wall and a large $h$; a thick
or stagnant one insulates the surface. The
\texttt{convection\_correlations.py} module evaluates the Nusselt-number
correlations below and converts them to $h$, and
\texttt{convection\_coefficients.csv} lists the working ranges by
regime.

Two boundary layers grow together at the wall, and the Prandtl number
sets their relative thickness. The velocity boundary layer is the region
of slowed fluid; the thermal boundary layer is the region of changed
temperature. Their thickness ratio scales as $\delta/\delta_t \sim
\text{Pr}^{1/3}$, so a fluid with $\text{Pr} \approx 1$ (air, near
$0.7$) carries the two layers at comparable thickness, a high-Prandtl
oil ($\text{Pr}$ in the hundreds) has a thin thermal layer buried deep
inside a thick velocity layer, and a liquid metal ($\text{Pr} \sim
0.01$) has a thermal layer far thicker than the velocity layer because
heat diffuses faster than momentum. This is why the same correlation
carries $\text{Pr}$ to a fractional power, and why liquid metals,
despite their thick thermal layers, deliver enormous convection
coefficients: their high conductivity dominates the thin-layer
disadvantage. The Nusselt number these correlations return is itself a
ratio worth reading physically, the convective heat transfer at the wall
divided by the heat that pure conduction across the same fluid layer
would carry; $\text{Nu} = 1$ means the moving fluid does no better than
a still layer, and the large Nusselt numbers of turbulent flow measure
how much the motion improves on conduction alone.

% Figure: the relative thickness of the velocity and thermal boundary layers
% for three Prandtl-number regimes. For Pr ~ 1 (gases) the two layers grow
% together; for Pr >> 1 (oils) the thermal layer is buried inside a thick
% velocity layer; for Pr << 1 (liquid metals) the thermal layer outruns the
% velocity layer. The ratio delta/delta_t ~ Pr^(1/3) is the source of the
% Prandtl power in every forced-convection correlation.
\begin{figure}[htbp]
\centering
\begin{tikzpicture}
\begin{axis}[
  width=13cm, height=5.2cm,
  xlabel={distance along plate $x$ (arb.)},
  ylabel={layer thickness (arb.)},
  xmin=0, xmax=10, ymin=0, ymax=3.2,
  axis lines=left,
  legend pos=north west,
  legend cell align=left,
  tick label style={font=\scriptsize},
  label style={font=\small},
  legend style={font=\scriptsize},
  domain=0:10, samples=60,
]
% velocity layer (common reference, ~ sqrt(x))
\addplot[engblue, very thick] {0.55*sqrt(x)};
\addlegendentry{velocity layer $\delta$}
% thermal layer, Pr ~ 1 (air): comparable
\addplot[englime, very thick, dashed] {0.60*sqrt(x)};
\addlegendentry{thermal $\delta_t$, $\text{Pr}\approx 1$ (air)}
% thermal layer, Pr >> 1 (oil): thinner
\addplot[engred, very thick, dotted] {0.24*sqrt(x)};
\addlegendentry{thermal $\delta_t$, $\text{Pr}\gg 1$ (oil)}
% thermal layer, Pr << 1 (liquid metal): thicker
\addplot[engorange, very thick, dash dot] {1.05*sqrt(x)};
\addlegendentry{thermal $\delta_t$, $\text{Pr}\ll 1$ (liquid metal)}
\end{axis}
\end{tikzpicture}
\caption[Velocity and thermal boundary-layer thickness by Prandtl number]{The
velocity boundary layer (the region of slowed fluid) and the thermal boundary
layer (the region of changed temperature) grow together along a flat plate,
each thickening as the square root of distance. The Prandtl number fixes their
ratio, $\delta/\delta_t \sim \text{Pr}^{1/3}$. A gas with $\text{Pr}\approx 1$
carries the two layers at comparable thickness; a high-Prandtl oil buries a
thin thermal layer deep inside a thick velocity layer, so the wall temperature
gradient is steep and the convection coefficient large for the same velocity
field; a liquid metal with $\text{Pr}\ll 1$ has a thermal layer far thicker than
the velocity layer because heat diffuses faster than momentum. The fractional
power of $\text{Pr}$ in every flat-plate Nusselt correlation is this thickness
ratio in disguise.}
\label{fig:vol04ch05:prandtl-layers}
\end{figure}


\engterm{Where the correlation form comes from}. The result
$\text{Nu} = f(\text{Re}, \text{Pr})$ is forced by the structure of the
boundary-layer equations and confirmed by dimensional analysis; the numerical
coefficients are all the experiment or the exact solution must then supply. Two
arguments make this concrete. The dimensional argument lists the quantities that
fix the convection coefficient $h$ for forced flow over a surface: the velocity
$u$, the length $L$, the fluid density $\rho$, viscosity $\mu$, conductivity
$k$, and specific heat $c_p$. Seven quantities (counting $h$) in four base
dimensions of mass, length, time, and temperature reduce by the Buckingham
theorem to three dimensionless groups, and the natural choice of those groups is
the Nusselt number $\text{Nu} = hL/k$, the Reynolds number
$\text{Re} = \rho u L/\mu$, and the Prandtl number $\text{Pr} = c_p\mu/k$. No
relation among $h$ and its six causes can be anything other than
$\text{Nu} = f(\text{Re},\text{Pr})$; the experiment fixes the function, not its
arguments. The physical argument recovers the same form and explains the
exponents. The convection coefficient is the wall conduction into the fluid
divided by the driving difference, $h = -k(\partial T/\partial y)_{\text{wall}}/
(T_s - T_\infty)$, and the wall gradient is set by the temperature drop
$T_s - T_\infty$ spread across the thermal boundary-layer thickness $\delta_t$,
so $h \sim k/\delta_t$ and $\text{Nu} = hL/k \sim L/\delta_t$. The Nusselt
number is the plate length measured in thermal-boundary-layer thicknesses. The
velocity layer grows along a laminar flat plate as $\delta/L \sim
\text{Re}_L^{-1/2}$ from the balance of inertia against viscous diffusion of
momentum, and the thermal layer relates to it by the Prandtl ratio
$\delta/\delta_t \sim \text{Pr}^{1/3}$ of
Figure~\ref{fig:vol04ch05:prandtl-layers}, so
\[
\text{Nu}_L \sim \frac{L}{\delta_t} = \frac{L}{\delta}\cdot
\frac{\delta}{\delta_t} \sim \text{Re}_L^{1/2}\,\text{Pr}^{1/3},
\]
which is exactly the form of the laminar flat-plate correlation
$\text{Nu}_L = 0.664\,\text{Re}_L^{1/2}\,\text{Pr}^{1/3}$ below, with the
$0.664$ the only thing the scaling cannot give and the exact Blasius
boundary-layer solution does. The turbulent and pipe correlations carry larger
Reynolds exponents (near $0.8$) because turbulent mixing thins the layer faster
than laminar diffusion, but the template, a Reynolds power from the momentum
layer times a Prandtl power from the thermal-layer ratio, is the same across the
table. An engineer who remembers the form can reconstruct any of the
correlations to within its leading coefficient and can sanity-check a tabulated
one whose exponents look wrong.

\engterm{Forced convection} is convection driven by an external
flow (a fan, a pump, the wind, the motion of a vehicle). The relevant
dimensionless groups are the Reynolds number $\text{Re} = \rho u
L/\mu$ (the ratio of inertial to viscous forces), the Prandtl number
$\text{Pr} = c_p \mu/k$ (the ratio of momentum diffusivity to thermal
diffusivity), and the Nusselt number $\text{Nu} = hL/k$ (the
dimensionless convection coefficient). Correlations of the form $\text{Nu}
= f(\text{Re}, \text{Pr})$ cover the working geometries:

\begin{itemize}[leftmargin=*]
  \item Flow over a flat plate, laminar: $\text{Nu}_x = 0.332\,
    \text{Re}_x^{1/2}\,\text{Pr}^{1/3}$ (local), $\text{Nu}_L = 0.664\,
    \text{Re}_L^{1/2}\,\text{Pr}^{1/3}$ (averaged).
  \item Flow over a flat plate, turbulent: $\text{Nu}_x = 0.0296\,
    \text{Re}_x^{4/5}\,\text{Pr}^{1/3}$.
  \item Flow inside a pipe, turbulent (Dittus-Boelter): $\text{Nu}_D
    = 0.023\,\text{Re}_D^{0.8}\,\text{Pr}^{n}$, with $n = 0.4$ for
    heating and $n = 0.3$ for cooling.
  \item Flow inside a pipe, laminar (fully developed, constant wall
    temperature): $\text{Nu}_D = 3.66$.
\end{itemize}

Working values of $h$: still air over a building exterior, $5$ to
$25\,\text{W}/(\text{m}^2\cdot\si{\kelvin})$; air with a $5\,\text{m/s}$
wind, $30$ to $50\,\text{W}/(\text{m}^2\cdot\si{\kelvin})$; water in a
pipe at $1\,\text{m/s}$, $1000$ to $3000\,\text{W}/(\text{m}^2\cdot
\si{\kelvin})$; condensing steam on a vertical tube, $5000$ to $20000\,
\text{W}/(\text{m}^2\cdot\si{\kelvin})$; nucleate-boiling water, $5000$
to $50000\,\text{W}/(\text{m}^2\cdot\si{\kelvin})$. Six orders of
magnitude separate the lightest (still gas) from the heaviest (boiling
or condensing water) convection regimes; the engineer who chooses the
heat-transfer mode chooses the working surface area required for a
given duty.

% Figure: averaged Nusselt number versus Reynolds number for external flow
% over a flat plate, showing the laminar Re^{1/2} branch and the turbulent
% Re^{4/5} branch meeting near the transition Reynolds number. Computed
% from the two correlations at Pr = 0.71, log-log axes.
\begin{figure}[htbp]
\centering
\begin{tikzpicture}
\begin{axis}[
  width=12cm, height=7.5cm,
  xmode=log, ymode=log,
  xlabel={Reynolds number $\text{Re}_L$},
  ylabel={averaged Nusselt number $\text{Nu}_L$},
  xmin=1e3, xmax=1e7,
  ymin=10, ymax=2e4,
  axis lines=left,
  tick label style={font=\scriptsize},
  label style={font=\small},
  legend style={font=\scriptsize, at={(0.02,0.98)}, anchor=north west},
  legend cell align=left,
  legend entries={laminar $0.664\,\text{Re}^{1/2}\text{Pr}^{1/3}$,turbulent $0.037\,\text{Re}^{4/5}\text{Pr}^{1/3}$},
  clip=false,
]
% Pr = 0.71 -> Pr^{1/3} = 0.892
% laminar branch
\addplot[engblue, very thick, domain=1e3:5e5, samples=40] {0.664*sqrt(x)*0.892};
% turbulent branch
\addplot[engred, very thick, domain=5e5:1e7, samples=40] {0.037*(x^0.8)*0.892};
% transition marker
\draw[enggray, dotted] (axis cs:5e5,10) -- (axis cs:5e5,2e4);
\node[font=\scriptsize, enggray, anchor=south, rotate=90] at (axis cs:5e5,20) {transition $\text{Re}\approx 5\times10^5$};
\end{axis}
\end{tikzpicture}
\caption[Nusselt number across the laminar and turbulent branches]{The
averaged Nusselt number for flow over a flat plate, plotted against the
plate Reynolds number at a gas Prandtl number of $0.71$. The laminar
branch rises as $\text{Re}^{1/2}$, the slope the boundary-layer scaling
of the chapter predicts; past the transition Reynolds number near
$5\times10^5$ the turbulent branch takes over and climbs faster, as
$\text{Re}^{4/5}$, because turbulent mixing thins the boundary layer more
quickly than laminar diffusion. The break in slope is the visible
signature of the change of regime, and the design consequence is direct:
running a heat-transfer surface into its turbulent branch buys a
disproportionate gain in convection coefficient at the cost of a steeper
rise in pressure drop.}
\label{fig:vol04ch05:nusselt-regimes}
\end{figure}


Figure~\ref{fig:vol04ch05:nusselt-regimes} draws the two flat-plate
correlations across their Reynolds range and shows the change of slope at
the transition. The laminar branch climbs as $\text{Re}^{1/2}$ and the
turbulent branch as roughly $\text{Re}^{4/5}$, so the convection
coefficient jumps upward as a surface crosses into turbulence, which is the
reason a designer often deliberately trips the boundary layer, with a wire,
a step, or a rough leading strip, to force the turbulent branch and buy the
higher coefficient. The gain is not free: the same turbulence that thins
the thermal layer thickens the wake and raises the pressure drop, so the
pumping power to push the fluid climbs faster than the heat transfer, and
the operating point is chosen where the marginal heat gained justifies the
marginal power spent.

\engterm{A forced-convection duct heater}. The pipe correlation earns its
keep on any surface that heats or cools a flowing stream, and we work one
instance in full. Air at $20\,\degC$ flows at $6\,\text{m/s}$ through a
circular duct of diameter $D = 0.15\,\text{m}$ whose wall is held at
$60\,\degC$ by a surrounding jacket; we find the convection coefficient and
the length of duct needed to warm the air to $45\,\degC$. For air near the
mean temperature take $\rho = 1.15\,\text{kg/m}^3$, $\nu = 1.6\times10^{-5}
\,\text{m}^2/\text{s}$, $k = 0.027\,\text{W}/(\text{m}\cdot\si{\kelvin})$,
$c_p = 1007\,\text{J}/(\text{kg}\cdot\si{\kelvin})$, $\text{Pr} = 0.71$. The
Reynolds number is
\[
\text{Re}_D = \frac{u D}{\nu} = \frac{6\times0.15}{1.6\times10^{-5}} =
5.6\times10^4,
\]
turbulent, so Dittus-Boelter for heating ($n = 0.4$) gives
\[
\text{Nu}_D = 0.023\,\text{Re}_D^{0.8}\,\text{Pr}^{0.4} = 0.023\times
(5.6\times10^4)^{0.8}\times0.71^{0.4} = 0.023\times6.19\times10^3\times
0.867 = 123,
\]
and the coefficient is $h = \text{Nu}_D\,k/D = 123\times0.027/0.15 =
22.2\,\text{W}/(\text{m}^2\cdot\si{\kelvin})$. The air mass flow is
$\dot{m} = \rho u (\pi D^2/4) = 1.15\times6\times(\pi\times0.15^2/4) =
0.122\,\text{kg/s}$. A constant-wall-temperature duct warms its stream
exponentially toward the wall, the outlet approach following
\[
\frac{T_w - T_{\text{out}}}{T_w - T_{\text{in}}} = \exp\!\left(-\frac{h\,
\pi D\,L}{\dot{m}c_p}\right),
\]
the same first-order form the lumped body obeyed, with duct length playing
the role of time. The target sets the left side to $(60-45)/(60-20) =
15/40 = 0.375$, so
\[
\frac{h\,\pi D\,L}{\dot{m}c_p} = -\ln(0.375) = 0.981, \quad L =
\frac{0.981\times\dot{m}c_p}{h\,\pi D} = \frac{0.981\times0.122\times1007}
{22.2\times\pi\times0.15} = 11.5\,\text{m}.
\]
The duct must run about twelve metres to close three-quarters of the
temperature gap, and the exponential approach warns why chasing the last
few degrees is expensive: warming the air to $55\,\degC$ instead of $45$
would drive the approach ratio to $0.125$ and need $\ln(8)/\ln(2.67)$
times the length, more than twice as long, for a further ten degrees. The
design lever is the coefficient, and pushing the velocity to $12\,\text{m/s}$
would raise $\text{Re}$ and $h$ by $2^{0.8} = 1.74$ and shorten the duct in
proportion, at the cost of a pressure drop that rises with the square of
velocity.

\engterm{Free convection} (also called natural convection) is
convection driven by buoyancy: a fluid heated by a hot surface
expands, becomes less dense, and rises, drawing cooler fluid into
contact with the surface. The relevant dimensionless groups are the
Grashof number $\text{Gr} = g\beta(T_s - T_\infty)L^3/\nu^2$ (with
$\beta$ the volumetric thermal expansion coefficient and $\nu$ the
kinematic viscosity), the Rayleigh number $\text{Ra} = \text{Gr}\,
\text{Pr}$, and the Nusselt number as before. The Rayleigh number
sets the regime: $\text{Ra} < 10^9$ for laminar free convection,
$\text{Ra} > 10^{10}$ for turbulent. Correlations:

\begin{itemize}[leftmargin=*]
  \item Vertical isothermal plate: $\text{Nu}_L = \{0.825 + 0.387\,
    \text{Ra}_L^{1/6}/[1 + (0.492/\text{Pr})^{9/16}]^{8/27}\}^2$
    (Churchill and Chu).
  \item Horizontal hot plate facing up: $\text{Nu}_L = 0.54\,
    \text{Ra}_L^{1/4}$ for $\text{Ra}_L < 10^7$; $\text{Nu}_L = 0.15\,
    \text{Ra}_L^{1/3}$ for higher.
  \item Horizontal cylinder: $\text{Nu}_D = \{0.6 + 0.387\,
    \text{Ra}_D^{1/6}/[1 + (0.559/\text{Pr})^{9/16}]^{8/27}\}^2$.
\end{itemize}

Free-convection coefficients are typically $2$ to $25\,\text{W}/(
\text{m}^2\cdot\si{\kelvin})$ in air and $50$ to $1000\,\text{W}/(
\text{m}^2\cdot\si{\kelvin})$ in water. They are smaller than forced-
convection coefficients by factors of $2$ to $10$ at the same
temperature difference; free convection is the dominant heat-transfer
mode in passive systems (building exterior walls, electronics
cabinets without fans, hot drinks cooling on a counter) and a poor
choice when high heat-transfer rates are required.

\engterm{Boiling and condensation} are special convection regimes
involving phase change. The boiling curve (heat flux versus surface
superheat) has several regimes: free-convection boiling (low
superheat, no bubbles), nucleate boiling (bubbles form at active
sites on the heated surface), partial film boiling (the bubbles
merge into an unstable film), and stable film boiling (a continuous
vapour film blanks the surface, with much lower heat transfer). The
\engterm{critical heat flux} is the peak of the nucleate-boiling
regime, above which film boiling sets in abruptly. Operating above
the critical heat flux causes ``burnout'' in heated tubes (the wall
temperature jumps by hundreds of degrees as the cooling collapses);
the design discipline for water-cooled reactors and boilers is to
maintain margins below the critical heat flux at every point.

% Figure: the pool-boiling curve for saturated water at one atmosphere,
% heat flux versus wall superheat on log axes, with the four regimes
% (free convection, nucleate, transition, film) and the critical-heat-
% flux peak and the Leidenfrost minimum marked. Schematic; the curve is
% drawn from a small set of representative points, not a correlation fit.
\begin{figure}[htbp]
\centering
\begin{tikzpicture}
\begin{axis}[
  width=12cm, height=8cm,
  xmode=log, ymode=log,
  xlabel={wall superheat $\Delta T_e = T_s - T_{\text{sat}}$ ($\si{\kelvin}$)},
  ylabel={surface heat flux $q$ (W/m$^2$)},
  xmin=1, xmax=1000, ymin=1e3, ymax=3e6,
  axis lines=left,
  tick label style={font=\scriptsize},
  label style={font=\small},
  clip=false,
]
% free-convection + nucleate rising branch up to critical heat flux
\addplot[engblue, very thick] coordinates {
  (2,3e3) (5,2e4) (10,1e5) (20,5e5) (30,1e6) };
% critical heat flux peak marker
\addplot[only marks, mark=*, engred, mark options={fill=engred}]
  coordinates {(30,1.0e6)};
% transition boiling (falling branch), drawn dashed for instability
\addplot[engamber, very thick, dashed] coordinates {
  (30,1e6) (60,3e5) (120,3.5e4) };
% Leidenfrost minimum marker
\addplot[only marks, mark=square*, engamber, mark options={fill=engamber}]
  coordinates {(120,3.5e4)};
% stable film boiling rising branch
\addplot[enggray, very thick] coordinates {
  (120,3.5e4) (300,1.2e5) (700,6e5) };
% regime labels
\node[font=\scriptsize, engblue, anchor=west] at (axis cs:2.2,7e3) {free convection};
\node[font=\scriptsize, engblue, anchor=south east] at (axis cs:18,4e5) {nucleate};
\node[font=\scriptsize, engred, anchor=south] at (axis cs:30,1.25e6) {critical heat flux};
\node[font=\scriptsize, engamber, anchor=south west] at (axis cs:55,3.5e5) {transition (unstable)};
\node[font=\scriptsize, engamber, anchor=north] at (axis cs:120,2.6e4) {Leidenfrost minimum};
\node[font=\scriptsize, enggray, anchor=west] at (axis cs:300,1.4e5) {film boiling};
% the burnout jump at constant flux
\draw[engred, ->, thick] (axis cs:30,1.0e6) -- (axis cs:430,1.0e6);
\node[font=\scriptsize, engred, anchor=south] at (axis cs:110,1.05e6) {burnout (flux-controlled)};
\end{axis}
\end{tikzpicture}
\caption[Pool-boiling curve for saturated water at one
atmosphere]{The pool-boiling curve for saturated water at one
atmosphere, heat flux against wall superheat on logarithmic axes. The
heat flux climbs steeply through the nucleate-boiling branch, where
bubbles grow and detach from active sites, to the critical heat flux at
a superheat near $30\,\si{\kelvin}$. A surface driven harder than the
critical heat flux on a flux-controlled heater cannot follow the falling
transition branch; it jumps along the horizontal arrow to the film-
boiling branch, where the same flux now demands a superheat of hundreds
of kelvin, so the wall temperature leaps by hundreds of degrees and the
surface burns out. The dashed transition branch is reachable only on a
temperature-controlled surface. The design rule for flux-controlled
boilers and water-cooled reactors is to hold the operating flux a stated
margin below the critical heat flux at every point on the surface.}
\label{fig:vol04ch05:boiling-curve}
\end{figure}


The boiling curve of Figure~\ref{fig:vol04ch05:boiling-curve} carries a
distinction that decides equipment survival: whether the surface is
\emph{flux-controlled} or \emph{temperature-controlled}. An electric
resistance heater, a nuclear fuel rod, and a soldering iron set the heat
flux and let the surface find its own temperature. Driven past the
critical heat flux, such a surface has no stable operating point on the
falling transition branch; it leaps along a line of constant flux to the
film-boiling branch, where the same flux demands a superheat of hundreds
of kelvin, and the wall temperature jumps to that value in a moment. A
copper surface survives the jump; a thin steel tube or a fuel cladding
melts, which is the burnout failure. A condenser surface or a steam-
heated jacket, by contrast, sets the surface temperature and reads off
whatever flux the curve allows, so it can sit on the transition branch
without the catastrophic jump. The water-cooled reactor and the fired
boiler are the flux-controlled cases, and the departure-from-nucleate-
boiling ratio (the operating flux divided by the local critical heat
flux) is the margin those designs are required to hold above unity at
every point on the heated surface and through every transient the plant
can reach. The same curve governs quenching in the lumped-capacitance
example above: a part plunged into water first sits in film boiling
(slow cooling under a vapour blanket), passes the Leidenfrost minimum as
it cools, and then cools fastest through the nucleate-boiling regime, so
the heat-transfer coefficient a quench delivers is not a constant but a
strong function of the part's own surface temperature.

Condensation is the mirror image of boiling and carries its own large
coefficients, which is why steam is the workhorse heating fluid of process
plant. When vapour meets a surface below its saturation temperature it
gives up its latent heat and the condensate drains away, and the mode of
drainage sets the coefficient. In \engterm{film condensation} the liquid
wets the surface and runs off as a continuous sheet, so the fresh vapour
must conduct its latent heat across a growing liquid film that is itself
the controlling resistance; coefficients run from $5000$ to $20000\,
\text{W}/(\text{m}^2\cdot\si{\kelvin})$ on a vertical tube, an order above
the boiling-water figures and three orders above still air. In
\engterm{dropwise condensation} the surface does not wet, the condensate
beads and rolls off before a film can build, and fresh vapour reaches bare
metal over most of the surface, lifting the coefficient by a further factor
of five to ten. Dropwise condensation is the highest ordinary
convective mode in the chapter's table, but it is hard to sustain because
the non-wetting promoter coatings degrade, so condenser design assumes the
film mode and treats any dropwise gain as margin. The film mode is also
the one that fails silently: a trace of non-condensable gas, air leaking
into a steam condenser, blankets the surface and collapses the coefficient,
because the incoming vapour must now diffuse through a stagnant gas layer to
reach the film, which is why condensers carry vents and why the mass-
transfer resistance of chapter~7 reappears here as a heat-transfer penalty.

\begin{mastery}
Nusselt's 1916 analysis of laminar film condensation on a vertical plate
derives the coefficient from a momentum-and-energy balance on the draining
film, and it is the cleanest place to see a convection coefficient built
from first principles rather than correlated. Take a vertical plate at
uniform temperature $T_s$ below the saturation temperature $T_{\text{sat}}$
of a quiescent vapour. Condensate collects as a film of local thickness
$\delta(x)$ that grows down the plate under gravity. In the film the flow
is slow enough to be a balance of gravity against viscous shear, so the
momentum equation reduces to $\mu\,d^2u/dy^2 = -\rho_l g$, which with a
no-slip wall and a shear-free vapour interface integrates to a parabolic
velocity profile whose mass flow per unit width is $\Gamma = \rho_l g
\delta^3/(3\mu_l)$, taking the vapour density as negligible against the
liquid. Heat crosses the film by pure conduction, since the film is thin
and nearly stagnant, so the local flux is $q'' = k_l(T_{\text{sat}} -
T_s)/\delta$, and each increment of that flux condenses vapour whose latent
heat $h_{fg}$ must be carried off by the thickening film. Equating the
conducted heat to the latent heat of the added condensate, $k_l\,\Delta T/
\delta\,dx = h_{fg}\,d\Gamma$, and substituting $\Gamma(\delta)$ gives a
first-order equation for $\delta(x)$ whose solution is
\[
\delta(x) = \left[\frac{4 k_l \mu_l (T_{\text{sat}} - T_s)\,x}{g \rho_l^2
h_{fg}}\right]^{1/4},
\]
the film growing as the fourth root of distance down the plate. The local
coefficient is $h_x = k_l/\delta$, so it \emph{falls} as $x^{-1/4}$, and
averaging over a plate of height $L$ gives the classic result
\[
\bar{h} = 0.943\left[\frac{g \rho_l^2 h_{fg} k_l^3}{\mu_l (T_{\text{sat}} -
T_s) L}\right]^{1/4}.
\]
The coefficient rises with the cube root of the liquid conductivity and
falls with the driving subcooling to the quarter power, so a condenser runs
best at a small temperature difference across a short drainage length, which
is why condenser tubes are mounted horizontally, cutting the drainage length
from the plate height to the tube diameter and holding the film thin. The
same quarter-power film-thickness law explains the non-condensable-gas
penalty of the paragraph above: a gas layer adds a diffusion resistance in
series with the conduction across the film, and the analysis that produced
$\bar{h}$ from a pure-conduction film no longer holds.
\end{mastery}

\section{Radiation: Stefan-Boltzmann, view factors}\pathtag{core}

Thermal radiation is the emission of electromagnetic energy by all
bodies at non-zero absolute temperature. The radiation spectrum,
intensity, and angular distribution are governed by the body's
temperature, surface emissivity, and geometry.

\engterm{The Stefan-Boltzmann law}. A blackbody (an idealised perfect
absorber and emitter) at absolute temperature $T$ emits total
radiative flux
\[
E_b = \sigma T^4,
\]
where $\sigma = 5.67 \times 10^{-8}\,\text{W}/(\text{m}^2\cdot\si{
\kelvin}^4)$ is the Stefan-Boltzmann constant. The flux scales as the
fourth power of absolute temperature, the strongest scaling in
engineering heat transfer. Working values: $E_b$ at $300\,\si{\kelvin}$
(room temperature) is $459\,\text{W}/\text{m}^2$; at $1000\,\si{\kelvin}$
(red-hot metal), $57000\,\text{W}/\text{m}^2$; at $5800\,\si{\kelvin}$
(sun's surface), $6.4 \times 10^7\,\text{W}/\text{m}^2$. The
\texttt{radiation\_exchange.py} module evaluates these and the
two-surface exchange formulas below.

% Figure: blackbody spectral emissive power (Planck) vs wavelength for a
% few temperatures, illustrating the Wien shift of the peak to shorter
% wavelength as T rises and the strong growth of total radiated power.
% Curves are normalised per curve to peak near unity for display on one
% set of axes; the normalisation constants are the Planck peak values.
\begin{figure}[htbp]
\centering
\begin{tikzpicture}
\begin{axis}[
  width=12cm, height=7.5cm,
  xlabel={wavelength $\lambda$ ($\mu$m)},
  ylabel={spectral emissive power (normalised)},
  xmin=0, xmax=12, ymin=0, ymax=1.15,
  axis lines=left,
  legend pos=north east,
  legend cell align=left,
  tick label style={font=\scriptsize},
  label style={font=\small},
  legend style={font=\scriptsize},
  clip=false,
]
% Planck shape E_lambda ~ 1/(lambda^5 (exp(c2/(lambda T)) - 1)), c2=14388
% um*K, each curve normalised to its own peak. Coordinates precomputed to
% keep pgfmath out of its fixed-point overflow range (lambda^5 overflows).
\addplot[engred, very thick] coordinates {
(0.40,0.0000) (0.80,0.0000) (1.00,0.0014) (1.20,0.0110) (1.40,0.0435)
(1.60,0.1111) (1.80,0.2150) (2.00,0.3449) (2.20,0.4851) (2.40,0.6207)
(2.60,0.7406) (2.80,0.8386) (3.00,0.9122) (3.20,0.9619) (3.40,0.9901)
(3.60,0.9999) (3.80,0.9946) (4.00,0.9774) (4.40,0.9185) (4.80,0.8416)
(5.20,0.7585) (5.60,0.6766) (6.00,0.5996) (6.60,0.4971) (7.20,0.4111)
(8.00,0.3198) (9.00,0.2357) (10.00,0.1761) (11.00,0.1335) (12.00,0.1026)
};
\addlegendentry{$T=800$ K (peak $\approx 3.6\,\mu$m)}
\addplot[engamber, very thick] coordinates {
(0.40,0.0000) (0.60,0.0003) (0.80,0.0110) (1.00,0.0725) (1.20,0.2150)
(1.40,0.4147) (1.60,0.6207) (1.80,0.7926) (2.00,0.9122) (2.20,0.9785)
(2.40,0.9999) (2.60,0.9873) (2.80,0.9512) (3.00,0.9004) (3.40,0.7794)
(3.80,0.6568) (4.20,0.5463) (4.60,0.4521) (5.00,0.3740) (5.60,0.2826)
(6.20,0.2156) (7.00,0.1530) (8.00,0.1026) (9.00,0.0710) (10.00,0.0504)
(11.00,0.0368) (12.00,0.0274)
};
\addlegendentry{$T=1200$ K (peak $\approx 2.4\,\mu$m)}
\addplot[englime!70!black, very thick] coordinates {
(0.40,0.0003) (0.60,0.0324) (0.80,0.2150) (1.00,0.5199) (1.20,0.7926)
(1.40,0.9516) (1.60,0.9999) (1.80,0.9716) (2.00,0.9004) (2.20,0.8107)
(2.40,0.7171) (2.60,0.6277) (2.80,0.5463) (3.00,0.4741) (3.40,0.3567)
(3.80,0.2700) (4.20,0.2064) (4.60,0.1595) (5.00,0.1248) (5.60,0.0882)
(6.20,0.0639) (7.00,0.0429) (8.00,0.0274) (9.00,0.0182) (10.00,0.0126)
(11.00,0.0089) (12.00,0.0065)
};
\addlegendentry{$T=1800$ K (peak $\approx 1.6\,\mu$m)}
\node[font=\scriptsize, engblue, anchor=west] at (axis cs:4.6,0.92)
  {Wien: $\lambda_{\max}T=2898\,\mu\mathrm{m\cdot K}$};
\draw[engblue, dashed, ->] (axis cs:4.55,0.92) -- (axis cs:1.7,0.95);
\end{axis}
\end{tikzpicture}
\caption[Blackbody spectral emissive power]{Blackbody spectral emissive
power against wavelength for three temperatures, each curve normalised to
its own peak so all three fit on one set of axes. The peak wavelength
moves to the left as temperature rises, following the Wien displacement
law $\lambda_{\max}T = 2898\,\mu\mathrm{m\cdot K}$: a body at $800$ K
radiates mostly in the infrared, while at $1800$ K an appreciable part of
the emission has crossed into the near infrared and the body begins to
glow visibly. The total area under each unnormalised curve grows as
$T^4$, the Stefan-Boltzmann result, so the hotter body radiates far more
power than the normalised plot suggests.}
\label{fig:vol04ch05:blackbody-spectrum}
\end{figure}


The total flux is the area under the spectral curve of
Figure~\ref{fig:vol04ch05:blackbody-spectrum}. As temperature rises the
peak of that curve shifts to shorter wavelength (the Wien displacement
law $\lambda_{\max} T = 2898\,\mu\text{m}\cdot\si{\kelvin}$) while the
area grows as $T^4$, which is why a heated body first glows dull red,
then orange, then white as it climbs through visible wavelengths.

\begin{mastery}
The Stefan-Boltzmann law is the wavelength integral of Planck's
spectral distribution. The spectral emissive power of a blackbody is
\[
E_{b,\lambda}(\lambda, T) = \frac{2\pi h c^2}{\lambda^5}
\frac{1}{e^{hc/(\lambda k_B T)} - 1},
\]
with $h$ Planck's constant, $c$ the speed of light, and $k_B$
Boltzmann's constant. Integrating over all wavelengths and substituting
$u = hc/(\lambda k_B T)$ collapses the temperature dependence to a
single factor $T^4$ times the dimensionless integral
$\int_0^\infty u^3/(e^u - 1)\,du = \pi^4/15$, which yields
$\sigma = 2\pi^5 k_B^4/(15 h^3 c^2)$. The fourth-power law is therefore
not an empirical fit; it follows from the $\lambda^{-5}$ prefactor and
the Bose-Einstein form of the photon occupation. Differentiating
$E_{b,\lambda}$ with respect to $\lambda$ and setting the derivative to
zero gives the Wien displacement law, with the constant
$2898\,\mu\text{m}\cdot\si{\kelvin}$ emerging as $hc/(4.965\,k_B)$.
\end{mastery}

A real (non-black) surface emits a fraction of the blackbody flux,
$E = \epsilon \sigma T^4$, where $\epsilon$ is the
\engterm{emissivity} of the surface. Emissivity ranges from near
zero (polished aluminium, $\epsilon \approx 0.04$) to near unity
(rough black surfaces, oxidised metals, paint, wood, water, $\epsilon
\approx 0.9$). The emissivity of a surface for thermal radiation is
generally close to its absorptivity at the same wavelength (Kirchhoff's
law), so a good emitter is also a good absorber and a polished
metal is both a poor emitter and a poor absorber.

\engterm{Net radiative exchange between two surfaces}. For two
surfaces at temperatures $T_1$ and $T_2$, the net radiative heat
flow from 1 to 2 is
\[
\dot{Q}_{1 \to 2} = \frac{\sigma (T_1^4 - T_2^4)}{(1 - \epsilon_1)/(
\epsilon_1 A_1) + 1/(A_1 F_{12}) + (1 - \epsilon_2)/(\epsilon_2 A_2)},
\]
where $F_{12}$ is the \engterm{view factor} from surface 1 to surface
2 (the fraction of radiation leaving 1 that arrives at 2). View
factors satisfy reciprocity ($A_1 F_{12} = A_2 F_{21}$) and
summation ($\sum_j F_{ij} = 1$, with the sum including all surfaces
visible from $i$, including itself if $i$ is concave).

% Figure: view-factor geometry. Two finite surfaces A1 and A2 with their
% normals, a connecting ray of length r, and the angles theta1, theta2 the
% ray makes with each normal. The differential view-factor integrand
% cos(theta1)cos(theta2)/(pi r^2) is what gets integrated over both areas.
\begin{figure}[htbp]
\centering
\begin{tikzpicture}[scale=1.0, >=latex]
% surface 1 (lower left), tilted
\draw[thick, engblue, fill=engblue!10] (0,0) -- (2.0,0.5) -- (2.3,1.3) -- (0.3,0.8) -- cycle;
\node[font=\small, engblue] at (0.6,0.2) {$A_1$};
% point on surface 1
\coordinate (p1) at (1.1,0.65);
\fill (p1) circle (1.3pt);
% surface 2 (upper right), tilted differently
\draw[thick, engred, fill=engred!10] (6.0,3.2) -- (8.0,3.0) -- (8.2,4.0) -- (6.2,4.3) -- cycle;
\node[font=\small, engred] at (7.6,3.5) {$A_2$};
\coordinate (p2) at (7.0,3.65);
\fill (p2) circle (1.3pt);
% connecting ray
\draw[thick, enggray] (p1) -- (p2);
\node[font=\scriptsize, enggray] at ($(p1)!0.5!(p2)$) [above left] {$r$};
% normal at p1 (roughly upward-left perpendicular to surface 1)
\draw[->, engblue, thick] (p1) -- ($(p1) + (-0.5,1.4)$);
\node[font=\scriptsize, engblue] at ($(p1) + (-0.55,1.5)$) {$\hat{n}_1$};
% normal at p2 (roughly downward-left)
\draw[->, engred, thick] (p2) -- ($(p2) + (-0.6,-1.3)$);
\node[font=\scriptsize, engred] at ($(p2) + (-0.7,-1.45)$) {$\hat{n}_2$};
% angle labels
\node[font=\scriptsize] at ($(p1) + (0.15,0.75)$) {$\theta_1$};
\node[font=\scriptsize] at ($(p2) + (0.05,-0.7)$) {$\theta_2$};
% integrand annotation
\node[font=\scriptsize, anchor=west] at (3.4,1.6)
  {$dF_{1\to 2}=\dfrac{\cos\theta_1\,\cos\theta_2}{\pi r^2}\,dA_2$};
\end{tikzpicture}
\caption[View-factor geometry between two surfaces]{The geometry that
defines the view factor between two surfaces. A ray of length $r$ joins a
point on $A_1$ to a point on $A_2$, making angle $\theta_1$ with the
normal $\hat{n}_1$ and angle $\theta_2$ with the normal $\hat{n}_2$. The
fraction of diffuse radiation leaving the first point that reaches the
second falls with the cosine of each angle (a surface viewed obliquely
presents less projected area) and with the inverse square of the
separation. Integrating $\cos\theta_1\cos\theta_2/(\pi r^2)$ over both
areas gives the view factor $F_{12}$, which obeys the reciprocity
relation $A_1F_{12}=A_2F_{21}$ and the summation rule $\sum_j F_{ij}=1$
used in the radiation-network calculations of section 5.4.}
\label{fig:vol04ch05:view-factor-geometry}
\end{figure}


Figure~\ref{fig:vol04ch05:view-factor-geometry} shows the geometry that
defines the view factor: the cosine of the angle each surface makes with
the connecting ray, weighted by the inverse square of the separation,
integrated over both areas. Representative emissivities for the surface
resistance terms are tabulated in \texttt{surface\_emissivity.csv}.

For the simple case of two large parallel plates with $A_1 = A_2 = A$,
$F_{12} = 1$,
\[
\dot{Q} = \frac{A \sigma (T_1^4 - T_2^4)}{1/\epsilon_1 + 1/\epsilon_2
- 1}.
\]
For a small body (1) in a large enclosure (2), $F_{12} = 1$ and the
enclosure radiation reaches the body with no significant return,
giving
\[
\dot{Q}_{1 \to 2} = \epsilon_1 A_1 \sigma (T_1^4 - T_2^4).
\]
The two formulas are the working tools for most engineering radiation
problems.

\engterm{View-factor algebra for the common geometries}. The two
formulas above are the endpoints of a single algebra, and the
intermediate cases are reached by the same two rules the view factor
obeys. A convex or flat surface cannot see itself, so $F_{ii} = 0$; a
concave surface sees part of its own radiation return, so $F_{ii} > 0$.
The summation rule $\sum_j F_{ij} = 1$ closes the radiation leaving any
surface of an enclosure, and reciprocity $A_i F_{ij} = A_j F_{ji}$ ties
each pair. For a convex body (1) wholly inside an enclosure (2),
every ray leaving the inner surface lands on the enclosure, so
$F_{12} = 1$ exactly, and reciprocity gives $F_{21} = A_1/A_2$ while the
summation rule fixes the enclosure self-view at
$F_{22} = 1 - A_1/A_2$ (Figure~\ref{fig:vol04ch05:concentric-cylinders}).
Feeding $F_{12} = 1$ into the general exchange formula collapses the
denominator and leaves
\[
\dot{Q}_{1 \to 2} = \frac{A_1 \sigma (T_1^4 - T_2^4)}{
\dfrac{1}{\epsilon_1} + \dfrac{A_1}{A_2}\!\left(\dfrac{1}{\epsilon_2} -
1\right)},
\]
the working form for a satellite radiating to space, a hot tank inside
a chamber, or a pipe inside a duct. Two limits are worth holding. When
the enclosure is much larger than the body, $A_1/A_2 \to 0$, the second
term vanishes and the result returns to the small-body expression
$\dot{Q} = \epsilon_1 A_1 \sigma(T_1^4 - T_2^4)$, the enclosure
emissivity dropping out because the body's own radiation never returns.
When the two areas are equal, $A_1 = A_2$, the denominator becomes
$1/\epsilon_1 + 1/\epsilon_2 - 1$ and the formula returns to the
parallel-plate result. The same area-ratio factor governs concentric
cylinders and concentric spheres, with $A_1/A_2 = r_1/r_2$ for cylinders
and $(r_1/r_2)^2$ for spheres, so a small hot inner surface inside a
large cold outer one approaches the small-body limit while two close
nested surfaces approach the parallel-plate limit.

% Figure: the concentric-cylinder (or sphere-in-enclosure) radiation
% geometry used to deepen the two-surface exchange algebra. An inner
% convex surface (1) sees only the outer surface (2), so F12 = 1; the
% outer surface sees both the inner body and itself, so F22 = 1 - A1/A2.
% Drawn as two nested circles with rays to make the self-view explicit.
\begin{figure}[htbp]
\centering
\begin{tikzpicture}[scale=1.0]
% outer enclosure
\draw[engblue, very thick] (0,0) circle (2.6);
% inner convex body
\draw[engred, very thick, fill=engred!10] (0,0) circle (1.0);
\node[font=\small, engred] at (0,0) {$A_1,\ T_1$};
\node[font=\small, engblue] at (0,2.05) {$A_2,\ T_2$};
% a ray leaving the inner body, reaching the outer surface (F12 = 1)
\draw[engblack, ->, thick] (0.71,0.71) -- (1.84,1.84);
\node[font=\scriptsize, engblack, anchor=west] at (1.5,1.3) {$F_{12}=1$};
% a ray leaving the outer surface and returning to the outer surface
% (the self-view F22), shown as an arc-chord that misses the inner body
\draw[enggray, ->, thick] (-2.45,0.86) to[bend right=35] (-0.86,-2.45);
\node[font=\scriptsize, enggray, anchor=east] at (-2.55,-1.0)
  {$F_{22}=1-\dfrac{A_1}{A_2}$};
% separation gap label
\draw[englime!50!black, <->] (1.0,-0.0) -- (2.6,-0.0);
\node[font=\scriptsize, englime!50!black, anchor=north] at (1.8,-0.1)
  {evacuated gap};
\end{tikzpicture}
\caption[Concentric-body radiation geometry and its view
factors]{The radiation geometry of a convex body inside an enclosure,
the case the two-surface exchange algebra reduces to most often. Every
ray leaving the inner convex surface lands on the enclosure, so the view
factor from the inner body to the enclosure is exactly one. The enclosure
sees both the inner body and itself, and the summation rule fixes its
self-view factor at $F_{22}=1-A_1/A_2$ once reciprocity gives
$F_{21}=A_1/A_2$. Substituting these into the general two-surface formula
collapses it to a single expression in the area ratio and the two
emissivities, the working form for a tank in a chamber, a satellite in
space, and a wire in a duct.}
\label{fig:vol04ch05:concentric-cylinders}
\end{figure}


\engterm{The linearised radiation coefficient}. For small temperature
differences, the radiation flux is approximately linear in $\Delta T$:
\[
\dot{Q} \approx h_{\text{rad}} A (T_1 - T_2), \quad h_{\text{rad}}
\approx 4 \epsilon \sigma T_{\text{avg}}^3,
\]
with $T_{\text{avg}}$ the average of $T_1$ and $T_2$. At room
temperature ($T_{\text{avg}} = 300\,\si{\kelvin}$, $\epsilon = 0.9$),
$h_{\text{rad}} \approx 5.5\,\text{W}/(\text{m}^2\cdot\si{\kelvin})$,
comparable to free-convection coefficients in still air. Radiation
matters at room temperature whenever still-air convection is the
dominant alternative; it is the principal heat-loss mode from human
skin, from low-velocity radiators, and from cold-night terrestrial
surfaces to space.

\section{Combined heat transfer}\pathtag{core}

Real engineering surfaces transfer heat by several modes
simultaneously. The combined-resistance machinery generalises the
electrical-analog of section 5.1.

\engterm{Overall heat-transfer coefficient}. For a surface separating
two fluids, with convection on both sides and conduction through the
wall, the overall coefficient $U$ is defined by
\[
\dot{Q} = U A \Delta T_{\text{total}},
\]
with $\Delta T_{\text{total}} = T_{\infty,1} - T_{\infty,2}$ the bulk-
fluid temperature difference. The $U$ satisfies
\[
\frac{1}{U A} = \frac{1}{h_1 A_1} + \frac{L_w}{k_w A_w} + \frac{1}{h_2
A_2},
\]
for a flat wall with two convection coefficients $h_1, h_2$ and a wall
of thickness $L_w$ and conductivity $k_w$. For a tube, the formula
modifies to use the cylindrical-conduction resistance.

\engterm{Building U-values}. Building thermal performance is
specified by the U-value (overall heat-transfer coefficient) of the
envelope components. Typical values: single-pane window, $U \sim 5\,
\text{W}/(\text{m}^2\cdot\si{\kelvin})$; double-pane insulated glass,
$U \sim 2.5$; triple-pane low-emissivity glass, $U \sim 0.8$;
uninsulated brick wall, $U \sim 1.5$; well-insulated wall ($150\,
\si{\milli\meter}$ of polyurethane), $U \sim 0.2$; passive-house wall,
$U \sim 0.1$. The U-value enters the heat-loss calculation $\dot{Q} =
UA \Delta T$ for each component, and the total heat-loss rate of a
building is the sum over all envelope components plus the air-
infiltration contribution.

The two air films that bracket every envelope component are not left to the
designer to correlate; building standards fix them as \engterm{standardised
surface resistances} so that two engineers computing the same wall arrive at
the same U-value. The convention writes the total as $R = R_{si} + \sum_j
L_j/k_j + R_{se}$, with $R_{si}$ the inside surface resistance and $R_{se}$
the outside one, and the tabulated values fold the free-convection and
radiation films into a single number: a horizontal heat flow through a wall
takes $R_{si} = 0.13$ and $R_{se} = 0.04\,\si{\kelvin}\cdot\text{m}^2/\text{W}$,
upward flow through a roof takes a smaller $R_{si} = 0.10$ because the warm air
rises freely off the ceiling, and downward flow through a floor takes a larger
$R_{si} = 0.17$ because the still layer under the surface resists the buoyant
escape. The outside film is small and nearly fixed because wind keeps the
outer convection coefficient high, so $R_{se}$ barely moves; the inside film
is the larger and the one that shifts with flow direction. The reason the
films are standardised rather than computed is that they combine a free-
convection coefficient that depends weakly on the temperature difference with
a radiation coefficient of comparable size, both of which the linearised
$h_{\text{rad}} = 4\epsilon\sigma T^3$ estimate of section 5.4 recovers to
within the tolerance a building calculation needs; pinning them as constants
removes an argument over a term that changes the answer by a few percent. The
practical consequence is that a well-insulated component has almost all its
resistance in the insulation and the films are a rounding term, while a
single-glazed window or a bare metal panel has so little fabric resistance
that the two films set the U-value outright, which is why the achievable
U-value of any glazing is bounded below by the sum of its surface resistances
no matter how the panes are built.

\engterm{Fins (extended surfaces)}. When the convection resistance
$1/(hA)$ is the dominant thermal resistance, attaching fins to the
surface extends the area and reduces the resistance. The
effectiveness of a fin depends on the ratio of conduction along the
fin to convection from the fin's surface, captured by the
dimensionless parameter $mL = L\sqrt{hP/(kA_c)}$, with $P$ the fin
perimeter and $A_c$ its cross-section. A long thin fin ($mL > 5$)
acts approximately as $kA_c m \tanh(mL) \approx kA_c m$; a short fat
fin ($mL < 1$) acts as if its temperature equals the base
temperature. The fin design is the working tool for electronics
cooling, automotive radiators, and air-conditioner condensers.

% Figure: temperature distribution along a convecting fin, as excess
% temperature theta/theta_b = cosh(m(L-x))/cosh(mL) for several values of
% the fin parameter mL. Larger mL means the fin cools off faster and the
% tip approaches ambient.
\begin{figure}[htbp]
\centering
\begin{tikzpicture}
\begin{axis}[
  width=12cm, height=7.5cm,
  xlabel={normalised position $x/L$},
  ylabel={excess temperature $\theta/\theta_b$},
  xmin=0, xmax=1, ymin=0, ymax=1.02,
  axis lines=left,
  legend pos=south west,
  legend cell align=left,
  tick label style={font=\scriptsize},
  label style={font=\small},
  legend style={font=\scriptsize},
  domain=0:1, samples=80,
  clip=false,
]
% theta/theta_b = cosh(mL*(1-s))/cosh(mL), s=x/L, written via exp() for
% robust pgfmath evaluation: cosh(z) = (exp(z)+exp(-z))/2.
\addplot[engblue, very thick]
  {(exp(0.5*(1-x))+exp(-0.5*(1-x)))/(exp(0.5)+exp(-0.5))};
\addlegendentry{$mL=0.5$ (short, fat fin)}
\addplot[englime!70!black, very thick]
  {(exp(1.0*(1-x))+exp(-1.0*(1-x)))/(exp(1.0)+exp(-1.0))};
\addlegendentry{$mL=1$}
\addplot[engamber, very thick]
  {(exp(2.0*(1-x))+exp(-2.0*(1-x)))/(exp(2.0)+exp(-2.0))};
\addlegendentry{$mL=2$}
\addplot[engred, very thick]
  {(exp(4.0*(1-x))+exp(-4.0*(1-x)))/(exp(4.0)+exp(-4.0))};
\addlegendentry{$mL=4$ (long, slender fin)}
% base and tip annotations
\node[font=\scriptsize, anchor=west] at (axis cs:0.02,0.98) {base ($\theta=\theta_b$)};
\node[font=\scriptsize, anchor=east] at (axis cs:0.98,0.10) {tip};
\end{axis}
\end{tikzpicture}
\caption[Temperature distribution along a convecting fin]{Temperature
along an insulated-tip fin, plotted as excess temperature
$\theta/\theta_b = \cosh\!\big(m(L-x)\big)/\cosh(mL)$ against normalised
position, for several values of the fin parameter $mL = L\sqrt{hP/(kA_c)}$.
A short, conductive fin ($mL=0.5$) stays close to the base temperature
along its whole length and acts almost as added base area. A long,
slender fin ($mL=4$) cools to near ambient well before the tip, so its
outer length contributes little and the extra material is wasted. The
design target is an $mL$ near unity, where the fin runs warm enough to
move heat yet long enough to add useful area.}
\label{fig:vol04ch05:fin-temperature}
\end{figure}


Figure~\ref{fig:vol04ch05:fin-temperature} plots the fin temperature for
several values of $mL$, showing why the slender fin wastes its outer
length: it has cooled to near ambient before the tip. The
\texttt{fin\_analysis.py} module computes the tip temperature, heat
rate, and efficiency for a straight fin. Building U-values for the
combined-resistance examples are tabulated in \texttt{wall\_uvalues.csv}.

\begin{mastery}
The fin temperature, heat rate, efficiency, and effectiveness all follow
from one ordinary differential equation. An energy balance on a slice of
a constant-cross-section fin balances the conducted gradient against the
convective loss from the slice surface, giving, for the excess
temperature $\theta = T - T_\infty$,
\[
\frac{d^2\theta}{dx^2} - m^2 \theta = 0, \quad m^2 = \frac{hP}{k A_c},
\]
a linear equation with solution $\theta = C_1 e^{mx} + C_2 e^{-mx}$. The
base condition $\theta(0) = \theta_b$ and an insulated tip
$\theta'(L) = 0$ fix the constants and collapse the solution to
$\theta(x)/\theta_b = \cosh\!\big(m(L-x)\big)/\cosh(mL)$, the curve of
Figure~\ref{fig:vol04ch05:fin-temperature}. The fin heat rate is the
conduction into the base, $\dot{Q}_f = -k A_c \theta'(0) = \sqrt{h P k
A_c}\,\theta_b \tanh(mL)$. Two ratios make the result a design tool. The
\engterm{fin efficiency} compares the actual heat rate to the rate an
ideal fin held entirely at the base temperature would dissipate over the
same surface, $\eta_f = \dot{Q}_f/(h P L \theta_b) = \tanh(mL)/(mL)$,
which falls from unity for a short conductive fin toward zero for a long
slender one. The \engterm{fin effectiveness} compares the finned heat
rate to the bare-base rate the fin's footprint would have dissipated
without it, $\varepsilon_f = \dot{Q}_f/(h A_c \theta_b) =
\sqrt{kP/(hA_c)}\tanh(mL)$. Effectiveness above unity is the condition
for a fin to be worth adding at all, and it grows with the conductivity
ratio $kP/(hA_c)$, which is why fins pay off most on the low-$h$ gas side
of a heat exchanger and why they are made of aluminium or copper rather
than steel. An array's overall surface efficiency then blends the finned
and unfinned area, $\eta_o = 1 - (A_f/A_t)(1 - \eta_f)$, the factor that
turned the bare fin area into effective area in the processor case study.
\end{mastery}

\section{Worked examples: insulated pipe, electronics cooling, building heat loss}\pathtag{standard}

\subsection{Heat loss from an insulated steam pipe}

A steel pipe of internal diameter $50\,\si{\milli\meter}$, wall
thickness $5\,\si{\milli\meter}$, carries saturated steam at $150\,\degC$
(internal convection coefficient $h_i = 5000\,\text{W}/(\text{m}^2
\cdot\si{\kelvin})$). The pipe is in a $20\,\degC$ environment with
$h_o = 8\,\text{W}/(\text{m}^2\cdot\si{\kelvin})$. Insulation of
thermal conductivity $k = 0.04\,\text{W}/(\text{m}\cdot\si{\kelvin})$
is added. Find the heat-loss rate per metre of pipe for insulation
thicknesses of $0$, $25$, $50$, and $100\,\si{\milli\meter}$.

For the cylindrical geometry, the resistances per metre are:
\begin{itemize}[leftmargin=*]
  \item Internal convection: $R_i = 1/(h_i \pi D_i) = 1/(5000 \times
    \pi \times 0.05) = 1.27 \times 10^{-3}\,\si{\kelvin}\cdot\text{m}/
    \text{W}$.
  \item Pipe wall: $R_w = \ln(D_o/D_i)/(2\pi k_w) = \ln(0.060/0.050)
    /(2\pi \times 50) = 5.8 \times 10^{-4}\,\si{\kelvin}\cdot\text{m}/
    \text{W}$.
  \item Insulation: $R_{\text{ins}}(t) = \ln((D_o + 2t)/D_o)/(2\pi k)
    = \ln((60 + 2t)/60)/(2\pi \times 0.04)$.
  \item External convection: $R_o(t) = 1/(h_o \pi (D_o + 2t)) = 1/(8
    \times \pi \times (D_o + 2t))$.
\end{itemize}

Tabulating for the four cases (with the pipe-wall and internal
convection resistances negligible compared to the others):

\begin{center}
\begin{tabular}{ccccc}
\toprule
$t$ (mm) & $R_{\text{ins}}$ & $R_o$ & $R_{\text{tot}}$ & $\dot{Q}/L$ \\
& (K$\cdot$m/W) & (K$\cdot$m/W) & (K$\cdot$m/W) & (W/m) \\
\midrule
$0$ & $0$ & $0.66$ & $0.66$ & $197$ \\
$25$ & $2.41$ & $0.36$ & $2.77$ & $47$ \\
$50$ & $3.90$ & $0.25$ & $4.15$ & $31$ \\
$100$ & $5.83$ & $0.15$ & $5.98$ & $22$ \\
\bottomrule
\end{tabular}
\end{center}

% Figure: heat loss per metre from the insulated steam pipe versus
% insulation thickness, the diminishing-returns curve from the worked
% example. Data points at 0, 25, 50, 100 mm from the chapter table, with a
% smooth interpolating curve. Annotates the diminishing marginal benefit.
\begin{figure}[htbp]
\centering
\begin{tikzpicture}
\begin{axis}[
  width=12cm, height=7.5cm,
  xlabel={insulation thickness $t$ (mm)},
  ylabel={heat loss per metre $\dot{Q}/L$ (W/m)},
  xmin=0, xmax=110, ymin=0, ymax=210,
  axis lines=left,
  legend pos=north east,
  legend cell align=left,
  tick label style={font=\scriptsize},
  label style={font=\small},
  legend style={font=\scriptsize},
  clip=false,
]
% chapter table values
\addplot[engblue, very thick, mark=*, mark options={fill=engblue}]
  coordinates {(0,197) (25,71) (50,49) (100,32)};
\addlegendentry{$\dot{Q}/L$ vs thickness}
% annotate the steep first increment and the flat tail
\draw[engred, ->, thick] (axis cs:6,150) -- (axis cs:6,90);
\node[font=\scriptsize, engred, anchor=west] at (axis cs:8,160)
  {first $25$ mm removes $64\%$ of the loss};
\draw[enggray, dashed] (axis cs:50,0) -- (axis cs:50,49);
\node[font=\scriptsize, enggray, anchor=south west] at (axis cs:52,52)
  {diminishing returns};
\end{axis}
\end{tikzpicture}
\caption[Heat loss versus insulation thickness for the steam pipe]{Heat
loss per metre from the insulated steam pipe of the worked example, as a
function of insulation thickness. The first $25\,$mm of insulation
removes nearly two thirds of the bare-pipe loss; each further increment
removes less, because the added conductive resistance is a shrinking
fraction of a total that already includes the earlier layers and the
outer-film resistance. The economic thickness sits where the marginal
value of the heat saved over the system lifetime equals the marginal cost
of insulation, well to the left of the point of vanishing returns. For
this pipe the outer radius is far above the critical insulation radius
$r_c=k/h$, so insulation always reduces the loss; on a thin wire the
opposite can hold.}
\label{fig:vol04ch05:insulation-thickness}
\end{figure}


The heat-loss rate falls from $197\,\text{W}/\text{m}$ uninsulated to
$22\,\text{W}/\text{m}$ with $100\,\si{\milli\meter}$ of insulation, an
$89\,\%$ reduction, plotted in
Figure~\ref{fig:vol04ch05:insulation-thickness}. The marginal
effectiveness diminishes as
insulation is added; the engineering choice of thickness is set by
the cost of insulation balanced against the cost of the heat lost
over the system's lifetime, plus practical constraints (pipe spacing,
weight, aesthetics).

A subtle point: for small radii, adding insulation can \emph{increase}
the heat loss, because the increase in outer surface area for
convection can exceed the conductive resistance of the insulation. The
\engterm{critical insulation radius} $r_c = k/h$ is the radius below
which insulation increases heat loss. For the example above, $r_c =
0.04/8 = 0.005\,\text{m} = 5\,\si{\milli\meter}$, much smaller than
the pipe's outer radius of $30\,\si{\milli\meter}$, so insulation
unambiguously reduces heat loss. For small-diameter wires (electrical
insulation on copper wire), the critical radius can exceed the
actual radius and the insulation actually increases the wire's heat
dissipation.

\subsection{Electronics cooling: a microprocessor heat sink}

A microprocessor dissipates $50\,\text{W}$ over a $1\,\si{\centi\meter}
\times 1\,\si{\centi\meter}$ die. The die is mounted to an aluminium
heat sink whose base measures $5\,\si{\centi\meter} \times 5\,\si{
\centi\meter}$. A fan blows air at $25\,\degC$ over the heat-sink
fins. Specify the heat-sink fin geometry to keep the die temperature
below $85\,\degC$.

The die-to-base temperature drop comes from the thermal-interface
material (TIM) and the spreading resistance in the heat-sink base.
For a TIM resistance of $R_{\text{TIM}} = 0.1\,\si{\kelvin}/\text{W}$
(typical for a thermal-grease interface), the die-to-base drop is
$50 \times 0.1 = 5\,\si{\kelvin}$. The spreading resistance from the
$1\,\si{\centi\meter}^2$ die to the $5\,\si{\centi\meter} \times 5\,\si{
\centi\meter}$ heat-sink base is $\sim 0.05\,\si{\kelvin}/\text{W}$ for
aluminium, adding another $2.5\,\si{\kelvin}$. The base temperature is
therefore $\sim 7.5\,\si{\kelvin}$ below the die, so the base must be
at or below $77.5\,\degC$ for the die to be at $85\,\degC$.

The heat-sink-to-ambient temperature drop is $77.5 - 25 = 52.5\,\si{
\kelvin}$. The required heat-sink thermal resistance is $52.5/50 = 1.05
\,\si{\kelvin}/\text{W}$. A heat sink with $50$ fins of $30\,\si{
\milli\meter}$ height, $1\,\si{\milli\meter}$ thickness, with $5\,
\text{m/s}$ air velocity, gives an overall thermal resistance of
$\sim 0.5\,\si{\kelvin}/\text{W}$ (computed from fin theory with the
forced-convection coefficient for the air over the fins), well within
the requirement.

The example illustrates the typical electronics-cooling stack: the
die-to-junction TIM, the spreading resistance in the heat-sink base,
the fin-bundle convection, and the bulk-fluid temperature rise (for
sealed enclosures, the fluid-temperature rise is the limiting
resistance and a larger enclosure or active venting is needed).

\subsection{Building heat loss}

A residential building has the following envelope: $200\,\text{m}^2$
of wall with $U = 0.3\,\text{W}/(\text{m}^2\cdot\si{\kelvin})$, $30\,
\text{m}^2$ of double-pane windows with $U = 2.5$, $100\,\text{m}^2$
of roof with $U = 0.2$, $100\,\text{m}^2$ of slab-on-grade floor with
effective $U = 0.4$. Air infiltration is $0.5$ air-changes per hour
on a heated volume of $250\,\text{m}^3$. Find the heat-loss rate at
an indoor-outdoor temperature difference of $20\,\si{\kelvin}$.

The transmission losses are:
\begin{itemize}[leftmargin=*]
  \item Walls: $200 \times 0.3 \times 20 = 1200\,\text{W}$.
  \item Windows: $30 \times 2.5 \times 20 = 1500\,\text{W}$.
  \item Roof: $100 \times 0.2 \times 20 = 400\,\text{W}$.
  \item Floor: $100 \times 0.4 \times 20 = 800\,\text{W}$.
\end{itemize}
Total transmission: $3900\,\text{W}$.

The infiltration loss is
\[
\dot{Q}_{\text{inf}} = \dot{V} \rho c_p \Delta T = (0.5 \times 250/3600)
\times 1.2 \times 1005 \times 20 \approx 840\,\text{W}.
\]

The total heat-loss rate is $4740\,\text{W}$ at $\Delta T = 20\,\si{
\kelvin}$, or $237\,\text{W}/\si{\kelvin}$. At the design outdoor
condition of $-15\,\degC$ (with indoor $20\,\degC$), $\Delta T = 35\,
\si{\kelvin}$ and the design heat-loss rate is $8300\,\text{W}$. The
windows (with their high U-value) contribute $32\,\%$ of the heat loss
on $13\,\%$ of the envelope area; upgrading the glazing is the highest-
leverage improvement.

\subsection{Quenching a steel shaft: lumped capacitance with a Biot
check}

A cylindrical steel shaft of diameter $40\,\si{\milli\meter}$ and
length $300\,\si{\milli\meter}$ leaves a furnace at $820\,\degC$ and is
quenched in an agitated oil bath at $60\,\degC$ with a convection
coefficient $h = 600\,\text{W}/(\text{m}^2\cdot\si{\kelvin})$. Take
steel properties $k = 40\,\text{W}/(\text{m}\cdot\si{\kelvin})$, $\rho =
7850\,\text{kg/m}^3$, $c_p = 480\,\text{J}/(\text{kg}\cdot\si{\kelvin})$.
We want the time to reach a core temperature of $200\,\degC$, and we
want to know first whether the lumped-capacitance model is even allowed.

The characteristic length for a body is its volume over its surface
area. For the finite cylinder,
\begin{align*}
V &= \pi r^2 L = \pi (0.020)^2 (0.300) = 3.77 \times 10^{-4}\,\text{m}^3,
\\
A_s &= 2\pi r L + 2\pi r^2 = 0.0377 + 0.00251 = 0.0402\,\text{m}^2,
\end{align*}
so $L_c = V/A_s = 9.38 \times 10^{-3}\,\text{m}$. The Biot number is
\[
\text{Bi} = \frac{h L_c}{k} = \frac{600 \times 9.38 \times 10^{-3}}{40}
= 0.141.
\]
This sits just above the $0.1$ threshold, so the lumped model carries a
real but modest error: the core lags the surface by a temperature drop
of order $\text{Bi}$ times the surface-to-bath difference, a few percent
of the total swing. We proceed with the lumped estimate and flag that it
runs slightly optimistic on the core temperature. The time constant is
\[
\tau = \frac{\rho V c_p}{h A_s} = \frac{\rho L_c c_p}{h} =
\frac{7850 \times 9.38 \times 10^{-3} \times 480}{600} = 58.9\,\text{s}.
\]
The cooling follows $(T - T_\infty)/(T_0 - T_\infty) = e^{-t/\tau}$. The
target ratio is $(200 - 60)/(820 - 60) = 140/760 = 0.184$, so
\[
t = \tau \ln\!\left(\frac{1}{0.184}\right) = 58.9 \times 1.69 =
99.6\,\text{s},
\]
about a minute and forty seconds. Because $\text{Bi}$ is above $0.1$, a
one-term series solution for the finite cylinder would give a core that
reaches $200\,\degC$ a little later than this, on the order of ten to
twenty percent longer; the lumped estimate is the right size and the
right tool for sizing the quench tank and the handling time, with the
series solution reserved for the metallurgical question of whether the
core and the surface cross the transformation nose together.

\subsection{A multilayer wall with contact resistance}

The series-resistance formula assumes each layer touches the next with
no temperature jump at the interface. Real joined surfaces touch only at
their high spots, so heat crowds through the actual contact points and a
measurable temperature drop appears across the joint, captured by a
\engterm{thermal contact resistance} $R''_c$ (units
$\si{\kelvin}\cdot\text{m}^2/\text{W}$) that adds in series like any
other layer. We take a wall built of a $6\,\si{\milli\meter}$ aluminium
plate ($k = 237$) bolted to a $20\,\si{\milli\meter}$ steel plate ($k =
50$), carrying a flux from a $200\,\degC$ inner gas through inner film
$h_i = 50\,\text{W}/(\text{m}^2\cdot\si{\kelvin})$ to $25\,\degC$ outer
air through outer film $h_o = 15\,\text{W}/(\text{m}^2\cdot\si{\kelvin})$.
The bolted aluminium-to-steel joint has a contact resistance $R''_c =
2.5 \times 10^{-4}\,\si{\kelvin}\cdot\text{m}^2/\text{W}$, typical of a
moderately loaded metal-to-metal joint in air.

Per unit area, the resistances in series are: inner film $1/50 = 0.0200$;
aluminium $0.006/237 = 2.5 \times 10^{-5}$; the contact joint $2.5
\times 10^{-4}$; steel $0.020/50 = 4.0 \times 10^{-4}$; outer film
$1/15 = 0.0667$, all in $\si{\kelvin}\cdot\text{m}^2/\text{W}$. The
total is $R'' = 0.0873$, so the flux is
\[
q = \frac{\Delta T}{R''} = \frac{200 - 25}{0.0873} = 2005\,\text{W}/
\text{m}^2.
\]
The instructive part is where the drops land. The two films carry
$q \times 0.0200 = 40.1\,\si{\kelvin}$ and $q \times 0.0667 =
133.7\,\si{\kelvin}$, together almost the whole $175\,\si{\kelvin}$. The
contact joint drops $q \times 2.5 \times 10^{-4} = 0.50\,\si{\kelvin}$,
the steel $0.80\,\si{\kelvin}$, the aluminium $0.05\,\si{\kelvin}$. The
contact resistance is real and would matter in a high-flux package
(a power module, where it can dwarf the metal layers), but against two
gas films it is a rounding term. The lesson the numbers teach is the one
the chapter's principle states: rank the resistances before redesigning,
and here the design lever is the outer film, not the joint. Had the same
joint sat between two water-cooled plates with $h \sim 5000$, the films
would shrink to $2 \times 10^{-4}$ each and the $2.5 \times 10^{-4}$
contact resistance would become the controlling term.

\subsection{A radiation shield between hot plates}

Two large parallel plates at $T_1 = 800\,\si{\kelvin}$ and $T_2 =
400\,\si{\kelvin}$, each with emissivity $\epsilon = 0.8$, exchange heat
by radiation across an evacuated gap. We compute the net flux bare, then
insert a single thin shield of emissivity $\epsilon_s = 0.05$ on both
faces and recompute, to size the reduction a polished shield buys
(Figure~\ref{fig:vol04ch05:radiation-shield}).

% Figure: radiation shielding between two parallel plates. The left panel
% shows the bare two-plate exchange; the right shows one low-emissivity
% shield interposed, halving the flux for equal-emissivity surfaces and
% far more for a polished shield. Schematic plates with arrows whose
% width suggests the flux magnitude.
\begin{figure}[htbp]
\centering
\begin{tikzpicture}[every node/.style={font=\small}]
% ==== left panel: no shield ====
\begin{scope}[shift={(0,0)}]
\draw[engred, very thick, fill=engred!12] (0,0) rectangle (0.3,3);
\draw[engblue, very thick, fill=engblue!12] (2.4,0) rectangle (2.7,3);
\node[font=\scriptsize, engred, anchor=south] at (0.15,3.05) {$T_1$};
\node[font=\scriptsize, engblue, anchor=south] at (2.55,3.05) {$T_2$};
% thick flux arrow
\draw[engamber, line width=3pt, ->] (0.45,1.5) -- (2.25,1.5);
\node[font=\scriptsize, engamber, anchor=south] at (1.35,1.7) {$\dot{q}_0$};
\node[font=\scriptsize, anchor=north] at (1.35,-0.15) {no shield};
\end{scope}
% ==== right panel: one shield ====
\begin{scope}[shift={(5.5,0)}]
\draw[engred, very thick, fill=engred!12] (0,0) rectangle (0.3,3);
\draw[enggray, very thick, fill=enggray!20] (1.35,0) rectangle (1.55,3);
\draw[engblue, very thick, fill=engblue!12] (2.6,0) rectangle (2.9,3);
\node[font=\scriptsize, engred, anchor=south] at (0.15,3.05) {$T_1$};
\node[font=\scriptsize, enggray, anchor=south] at (1.45,3.05) {$T_s$};
\node[font=\scriptsize, engblue, anchor=south] at (2.75,3.05) {$T_2$};
% thinner flux arrows
\draw[engamber, line width=1.5pt, ->] (0.45,1.5) -- (1.2,1.5);
\draw[engamber, line width=1.5pt, ->] (1.7,1.5) -- (2.45,1.5);
\node[font=\scriptsize, engamber, anchor=south] at (1.45,1.7) {$\dot{q}_1$};
\node[font=\scriptsize, anchor=north] at (1.45,-0.15) {one shield};
\end{scope}
% relation between the two
\node[font=\scriptsize, engblack, anchor=west] at (9.0,1.5)
  {$\dot{q}_1=\tfrac12\,\dot{q}_0$};
\node[font=\scriptsize, enggray, anchor=west] at (9.0,0.9)
  {(equal $\epsilon$)};
\end{tikzpicture}
\caption[Radiation shielding between parallel plates]{Radiation
shielding between two parallel plates. With no shield (left) the net
flux $\dot{q}_0$ crosses the gap set by the two surface emissivities.
Interposing a single thin shield (right) splits the path into two
exchanges in series, each with its own emissivity resistance, so the
shield floats at an intermediate temperature $T_s$ and the net flux
falls. For a shield whose two faces share the emissivity of the plates,
one shield halves the flux and $N$ shields cut it by a factor $N+1$. A
polished shield of low emissivity does far better, because each shield
face adds a large emissivity resistance $1/\epsilon-1$ to the chain,
which is why multilayer insulation stacks dozens of aluminised films to
reach the very low fluxes needed in cryogenic and spacecraft thermal
control.}
\label{fig:vol04ch05:radiation-shield}
\end{figure}


Bare, the parallel-plate formula gives
\[
\dot{q}_0 = \frac{\sigma (T_1^4 - T_2^4)}{1/\epsilon_1 + 1/\epsilon_2 -
1} = \frac{5.67 \times 10^{-8} (800^4 - 400^4)}{1/0.8 + 1/0.8 - 1}.
\]
With $800^4 = 4.096 \times 10^{11}$ and $400^4 = 2.56 \times 10^{10}$,
the numerator is $5.67 \times 10^{-8} \times 3.84 \times 10^{11} =
2.18 \times 10^4$, and the denominator is $1.5$, so $\dot{q}_0 = 1.45
\times 10^4\,\text{W}/\text{m}^2$. With the shield, the path becomes two
gaps in series, plate~1 to shield and shield to plate~2, each carrying
the same net flux at steady state. The radiative resistance per unit
area of a gap between surfaces $a$ and $b$ is $1/\epsilon_a + 1/\epsilon_b
- 1$. The two gaps contribute
\[
\mathcal{R}_{1s} = \frac{1}{0.8} + \frac{1}{0.05} - 1 = 20.25, \quad
\mathcal{R}_{s2} = \frac{1}{0.05} + \frac{1}{0.8} - 1 = 20.25,
\]
so the total resistance is $40.5$ against the bare $1.5$, and
\[
\dot{q}_1 = \frac{\sigma(T_1^4 - T_2^4)}{40.5} = \frac{2.18 \times
10^4}{40.5} = 538\,\text{W}/\text{m}^2.
\]
The single low-emissivity shield cuts the flux by a factor of $27$, not
the factor of two a same-emissivity shield would give, because each
polished face adds a large resistance $1/\epsilon_s - 1 = 19$ to the
chain. The shield floats at the temperature that makes the two gap
fluxes equal; solving $\sigma(T_1^4 - T_s^4)/\mathcal{R}_{1s} =
\sigma(T_s^4 - T_2^4)/\mathcal{R}_{s2}$ with equal resistances gives
$T_s^4 = (T_1^4 + T_2^4)/2$, so $T_s = 683\,\si{\kelvin}$. Stacking
$N$ such shields multiplies the resistance again, the principle behind
the multilayer insulation that wraps cryogenic tanks and spacecraft.

\subsection{Critical insulation radius on a current-carrying wire}

The critical-radius warning of the steam-pipe example reverses on a thin
conductor, where adding insulation can raise the heat the surface
dissipates rather than lower it. We take a bare copper wire of diameter
$3\,\si{\milli\meter}$ carrying current that dissipates $\dot{Q}/L =
8\,\text{W}/\text{m}$ of Ohmic heat, sitting in still air at $25\,\degC$
with an outer film $h = 12\,\text{W}/(\text{m}^2\cdot\si{\kelvin})$. The
wire is to be coated with a plastic sheath of conductivity $k =
0.20\,\text{W}/(\text{m}\cdot\si{\kelvin})$, and we ask whether the
sheath makes the wire run hotter or cooler.

The critical radius is $r_c = k/h = 0.20/12 = 0.0167\,\text{m} =
16.7\,\si{\milli\meter}$, far larger than the bare wire radius of
$1.5\,\si{\milli\meter}$. Adding insulation from the bare radius outward
toward $r_c$ lowers the total outer resistance, because the gain in
convecting surface area outruns the small conductive resistance of a thin
plastic layer, so the wire runs cooler for a fixed dissipation. With the
bare wire, the only outer resistance per metre is the film,
\[
R_{\text{bare}} = \frac{1}{h\,\pi D} = \frac{1}{12 \times \pi \times
0.003} = 8.84\,\si{\kelvin}\cdot\text{m}/\text{W},
\]
so the bare wire sits at $\Delta T = (\dot{Q}/L)\,R_{\text{bare}} = 8
\times 8.84 = 70.7\,\si{\kelvin}$ above air, a surface near $96\,\degC$.
Coat it to the critical radius $r_c = 16.7\,\si{\milli\meter}$ (a sheath
$15.2\,\si{\milli\meter}$ thick, larger than any real wire would carry,
shown here to locate the minimum). The two outer resistances per metre
are the sheath conduction and the film on the enlarged surface,
\[
R(r_c) = \frac{\ln(r_c/r_i)}{2\pi k} + \frac{1}{2\pi r_c h} =
\frac{\ln(0.0167/0.0015)}{2\pi \times 0.20} + \frac{1}{2\pi \times
0.0167 \times 12},
\]
which evaluates to $1.92 + 0.795 = 2.71\,\si{\kelvin}\cdot\text{m}/
\text{W}$, well below the bare $8.84$. At the critical radius the wire
sits only $8 \times 2.71 = 21.7\,\si{\kelvin}$ above air, near
$47\,\degC$. The thin-conductor result is the design rule behind the
ampacity gain that a sheath gives a small wire and behind the warning
that the same plastic on a fat busbar, whose radius already exceeds
$r_c$, only insulates and raises the metal temperature. The sign of the
effect is read off one comparison: whether the starting radius is below
or above $k/h$.

\subsection{Frost depth and a buried water line}

The semi-infinite-solid result sizes the depth at which a buried pipe
escapes the seasonal frost. We take ground initially at $+8\,\degC$ whose
surface is held at $-10\,\degC$ through a cold spell, soil diffusivity
$\alpha = 5 \times 10^{-7}\,\text{m}^2/\text{s}$, and ask how deep the
$0\,\degC$ frost line reaches after thirty and after ninety days, then
where a water line must sit to stay unfrozen.

The error-function profile of section 5.2 gives the temperature at depth
$x$ and time $t$,
\[
\frac{T(x,t) - T_s}{T_0 - T_s} = \text{erf}\!\left(\frac{x}{2\sqrt{\alpha
t}}\right),
\]
with $T_s = -10\,\degC$ the surface and $T_0 = +8\,\degC$ the deep ground.
The frost line is the depth where $T = 0\,\degC$, so the left side is
$(0 - (-10))/(8 - (-10)) = 10/18 = 0.556$, and the frost depth solves
$\text{erf}(\eta) = 0.556$ with $\eta = x_f/(2\sqrt{\alpha t})$. Inverting
the error function, $\eta = 0.539$, so
\[
x_f = 0.539 \times 2\sqrt{\alpha t} = 1.078\sqrt{\alpha t}.
\]
At thirty days, $\alpha t = 5\times10^{-7} \times 2.59\times10^6 = 1.30\,
\text{m}^2$, so $\sqrt{\alpha t} = 1.14\,\text{m}$ and $x_f = 1.23\,
\text{m}$. At ninety days, $\alpha t = 3.89\,\text{m}^2$, $\sqrt{\alpha t}
= 1.97\,\text{m}$, and $x_f = 2.13\,\text{m}$ for this severe, unbroken
cold spell; real climates relieve the surface before the frost reaches
that depth, which is why published frost lines for temperate regions run
nearer $1\,\text{m}$. The profiles deepen as the square root of time
(Figure~\ref{fig:vol04ch05:semi-infinite-freeze}), so doubling the frost
depth costs four times the cold duration, and a water line buried below
the local code frost depth, with a margin, stays above $0\,\degC$ through
the season. The daily surface swing penetrates with the same square-root
law on a far shorter clock, a depth $\sqrt{\alpha \times 86400} \approx
0.21\,\text{m}$, so it never reaches the buried line and only the seasonal
mean and the long cold spell matter for frost.

% Figure: semi-infinite-solid temperature profiles in soil after a step to a
% sub-freezing surface, plotted at three elapsed times. The error-function
% profiles deepen with the square root of time; the depth where each curve
% crosses 0 degrees C is the frost penetration, advancing as sqrt(t). Used by
% the buried-pipe / frost-depth worked example. Coordinates are precomputed
% from T(x,t) = Ts + (T0 - Ts) erf( x / (2 sqrt(alpha t)) ) with Ts=-10,
% T0=8, alpha=5e-7 m^2/s at t = 10, 30, 90 days.
\begin{figure}[htbp]
\centering
\begin{tikzpicture}
\begin{axis}[
  width=12cm, height=7.5cm,
  xlabel={depth $x$ (m)},
  ylabel={soil temperature $T$ ($^{\circ}$C)},
  xmin=0, xmax=2.0, ymin=-12, ymax=10,
  axis lines=left,
  legend pos=south east,
  legend cell align=left,
  tick label style={font=\scriptsize},
  label style={font=\small},
  legend style={font=\scriptsize},
  clip=false,
]
\addplot[engblue, very thick] coordinates {
  (0.0,-10.0) (0.1,-8.458) (0.2,-6.935) (0.3,-5.446) (0.4,-4.007)
  (0.5,-2.634) (0.6,-1.338) (0.7,-0.128) (0.8,0.987) (0.9,2.004)
  (1.0,2.921) (1.1,3.737) (1.2,4.456) (1.3,5.082) (1.4,5.621)
  (1.5,6.079) (1.6,6.465) (1.7,6.785) (1.8,7.048) (1.9,7.262) (2.0,7.433)};
\addlegendentry{10 days}
\addplot[engamber, very thick] coordinates {
  (0.0,-10.0) (0.1,-9.109) (0.2,-8.221) (0.3,-7.339) (0.4,-6.468)
  (0.5,-5.611) (0.6,-4.769) (0.7,-3.947) (0.8,-3.147) (0.9,-2.371)
  (1.0,-1.622) (1.1,-0.901) (1.2,-0.21) (1.3,0.45) (1.4,1.078)
  (1.5,1.672) (1.6,2.234) (1.7,2.761) (1.8,3.255) (1.9,3.716) (2.0,4.145)};
\addlegendentry{30 days}
\addplot[engred, very thick] coordinates {
  (0.0,-10.0) (0.1,-9.485) (0.2,-8.971) (0.3,-8.458) (0.4,-7.947)
  (0.5,-7.438) (0.6,-6.933) (0.7,-6.432) (0.8,-5.935) (0.9,-5.443)
  (1.0,-4.957) (1.1,-4.477) (1.2,-4.004) (1.3,-3.538) (1.4,-3.08)
  (1.5,-2.63) (1.6,-2.189) (1.7,-1.757) (1.8,-1.334) (1.9,-0.92) (2.0,-0.517)};
\addlegendentry{90 days}
% zero-degree (frost) line
\draw[enggray, dashed] (axis cs:0,0) -- (axis cs:2.0,0);
\node[font=\scriptsize, enggray, anchor=south west] at (axis cs:1.45,0.2)
  {frost line ($0\,^{\circ}$C)};
\end{axis}
\end{tikzpicture}
\caption[Soil temperature profiles after a step to a sub-freezing
surface]{Soil temperature against depth after the surface is stepped to
$-10\,^{\circ}$C over ground initially at $+8\,^{\circ}$C, evaluated at ten,
thirty, and ninety days with the semi-infinite-solid error-function
solution and a soil diffusivity $\alpha = 5\times10^{-7}\,\text{m}^2/$s. The
profiles deepen with the square root of time, and the depth at which each
curve crosses $0\,^{\circ}$C is the frost penetration. The slow advance,
under a metre even after three months of hard cold, is why burying a water
line below the local frost line protects it and why the daily surface swing,
which has a far shorter penetration depth, never reaches that line at all.}
\label{fig:vol04ch05:semi-infinite-freeze}
\end{figure}


\subsection{A three-surface enclosure by the radiosity network}

The two-surface formulas of section 5.4 stop at an enclosure where a
third surface neither holds a fixed temperature nor exchanges net heat.
We take a long enclosure of three surfaces (a heated floor at $T_1 =
400\,\si{\kelvin}$, a cooled ceiling at $T_2 = 300\,\si{\kelvin}$, and a
side wall that is well insulated on its back so it carries no net heat,
a \engterm{re-radiating surface}) and find the floor-to-ceiling heat flow.
Take the floor and ceiling each $1\,\text{m}$ wide per metre of length
with emissivity $\epsilon_1 = \epsilon_2 = 0.8$, and view factors
$F_{12} = F_{21} = 0.4$ between floor and ceiling, with the remainder of
each surface's radiation reaching the wall, $F_{13} = F_{23} = 0.6$.

The radiosity network of Figure~\ref{fig:vol04ch05:enclosure-network}
carries a surface resistance at the floor and the ceiling and a triangle
of space resistances joining the three radiosity nodes. The blackbody
emissive powers are $E_{b1} = \sigma T_1^4 = 5.67\times10^{-8} \times
400^4 = 1451\,\text{W}/\text{m}^2$ and $E_{b2} = 5.67\times10^{-8}
\times 300^4 = 459\,\text{W}/\text{m}^2$. Per metre of length, with
$A_1 = A_2 = 1\,\text{m}^2$, the resistances are
\[
\frac{1-\epsilon_1}{\epsilon_1 A_1} = \frac{0.2}{0.8} = 0.25, \quad
\frac{1}{A_1 F_{12}} = \frac{1}{0.4} = 2.5, \quad \frac{1}{A_1 F_{13}}
= \frac{1}{A_2 F_{23}} = \frac{1}{0.6} = 1.667,
\]
all in $\text{m}^{-2}$ (the floor and ceiling share these by symmetry).
The re-radiating wall has no surface-resistance branch, so its node
$J_3$ floats and only ties the two space resistances $1/(A_1 F_{13})$ and
$1/(A_2 F_{23})$ in series, a path of $1.667 + 1.667 = 3.333$ that runs in
parallel with the direct floor-to-ceiling space resistance of $2.5$. The
two parallel space paths combine to
\[
R_{\text{space}} = \left(\frac{1}{2.5} + \frac{1}{3.333}\right)^{-1} =
(0.40 + 0.30)^{-1} = 1.429,
\]
and the total resistance from $E_{b1}$ to $E_{b2}$ adds the two surface
resistances in series, $0.25 + 1.429 + 0.25 = 1.929$. The net floor heat
flow per metre is the emissive-power difference over this resistance,
\[
\dot{Q}_1 = \frac{E_{b1} - E_{b2}}{1.929} = \frac{1451 - 459}{1.929} =
514\,\text{W}/\text{m}.
\]
The re-radiating wall raises the floor-to-ceiling exchange above the
$1451-459$ over the direct path alone would give, because it intercepts
radiation that would otherwise miss the ceiling and returns part of it,
a contribution the two-surface formula cannot represent. The wall settles
at the radiosity $J_3$ that makes its two space currents equal, which for
equal wall resistances is the average of the two adjacent radiosities, and
its physical temperature follows from $J_3 = \sigma T_3^4$ since a
re-radiating surface emits as a blackbody at its own floating temperature.
The network method scales to any number of surfaces: write one node per
surface, a surface resistance for each surface of known temperature, and a
space resistance for each visible pair, then solve the linear system.

% Figure: the radiosity network for a three-surface enclosure. Each surface
% i carries a surface resistance (1-eps)/(eps A) from its blackbody emissive
% power E_b,i to its radiosity node J_i, and each pair of nodes is joined by a
% space resistance 1/(A_i F_ij). Left: the physical enclosure (a heated floor,
% a cool ceiling, a re-radiating side wall). Right: the resistor network.
\begin{figure}[htbp]
\centering
\begin{tikzpicture}[every node/.style={font=\small}]
% ==== left: physical enclosure ====
\begin{scope}[shift={(0,0)}]
\draw[engred, very thick] (0,0) -- (3,0);
\node[font=\scriptsize, engred, anchor=north] at (1.5,-0.1) {floor (1), hot};
\draw[engblue, very thick] (0,2.4) -- (3,2.4);
\node[font=\scriptsize, engblue, anchor=south] at (1.5,2.5) {ceiling (2), cool};
\draw[enggray, very thick] (3,0) -- (3,2.4);
\draw[enggray, very thick] (0,0) -- (0,2.4);
\node[font=\scriptsize, enggray, anchor=west] at (3.1,1.2) {wall (3)};
\end{scope}
% ==== right: resistor network ====
\begin{scope}[shift={(6.0,0)}]
% nodes
\node[circle, draw=engred, fill=engred!12, inner sep=1.5pt] (Eb1) at (0,0) {};
\node[font=\scriptsize, anchor=north] at (0,-0.15) {$E_{b1}$};
\node[circle, draw=engblack, fill=white, inner sep=1.5pt] (J1) at (1.4,0) {};
\node[font=\scriptsize, anchor=north] at (1.4,-0.15) {$J_1$};
\node[circle, draw=engblue, fill=engblue!12, inner sep=1.5pt] (Eb2) at (4.2,0) {};
\node[font=\scriptsize, anchor=north] at (4.2,-0.15) {$E_{b2}$};
\node[circle, draw=engblack, fill=white, inner sep=1.5pt] (J2) at (2.8,0) {};
\node[font=\scriptsize, anchor=north] at (2.8,-0.15) {$J_2$};
\node[circle, draw=enggray, fill=enggray!15, inner sep=1.5pt] (J3) at (2.1,1.8) {};
\node[font=\scriptsize, anchor=south] at (2.1,1.9) {$J_3$ (wall)};
% surface resistances
\draw[engred, thick] (Eb1) -- (J1)
  node[midway, above, font=\tiny] {$\frac{1-\epsilon_1}{\epsilon_1 A_1}$};
\draw[engblue, thick] (J2) -- (Eb2)
  node[midway, above, font=\tiny] {$\frac{1-\epsilon_2}{\epsilon_2 A_2}$};
% space resistances
\draw[engblack] (J1) -- (J2)
  node[midway, below, font=\tiny] {$\frac{1}{A_1 F_{12}}$};
\draw[enggray] (J1) -- (J3)
  node[midway, left, font=\tiny] {$\frac{1}{A_1 F_{13}}$};
\draw[enggray] (J2) -- (J3)
  node[midway, right, font=\tiny] {$\frac{1}{A_2 F_{23}}$};
\end{scope}
\end{tikzpicture}
\caption[Radiosity network for a three-surface enclosure]{The radiosity
network for a three-surface enclosure (left, a hot floor, a cool ceiling,
and a re-radiating side wall). Each surface carries a surface resistance
$(1-\epsilon_i)/(\epsilon_i A_i)$ from its blackbody emissive power
$E_{bi}=\sigma T_i^4$ to a radiosity node $J_i$, and every pair of nodes is
joined by a space resistance $1/(A_i F_{ij})$ set by the view factor. A
surface that exchanges no net heat (the re-radiating wall) has its
$E_{b}$ branch open, so its node floats and only redistributes radiation,
which collapses the wall to a single node $J_3$ tied to the other two. The
net heat at any surface is its surface-resistance current
$(E_{bi}-J_i)/[(1-\epsilon_i)/(\epsilon_i A_i)]$, and solving the small
node network is the working method for enclosures the two-surface formulas
cannot cover.}
\label{fig:vol04ch05:enclosure-network}
\end{figure}


\begin{mastery}
The enclosure network is the circuit form of the radiosity balance, and
writing the balance directly shows where the resistances come from. The
\engterm{radiosity} $J_i$ is the total radiation leaving surface $i$ per
unit area, the sum of its own emission and the reflected part of the
irradiation $G_i$ falling on it,
\[
J_i = \epsilon_i E_{bi} + (1 - \epsilon_i) G_i,
\]
where the surface is taken grey and diffuse so that reflectivity is
$1 - \epsilon_i$ by Kirchhoff's law. The net heat leaving surface $i$ is
the radiosity minus the irradiation times the area, $\dot{Q}_i =
A_i(J_i - G_i)$. Eliminating $G_i$ between the two relations gives
$\dot{Q}_i = (E_{bi} - J_i)/[(1-\epsilon_i)/(\epsilon_i A_i)]$, which is
exactly a current through the surface resistance $(1-\epsilon_i)/
(\epsilon_i A_i)$ driven by the potential drop from the blackbody emissive
power $E_{bi}$ to the radiosity node $J_i$. The irradiation on $i$ is the
sum of the radiosities of every surface it sees, weighted by view factor,
$A_i G_i = \sum_j A_j F_{ji} J_j = \sum_j A_i F_{ij} J_j$ by reciprocity,
so the net heat is also $\dot{Q}_i = \sum_j A_i F_{ij}(J_i - J_j)$, a sum
of currents through the space resistances $1/(A_i F_{ij})$ between node
$J_i$ and each node $J_j$. A surface of fixed temperature pins its
$E_{bi}$; a re-radiating surface sets $\dot{Q}_i = 0$, which forces
$J_i = G_i$ and leaves the node floating at the potential its space
resistances dictate. The whole apparatus is one linear system in the
radiosities, and the two-surface formula of section 5.4 is the special
case of two nodes, two surface resistances, and a single space resistance
in series.
\end{mastery}

\subsection{Radiative cooling of a small body in shade}

A transient with no fluid to convect into runs on radiation alone, and
the fourth-power loss makes the cooling depart from the clean exponential
of the lumped model. We take a solid aluminium sphere of diameter
$80\,\si{\milli\meter}$, painted to emissivity $\epsilon = 0.85$, that
starts at $T_0 = 350\,\si{\kelvin}$ and radiates to a $4\,\si{\kelvin}$
deep-space sink with no other heat path. Take aluminium $\rho =
2700\,\text{kg/m}^3$, $c_p = 900\,\text{J}/(\text{kg}\cdot\si{\kelvin})$,
and $k = 237\,\text{W}/(\text{m}\cdot\si{\kelvin})$. We want the time to
cool to $250\,\si{\kelvin}$ and a check that the body stays nearly
isothermal so the lumped balance is allowed.

The conduction inside the sphere is enormous against the radiative surface
loss, so the body is isothermal. The radiative Biot check uses a
linearised radiation coefficient at the starting temperature,
$h_{\text{rad}} = 4\epsilon\sigma T_0^3 = 4 \times 0.85 \times 5.67
\times 10^{-8} \times 350^3 = 8.27\,\text{W}/(\text{m}^2\cdot\si{\kelvin})$,
and the sphere characteristic length $L_c = D/6 = 0.0133\,\text{m}$, so
$\text{Bi} = h_{\text{rad}} L_c/k = 8.27 \times 0.0133/237 = 4.6 \times
10^{-4}$, far below $0.1$, and the lumped balance holds. The energy
balance is now nonlinear,
\[
\rho V c_p \frac{dT}{dt} = -\epsilon A \sigma T^4,
\]
with $T_{\text{space}}^4$ dropped as negligible. Writing $V/A = L_c$ and
separating, $dT/T^4 = -(\epsilon\sigma/(\rho c_p L_c))\,dt$, which
integrates to
\[
\frac{1}{3}\!\left(\frac{1}{T^3} - \frac{1}{T_0^3}\right) =
\frac{\epsilon\sigma}{\rho c_p L_c}\,t.
\]
The coefficient is $\epsilon\sigma/(\rho c_p L_c) = 0.85 \times 5.67
\times 10^{-8}/(2700 \times 900 \times 0.0133) = 1.49 \times 10^{-12}\,
\si{\kelvin}^3/\text{s}$. The inverse-cube difference is $1/250^3 -
1/350^3 = 6.40 \times 10^{-8} - 2.33 \times 10^{-8} = 4.07 \times
10^{-8}\,\si{\kelvin}^{-3}$, so
\[
t = \frac{4.07 \times 10^{-8}}{3 \times 1.49 \times 10^{-12}} = 9.1
\times 10^3\,\text{s},
\]
about two and a half hours. The lesson the inverse-cube law carries is
that radiative cooling slows sharply as the body cools: the same sphere
takes far longer to drop the next hundred kelvin than the first, because
the loss falls as $T^4$. A satellite component shadowed by its own
spacecraft cools fast while hot and then lingers, which is why survival
heaters must hold the cold case rather than wait for a runaway chill.

\subsection{An overall fin-array efficiency, aluminium against steel}

The mastery box defined the array overall surface efficiency $\eta_o = 1 -
(A_f/A_t)(1 - \eta_f)$; here we put numbers to it and show why fin
material is chosen on the conductivity ratio rather than on cost alone. We
take an air-cooled array of $40$ rectangular fins, each $25\,\si{
\milli\meter}$ tall, $1.0\,\si{\milli\meter}$ thick, $80\,\si{\milli\meter}$
deep, with a convection coefficient $h = 60\,\text{W}/(\text{m}^2\cdot
\si{\kelvin})$ from a fan, and we compare an aluminium array ($k =
237$) against a mild-steel one ($k = 50$).

Each fin has perimeter $P = 2(0.080 + 0.001) = 0.162\,\text{m}$ and
cross-section $A_c = 0.080 \times 0.001 = 8.0 \times 10^{-5}\,\text{m}^2$.
The fin parameter is $m = \sqrt{hP/(kA_c)}$. For aluminium,
\[
m_{\text{Al}} = \sqrt{\frac{60 \times 0.162}{237 \times 8.0 \times
10^{-5}}} = \sqrt{513} = 22.6\,\text{m}^{-1}, \quad m_{\text{Al}} L =
22.6 \times 0.025 = 0.566,
\]
so $\eta_{f,\text{Al}} = \tanh(0.566)/0.566 = 0.512/0.566 = 0.905$. For
steel,
\[
m_{\text{St}} = \sqrt{\frac{60 \times 0.162}{50 \times 8.0 \times
10^{-5}}} = \sqrt{2430} = 49.3\,\text{m}^{-1}, \quad m_{\text{St}} L =
1.23,
\]
so $\eta_{f,\text{St}} = \tanh(1.23)/1.23 = 0.843/1.23 = 0.685$. The fin
surface area is $A_f = 40 \times P L = 40 \times 0.162 \times 0.025 =
0.162\,\text{m}^2$, and with an exposed base of about $0.02\,\text{m}^2$
the total area is $A_t = 0.182\,\text{m}^2$, so $A_f/A_t = 0.890$. The
overall surface efficiencies are
\[
\eta_{o,\text{Al}} = 1 - 0.890(1 - 0.905) = 0.915, \quad
\eta_{o,\text{St}} = 1 - 0.890(1 - 0.685) = 0.720,
\]
and the array conductances are $h A_t \eta_o = 60 \times 0.182 \times
0.915 = 10.0\,\text{W}/\si{\kelvin}$ for aluminium against $7.9\,\text{W}/
\si{\kelvin}$ for steel. The aluminium array moves $27\,\%$ more heat from
the same geometry, because its higher conductivity keeps the fin tips
nearer the base temperature and so keeps more of the fin surface working.
The result is the quantitative form of the mastery box's claim that fins
pay off most when $kP/(hA_c)$ is large, and it is why a heat sink is cast
in aluminium or copper and almost never in steel.

\subsection{A selective surface and its equilibrium in sunlight}

A surface in sunlight reaches a temperature set by the contest between
its solar absorptance and its thermal emittance, and these two need not be
equal because they act in different wavelength bands. Sunlight arrives near
a $5800\,\si{\kelvin}$ spectrum peaking in the visible at about
$0.5\,\mu\text{m}$ by the Wien law $\lambda_{\max} = 2898/T$, while a
surface near $300\,\si{\kelvin}$ emits in the infrared peaking near
$10\,\mu\text{m}$. A surface whose absorptance differs between the two
bands is \engterm{spectrally selective}. We compute the equilibrium of an
isolated plate in space, illuminated by $1361\,\text{W}/\text{m}^2$ of
sunlight on one face and radiating from both faces, for two finishes: a
grey black paint ($\alpha_s = 0.95$, $\epsilon = 0.95$) and a selective
white paint ($\alpha_s = 0.20$, $\epsilon = 0.90$). The spectral picture
that lets the two coefficients differ is drawn in
Figure~\ref{fig:vol04ch05:selective-surface}: the solar band and the
surface's own emission band barely overlap, so a coating can hold a
different absorptance in each.

% Figure: the spectral basis of a selective surface. The 5800 K solar
% spectrum (short wavelength) and the 300 K surface spectrum (long
% wavelength), each normalised to unit peak, with a step-function
% spectral absorptance overlaid: high in the visible, low in the infrared.
% Curves are Planck shapes evaluated at the two temperatures, normalised.
\begin{figure}[htbp]
\centering
\begin{tikzpicture}
\begin{axis}[
  width=12cm, height=7cm,
  xmode=log,
  xlabel={wavelength $\lambda$ ($\mu$m)},
  ylabel={normalised spectral power / absorptance},
  xmin=0.2, xmax=40,
  ymin=0, ymax=1.15,
  axis lines=left,
  tick label style={font=\scriptsize},
  label style={font=\small},
  legend style={font=\scriptsize, at={(0.98,0.98)}, anchor=north east},
  legend cell align=left,
  legend entries={sun (5800 K),surface (300 K),selective absorptance},
  clip=false,
]
% Normalised Planck-shape proxies: use lognormal bumps peaking at Wien peaks
% solar peak ~ 0.5 um, surface peak ~ 9.7 um. Use exp(-((ln(l/lp))^2)/(2 s^2)).
\addplot[engamber, very thick, domain=0.2:40, samples=120]
  {exp(-((ln(x/0.5))^2)/(2*0.32))};
\addplot[engblue, very thick, domain=0.2:40, samples=120]
  {exp(-((ln(x/9.7))^2)/(2*0.30))};
% step-function selective absorptance: high below ~2.5 um, low above
\addplot[engred, thick, dashed] coordinates {(0.2,0.95) (2.5,0.95) (2.5,0.10) (40,0.10)};
\draw[enggray, dotted] (axis cs:2.5,0) -- (axis cs:2.5,1.15);
\node[font=\scriptsize, enggray, anchor=south, rotate=90] at (axis cs:2.5,0.30) {cut-off};
\end{axis}
\end{tikzpicture}
\caption[The spectral basis of a selective surface]{Why a surface can
absorb sunlight and yet emit weakly, or the reverse. The incoming solar
spectrum peaks in the visible near $0.5\,\mu$m, while a surface near room
temperature emits in the infrared near $10\,\mu$m; the two bands barely
overlap. A spectrally selective coating exploits the gap by carrying a
high absorptance below a cut-off wavelength and a low one above it (dashed
step). A solar absorber wants the profile drawn here, high in the solar
band and low in its own emission band, so it takes in sunlight and holds
the heat; a surface that must stay cool in sun wants the mirror image, low
solar absorptance and high infrared emittance. The single design number is
the ratio $\alpha_s/\epsilon$ of the band-averaged values, which sets the
equilibrium temperature in sunlight.}
\label{fig:vol04ch05:selective-surface}
\end{figure}


The balance equates absorbed solar on the one lit face to emitted infrared
from both faces,
\[
\alpha_s G_s A = \epsilon \sigma T^4 (2A), \quad\Rightarrow\quad
T = \left(\frac{\alpha_s G_s}{2 \epsilon \sigma}\right)^{1/4}.
\]
For the grey black paint, $\alpha_s/\epsilon = 1.0$, so
\[
T = \left(\frac{0.95 \times 1361}{2 \times 0.95 \times 5.67 \times
10^{-8}}\right)^{1/4} = \left(\frac{1293}{1.077 \times 10^{-7}}\right)^{1/4}
= (1.20 \times 10^{10})^{1/4} = 331\,\si{\kelvin},
\]
about $58\,\degC$. For the selective white paint, $\alpha_s/\epsilon =
0.20/0.90 = 0.222$, so the bracket scales by that ratio to the
quarter-power, $T = 331 \times 0.222^{1/4} = 331 \times 0.687 = 227\,
\si{\kelvin}$, about $-46\,\degC$. The same sunlight drives the two
finishes more than a hundred kelvin apart, set entirely by the ratio
$\alpha_s/\epsilon$. The design rule reads straight off the balance: a
surface that must stay cool in sun (a radiator, a tank, a rover deck)
wants a low $\alpha_s/\epsilon$, a finish that reflects the visible solar
band while still emitting strongly in the infrared, while a surface that
must collect heat wants the opposite. This ratio is the single number a
spacecraft thermal engineer reads off a coatings table, and it explains
why the radiator of the spacecraft case wore a white finish with
$\alpha_s/\epsilon$ well below one.

\subsection{Sizing a heat exchanger by the log-mean temperature
difference}

The combined-resistance machinery sizes a wall; a heat exchanger sizes the
area to transfer a duty between two flowing streams whose temperatures
change along the device, and the right driving difference is not the
inlet difference but the \engterm{log-mean temperature difference}. We
size a counterflow exchanger to cool $0.5\,\text{kg/s}$ of oil from
$120\,\degC$ to $60\,\degC$ using water entering at $20\,\degC$ at
$0.6\,\text{kg/s}$, with an overall coefficient $U = 350\,\text{W}/
(\text{m}^2\cdot\si{\kelvin})$. Take oil $c_{p,\text{oil}} = 2100\,\text{J}
/(\text{kg}\cdot\si{\kelvin})$ and water $c_{p,\text{w}} = 4180\,\text{J}/
(\text{kg}\cdot\si{\kelvin})$.

The duty is the oil's heat release,
\[
\dot{Q} = \dot{m}_{\text{oil}} c_{p,\text{oil}} \Delta T_{\text{oil}} =
0.5 \times 2100 \times (120 - 60) = 6.3 \times 10^4\,\text{W}.
\]
The water carries that heat away, so its outlet temperature is
\[
T_{\text{w,out}} = T_{\text{w,in}} + \frac{\dot{Q}}{\dot{m}_{\text{w}}
c_{p,\text{w}}} = 20 + \frac{6.3 \times 10^4}{0.6 \times 4180} = 20 +
25.1 = 45.1\,\degC.
\]
In counterflow the hot oil inlet faces the warmed water outlet and the
cooled oil outlet faces the cold water inlet, so the two end differences
are
\[
\Delta T_a = 120 - 45.1 = 74.9\,\si{\kelvin}, \quad \Delta T_b = 60 - 20 =
40\,\si{\kelvin}.
\]
The log-mean difference is
\[
\Delta T_{\text{lm}} = \frac{\Delta T_a - \Delta T_b}{\ln(\Delta T_a/
\Delta T_b)} = \frac{74.9 - 40}{\ln(74.9/40)} = \frac{34.9}{0.627} =
55.7\,\si{\kelvin},
\]
which sits between the two end differences and below their arithmetic mean
of $57.5$, the bias the logarithm always gives. The required area is then
\[
A = \frac{\dot{Q}}{U\,\Delta T_{\text{lm}}} = \frac{6.3 \times 10^4}{350
\times 55.7} = 3.23\,\text{m}^2.
\]
The log-mean form is the right average because the local driving
difference varies exponentially along a counterflow exchanger, not
linearly, so using the inlet difference of $100\,\si{\kelvin}$ would
undersize the area by nearly a factor of two. The same $UA = \dot{Q}/
\Delta T_{\text{lm}} = 1131\,\text{W}/\si{\kelvin}$ is the conductance the
chapter's resistance network would assign to the wall separating the two
streams, the heat exchanger being the combined-resistance idea applied to
two moving fluids.

\subsection{Centreline temperature of a current-carrying conductor}

Conduction with internal generation is the Poisson problem of section 5.2, and
its cleanest instance is a long cylindrical conductor heated uniformly by the
current it carries. We take a solid aluminium busbar of radius $R =
15\,\si{\milli\meter}$ carrying enough current to dissipate a volumetric heat
rate $\dot{q}_v = 4.0 \times 10^{5}\,\text{W}/\text{m}^3$, cooled at its surface
to $T_s = 70\,\degC$, with aluminium conductivity $k = 237\,\text{W}/(\text{m}
\cdot\si{\kelvin})$. We want the temperature at the centre and a check that the
internal rise is small enough that the bar will not soften.

At steady state with axial symmetry the heat equation reduces to the radial
Poisson form $\dfrac{1}{r}\dfrac{d}{dr}\!\left(r\dfrac{dT}{dr}\right) =
-\dfrac{\dot{q}_v}{k}$. Integrating once gives $r\,dT/dr = -\dot{q}_v r^2/(2k) +
C_1$, and finiteness at the centre forces $C_1 = 0$. Integrating again,
\[
T(r) = -\frac{\dot{q}_v r^2}{4k} + C_2,
\]
and applying $T(R) = T_s$ fixes $C_2$, leaving the parabolic profile
\[
T(r) = T_s + \frac{\dot{q}_v}{4k}\left(R^2 - r^2\right).
\]
The centre is hottest, with the centre-to-surface rise
\[
T_0 - T_s = \frac{\dot{q}_v R^2}{4k} = \frac{4.0\times10^{5}\times(0.015)^2}{4
\times 237} = \frac{90}{948} = 0.095\,\si{\kelvin}.
\]
The interior of a good conductor is nearly isothermal even while it generates
heat, because the same high conductivity that makes it carry current well makes
it conduct the generated heat to the surface against a tiny gradient. The
surface temperature, set by the convection the bar sees, is what an engineer
must control; the internal rise of a tenth of a kelvin is a rounding term. The
result inverts for a poor conductor: the same generation in a nuclear fuel
pellet ($k \approx 3\,\text{W}/(\text{m}\cdot\si{\kelvin})$) of the same size
would lift the centre by nearly $8\,\si{\kelvin}$ per unit of the busbar's
generation, and at reactor power densities the centreline-to-surface rise across
a fuel pellet runs into the hundreds of kelvin, which is why fuel-pellet
centreline melting is a design limit while busbar centreline temperature never
is.

\subsection{A hot plate losing heat by convection and radiation together}

When free convection and radiation both carry heat from a surface and the
surface temperature is the unknown, the energy balance is implicit because
radiation is quartic in temperature, and it is solved by iteration. We take a
horizontal electronics enclosure lid, $0.3\,\text{m}\times 0.3\,\text{m}$,
dissipating $\dot{Q} = 40\,\text{W}$ from its top face into a still room at
$T_\infty = 25\,\degC = 298\,\si{\kelvin}$, with surroundings also at
$298\,\si{\kelvin}$ and surface emissivity $\epsilon = 0.9$. We want the steady
lid temperature.

The balance sets the dissipated power equal to the sum of the convective and
radiative losses from the $A = 0.09\,\text{m}^2$ face,
\[
\dot{Q} = h_{\text{conv}}(T_s)\,A\,(T_s - T_\infty) + \epsilon\sigma A
\,(T_s^4 - T_\infty^4),
\]
where the free-convection coefficient itself depends on $T_s$ through the
horizontal-plate correlation $\text{Nu}_L = 0.54\,\text{Ra}_L^{1/4}$, so
$h_{\text{conv}} \propto (T_s - T_\infty)^{1/4}$ and the convective term scales
as $(T_s - T_\infty)^{5/4}$. We iterate. Guess $T_s = 330\,\si{\kelvin}$, a
$32\,\si{\kelvin}$ rise. For air at the film temperature take the lumped
free-convection estimate $h_{\text{conv}} \approx 1.3\,(\Delta T/L_c)^{1/4}$
with the plate characteristic length $L_c = A/P = 0.09/1.2 = 0.075\,\text{m}$,
giving $h_{\text{conv}} \approx 1.3\,(32/0.075)^{1/4} = 1.3 \times 4.55 =
5.9\,\text{W}/(\text{m}^2\cdot\si{\kelvin})$. The convective loss is $5.9 \times
0.09 \times 32 = 17.0\,\text{W}$. The radiative loss is $0.9 \times 5.67
\times10^{-8} \times 0.09 \times (330^4 - 298^4) = 0.9 \times 5.67\times10^{-8}
\times 0.09 \times (1.186\times10^{10} - 7.886\times10^{9}) = 0.9 \times 5.67
\times10^{-8}\times 0.09 \times 3.97\times10^{9} = 18.2\,\text{W}$. The two sum
to $35.2\,\text{W}$, short of the $40\,\text{W}$ supplied, so the lid must run
hotter. Raise the guess to $T_s = 335\,\si{\kelvin}$, a $37\,\si{\kelvin}$ rise:
$h_{\text{conv}} \approx 1.3\,(37/0.075)^{1/4} = 6.1$, convection $6.1 \times
0.09 \times 37 = 20.3\,\text{W}$, and radiation $0.9 \times 5.67\times10^{-8}
\times 0.09\times(335^4 - 298^4) = 21.5\,\text{W}$, summing to
$41.8\,\text{W}$, just over. Interpolating, the lid settles near $T_s =
334\,\si{\kelvin} = 61\,\degC$. The instructive feature is that radiation and
free convection carry almost equal shares of the load at this modest
temperature, so a calculation that kept only convection would have overstated
the lid temperature by tens of kelvin. The linearised radiation coefficient
$h_{\text{rad}} = 4\epsilon\sigma T_{\text{avg}}^3 \approx 4\times 0.9 \times
5.67\times10^{-8}\times 316^3 = 6.4\,\text{W}/(\text{m}^2\cdot\si{\kelvin})$
confirms the two paths are comparable and lets the next design iteration run as
a single linear resistance network.

\subsection{Charging time of a thermal-storage tank}

A lumped body driven by a fixed heat input rather than a fixed ambient warms
toward a moving target, and the same first-order balance gives the charging
time. We take a $300\,\text{L}$ water thermal-storage tank with a
$3\,\text{kW}$ immersion heater, starting at $15\,\degC$, losing heat to a
$20\,\degC$ plant room through a jacket of overall conductance $UA =
3.5\,\text{W}/\si{\kelvin}$, and ask how long it takes to reach $60\,\degC$ and
what steady temperature it would eventually hold if left on.

The energy balance on the well-mixed tank balances the heater input against the
storage and the jacket loss,
\[
m c_p \frac{dT}{dt} = \dot{Q}_{\text{heater}} - UA\,(T - T_{\text{room}}),
\]
with $m = 300\,\text{kg}$ and $c_p = 4180\,\text{J}/(\text{kg}\cdot\si{
\kelvin})$. This is a first-order linear equation whose steady state is the
temperature at which the loss equals the input,
\[
T_\infty^* = T_{\text{room}} + \frac{\dot{Q}_{\text{heater}}}{UA} = 20 +
\frac{3000}{3.5} = 20 + 857 = 877\,\degC,
\]
a temperature the tank can never physically reach (the water boils first), which
tells us the jacket loss is negligible over the working range and the tank warms
almost linearly. The time constant is $\tau = m c_p/(UA) = 300 \times
4180/3.5 = 3.58\times10^{5}\,\text{s}$, about $100$ hours, far longer than the
charge will take. The solution is
\[
T(t) = T_\infty^* - (T_\infty^* - T_0)\,e^{-t/\tau},
\]
and solving for the time to reach $60\,\degC$,
\[
t = \tau\ln\!\frac{T_\infty^* - T_0}{T_\infty^* - T} = 3.58\times10^{5}\,
\ln\!\frac{877 - 15}{877 - 60} = 3.58\times10^{5}\times\ln(1.055) =
3.58\times10^{5}\times 0.0535 = 1.91\times10^{4}\,\text{s},
\]
about $5.3$ hours. The simple energy check confirms it: heating $300\,\text{kg}$
of water by $45\,\si{\kelvin}$ takes $m c_p\,\Delta T = 300\times4180\times45 =
5.64\times10^{7}\,\text{J}$, which at the net input of roughly $3000 -
\overline{UA\,\Delta T} \approx 2900\,\text{W}$ takes $5.64\times10^{7}/2900 =
1.95\times10^{4}\,\text{s}$, matching the exponential solution because the loss
term is small. The lesson the numbers carry is that for a well-insulated store
the charge is set by thermal mass over heater power and the jacket loss only
trims it; the same balance with a larger $UA$ (a poorly lagged tank) would pull
the unreachable steady temperature down toward the working range and make the
loss term bite, lengthening the charge and capping the attainable temperature.

\subsection{Radiant exchange between two perpendicular surfaces}

The two-surface formulas of section 5.4 need a view factor, and for surfaces
that are neither parallel nor concentric the factor comes from a catalogue or a
crossed-strings construction. We take an oven with a heated floor strip and an
adjacent perpendicular side wall, modelled as two long perpendicular plates
sharing a common edge, each $0.4\,\text{m}$ wide per metre of length, with the
floor at $T_1 = 500\,\si{\kelvin}$ and emissivity $\epsilon_1 = 0.85$ and the
wall at $T_2 = 350\,\si{\kelvin}$ and emissivity $\epsilon_2 = 0.7$, and we find
the net radiant exchange per metre of length.

For two long plates of equal width $w$ meeting at a right angle along a common
edge, the crossed-strings method gives the floor-to-wall view factor as the sum
of the crossed string lengths minus the uncrossed, divided by twice the source
width. The crossed strings are the two diagonals of length $w\sqrt{2}$, the
uncrossed strings are the shared edge (length $0$) and the far gap (the straight
line between the two outer edges, length $w\sqrt{2}$ as well by the right-angle
geometry); the standard result for two equal perpendicular strips sharing an
edge is $F_{12} = 1 - \sin(45^\circ) = 1 - 0.707 = 0.293$. With $A_1 = A_2 =
0.4\,\text{m}^2$ per metre, reciprocity gives $F_{21} = F_{12} = 0.293$ since the
areas are equal. Feeding these into the two-surface exchange formula with the
surface resistances,
\[
\dot{Q}_{1\to2} = \frac{\sigma(T_1^4 - T_2^4)}{\dfrac{1-\epsilon_1}{\epsilon_1
A_1} + \dfrac{1}{A_1 F_{12}} + \dfrac{1-\epsilon_2}{\epsilon_2 A_2}}.
\]
The three resistances per metre are $\dfrac{1-0.85}{0.85\times0.4} = 0.441$,
$\dfrac{1}{0.4\times0.293} = 8.53$, and $\dfrac{1-0.7}{0.7\times0.4} = 1.071$,
summing to $10.04\,\text{m}^{-2}$. The driving term is $\sigma(T_1^4 - T_2^4) =
5.67\times10^{-8}(500^4 - 350^4) = 5.67\times10^{-8}(6.25\times10^{10} -
1.501\times10^{10}) = 5.67\times10^{-8}\times 4.749\times10^{10} =
2693\,\text{W}/\text{m}^2$, so
\[
\dot{Q}_{1\to2} = \frac{2693}{10.04} = 268\,\text{W} \text{ per metre of
length}.
\]
The space resistance $8.53$ dominates the chain, because the perpendicular
geometry lets most of the floor's radiation miss the wall and strike the rest of
the enclosure; only $29\,\%$ of what leaves the floor reaches this wall
directly. The lesson is that geometry, through the view factor, can be the
controlling resistance in a radiation problem just as a thick insulation layer
is in a conduction problem, and that improving the emissivities buys little when
the view factor is the bottleneck.

\subsection{Vignette: sizing a car radiator by effectiveness-NTU}

A car radiator rejects engine waste heat to the air stream, and sizing it is the
heat-exchanger problem in its second canonical form. The log-mean method of the
preceding example sizes the area when the four terminal temperatures are known;
the \engterm{effectiveness-NTU} method sizes or rates the device when an outlet
temperature is the unknown, which is the usual case for a radiator whose air
outlet nobody measures. The two methods are the same physics in different
algebra, and a working engineer carries both.

A small engine rejects $\dot{Q} = 35\,\text{kW}$ to its coolant, which the water
pump circulates at $\dot{m}_c = 0.8\,\text{kg/s}$ entering the radiator at
$95\,\degC$, while the fan and vehicle motion drive air across the core at
$\dot{m}_a = 0.7\,\text{kg/s}$ entering at $35\,\degC$. Take coolant
$c_{p,c} = 3900\,\text{J}/(\text{kg}\cdot\si{\kelvin})$ (a $50/50$ glycol mix)
and air $c_{p,a} = 1007\,\text{J}/(\text{kg}\cdot\si{\kelvin})$. We size the
core conductance $UA$ that rejects the duty.

The two stream capacity rates are
\[
C_c = \dot{m}_c c_{p,c} = 0.8\times 3900 = 3120\,\text{W}/\si{\kelvin}, \quad
C_a = \dot{m}_a c_{p,a} = 0.7\times 1007 = 705\,\text{W}/\si{\kelvin},
\]
so the air is the minimum-capacity stream, $C_{\min} = 705$ and $C_{\max} =
3120$, with capacity ratio $C_r = C_{\min}/C_{\max} = 0.226$. The thermodynamic
maximum heat transfer is the minimum-capacity stream taken across the full inlet
difference,
\[
\dot{Q}_{\max} = C_{\min}(T_{c,\text{in}} - T_{a,\text{in}}) = 705\times(95 -
35) = 4.23\times10^{4}\,\text{W} = 42.3\,\text{kW}.
\]
The required effectiveness is the actual duty over this maximum,
\[
\varepsilon = \frac{\dot{Q}}{\dot{Q}_{\max}} = \frac{35.0}{42.3} = 0.828.
\]
For a crossflow core with both fluids unmixed, the effectiveness-NTU relation is
well approximated by $\varepsilon = 1 - \exp\!\big[(\text{NTU}^{0.22}/C_r)
(\exp(-C_r\,\text{NTU}^{0.78}) - 1)\big]$, but the near-unity capacity-ratio
behaviour and the value of $\varepsilon$ here are read more simply from the
counterflow form, which for $C_r = 0.226$ and $\varepsilon = 0.828$ inverts
through
\[
\text{NTU} = \frac{1}{C_r - 1}\ln\!\frac{\varepsilon - 1}{\varepsilon C_r - 1}
= \frac{1}{-0.774}\ln\!\frac{-0.172}{-0.813} = \frac{1}{-0.774}\ln(0.2115) =
\frac{-1.553}{-0.774} = 2.01.
\]
The crossflow core needs a slightly larger NTU than the counterflow figure for
the same effectiveness, near $\text{NTU}\approx 2.3$, and the core conductance
follows from the definition $\text{NTU} = UA/C_{\min}$,
\[
UA = \text{NTU}\times C_{\min} = 2.3 \times 705 = 1.62\times10^{3}\,\text{W}/
\si{\kelvin}.
\]
Figure~\ref{fig:vol04ch05:effectiveness-ntu} places this on the
effectiveness-NTU curves and shows the design lives on the flattening part of
the curve, where each further unit of area buys little effectiveness, so a
radiator sized much larger would reject barely more heat and a radiator sized
much smaller would overheat the engine quickly. The air outlet temperature
closes the picture: the air carries the full duty, so $T_{a,\text{out}} =
35 + 35000/705 = 35 + 49.6 = 84.6\,\degC$, and the coolant leaves at
$T_{c,\text{out}} = 95 - 35000/3120 = 95 - 11.2 = 83.8\,\degC$, the two outlets
nearly equal because the device runs close to its thermal limit. The
effectiveness-NTU method gave the size directly from the inlet conditions and
the duty, without the iterative guessing the log-mean method would have demanded
when an outlet is unknown, which is why radiator and intercooler engineers reach
for it first.

% Figure: heat-exchanger effectiveness versus number of transfer units for a
% counterflow exchanger at several capacity-rate ratios C_r = C_min/C_max. The
% curves rise steeply at low NTU and saturate, showing the diminishing return of
% adding area, and the marked point is the car-radiator vignette's operating
% point. Computed from the closed-form counterflow relation.
\begin{figure}[htbp]
\centering
\begin{tikzpicture}
\begin{axis}[
  width=12cm, height=7.2cm,
  xlabel={number of transfer units $\text{NTU}=UA/C_{\min}$},
  ylabel={effectiveness $\varepsilon$},
  xmin=0, xmax=5, ymin=0, ymax=1.02,
  axis lines=left,
  legend pos=south east,
  legend cell align=left,
  tick label style={font=\scriptsize},
  label style={font=\small},
  legend style={font=\scriptsize},
  domain=0.01:5, samples=80,
  clip=false,
]
% C_r = 0: epsilon = 1 - exp(-NTU)
\addplot[enggray, very thick] {1 - exp(-x)};
\addlegendentry{$C_r=0$ (phase change)}
% C_r = 0.5 counterflow: (1-exp(-NTU(1-Cr)))/(1-Cr exp(-NTU(1-Cr)))
\addplot[englime, very thick] {(1-exp(-x*0.5))/(1-0.5*exp(-x*0.5))};
\addlegendentry{$C_r=0.5$}
% C_r = 1 counterflow: NTU/(1+NTU)
\addplot[engblue, very thick] {x/(1+x)};
\addlegendentry{$C_r=1$}
% operating point of the radiator vignette, NTU ~ 1.4, Cr ~ 0.5
\addplot[engred, only marks, mark=*, mark options={fill=engred}]
  coordinates {(1.4,0.71)};
\node[engred, font=\scriptsize, anchor=west] at (axis cs:1.55,0.66)
  {radiator design point};
\end{axis}
\end{tikzpicture}
\caption[Effectiveness versus NTU for a counterflow exchanger]{Heat-exchanger
effectiveness $\varepsilon$ (the actual heat transfer divided by the
thermodynamic maximum) against the number of transfer units
$\text{NTU}=UA/C_{\min}$, for several capacity-rate ratios
$C_r=C_{\min}/C_{\max}$, in counterflow. The curves climb steeply at low NTU and
flatten as effectiveness approaches its ceiling, so the first units of area buy
most of the available duty and each further increment buys less. The
limiting curve $C_r=0$ is the phase-change case (a condenser or boiler), where
one stream holds its temperature; the $C_r=1$ curve is the balanced case, the
hardest to drive toward unity. The marked point is the car-radiator vignette,
where a moderate NTU near $1.4$ at $C_r\approx 0.5$ yields an effectiveness near
$0.71$, enough to reject the engine heat without an unreasonably large core.}
\label{fig:vol04ch05:effectiveness-ntu}
\end{figure}


\subsection{Heat loss from a buried pipe by the conduction shape factor}

Section 5.2 named the shape-factor catalogue as the tool for conduction
between two isothermal surfaces of fixed shape and then set it aside; here
we put it to work on the most common instance, a pipe or cable buried in
soil. The plane and cylindrical resistance formulas fail for a buried line
because the heat does not flow through a single-area path; it spreads
three-dimensionally into the ground and crowds toward the surface, as the
isotherms of Figure~\ref{fig:vol04ch05:shape-factor-buried} show. The
shape factor $S$ carries all of that geometry, reducing the two-surface
problem to $\dot{Q} = S k\,(T_p - T_s)$ with a resistance $R = 1/(Sk)$
that adds in series like any other.

% Figure: the conduction shape factor for a cylinder buried in a
% semi-infinite medium. A schematic section of ground with an isothermal
% buried pipe at depth z, the isotherms crowding near the surface, and
% the shape-factor formula that turns the two-surface geometry into a
% single conductance. Schematic; isotherms drawn by hand, not computed.
\begin{figure}[htbp]
\centering
\begin{tikzpicture}[>=Latex]
% ground surface
\draw[enggray, very thick] (-4,0) -- (4,0);
\node[font=\scriptsize, enggray, anchor=south west] at (-4,0.05)
  {ground surface, $T_s$};
% shade the soil
\begin{scope}
\clip (-4,0) rectangle (4,-3.4);
\foreach \r in {0.9,1.4,1.9,2.5,3.2} {
  \draw[engblue!45, thin] (0,-1.6) circle (\r);
}
\end{scope}
% buried pipe
\filldraw[engred!20, draw=engred, very thick] (0,-1.6) circle (0.45);
\node[font=\scriptsize, engred] at (0,-1.6) {$T_p$};
% depth dimension
\draw[<->, engamber, thick] (2.7,0) -- (2.7,-1.6);
\node[font=\scriptsize, engamber, anchor=west] at (2.75,-0.8) {$z$};
% diameter callout
\draw[<->, engred, thick] (-0.45,-2.35) -- (0.45,-2.35);
\node[font=\scriptsize, engred, anchor=north] at (0,-2.4) {$D$};
% isotherm label
\node[font=\scriptsize, engblue, anchor=west] at (1.1,-2.6) {isotherms};
% formula box
\node[font=\small, anchor=north west, align=left] at (-3.9,-2.6)
  {$\dot{Q} = S\,k\,(T_p - T_s)$\\[2pt]
   $S = \dfrac{2\pi L}{\cosh^{-1}(2z/D)}$};
\end{tikzpicture}
\caption[Conduction shape factor for a buried cylinder]{A pipe or cable
buried at depth $z$ in soil of conductivity $k$ conducts to the ground
surface along the curved paths sketched, the isotherms crowding toward the
surface where the area for spreading is least. The two isothermal surfaces
(pipe wall at $T_p$, ground surface at $T_s$) are joined by a single
conduction shape factor $S$ that carries all of the geometry, so the heat
rate is $\dot{Q} = S\,k\,(T_p - T_s)$ with $S = 2\pi L/\cosh^{-1}(2z/D)$
for a cylinder of length $L$ and diameter $D$. The shape factor is the
catalogue tool for conduction between two isothermal surfaces of fixed
shape, the conduction analogue of the view factor in radiation, and it
turns a two-dimensional spreading problem into the same resistance
$R = 1/(Sk)$ that the plane and cylindrical formulas provide for their
own geometries.}
\label{fig:vol04ch05:shape-factor-buried}
\end{figure}


We take a district-heating pipe of outer diameter $D = 0.20\,\text{m}$
carrying water that holds its outer wall at $T_p = 70\,\degC$, buried with
its axis $z = 0.80\,\text{m}$ below a ground surface at $T_s = 5\,\degC$,
in soil of conductivity $k = 1.5\,\text{W}/(\text{m}\cdot\si{\kelvin})$.
The shape factor for a cylinder of length $L$ buried at depth $z$ in a
semi-infinite medium is
\[
S = \frac{2\pi L}{\cosh^{-1}(2z/D)},
\]
valid for $z > D/2$, which holds here. The argument is $2z/D = 1.60/0.20 =
8.0$, and $\cosh^{-1}(8.0) = \ln(8.0 + \sqrt{8.0^2 - 1}) = \ln(8.0 +
7.94) = \ln(15.94) = 2.769$, so per metre of pipe $S/L = 2\pi/2.769 =
2.27\,\text{m}$. The heat loss per metre is
\[
\frac{\dot{Q}}{L} = \frac{S}{L}\,k\,(T_p - T_s) = 2.27 \times 1.5 \times
(70 - 5) = 221\,\text{W}/\text{m}.
\]
The depth enters only through the slowly varying inverse hyperbolic cosine,
so burying the pipe twice as deep, $z = 1.6\,\text{m}$, raises the argument
to $16$, $\cosh^{-1}(16) = 3.466$, and cuts the loss only to $2\pi/3.466
\times 1.5 \times 65 = 177\,\text{W}/\text{m}$, a $20\,\%$ reduction for
doubling the excavation. Depth is a weak lever on a bare buried line,
which is why district-heating pipes are lagged rather than merely buried
deep: a $50\,\si{\milli\meter}$ jacket of $k = 0.04$ insulation adds a
cylindrical resistance $\ln(0.30/0.20)/(2\pi\times0.04) = 1.61\,\si{
\kelvin}\cdot\text{m}/\text{W}$ in series with the soil path
$1/(Sk) = 1/(2.27\times1.5) = 0.294\,\si{\kelvin}\cdot\text{m}/\text{W}$,
dropping the loss to $65/(1.61 + 0.294) = 34\,\text{W}/\text{m}$, a
factor-of-six cut the depth could never buy. The shape factor is the
conduction analogue of the view factor: a single geometric number that
turns an awkward field problem into one more resistance in the network.

\subsection{A finless versus finned natural-convection sink}

The fin section argued that fins pay off most on a low-coefficient surface;
here we put numbers to a fanless enclosure that must shed its heat by
natural convection alone, and show both the gain a fin array buys and the
diminishing return that limits it. We take a sealed outdoor equipment box
whose one usable face is $0.20\,\text{m}\times0.20\,\text{m}$, dissipating
$\dot{Q} = 25\,\text{W}$ into still air at $T_a = 30\,\degC$, and ask the
face temperature bare, then with a modest fin array, taking a
natural-convection coefficient on a vertical surface of $h = 6\,\text{W}/
(\text{m}^2\cdot\si{\kelvin})$. Bare, the face area is $A = 0.04\,\text{m}^2$
and the rise is $\Delta T = \dot{Q}/(hA) = 25/(6\times0.04) = 104\,\si{
\kelvin}$, so the face reaches $134\,\degC$, far too hot for the
electronics inside. The convection resistance $1/(hA) = 4.17\,\si{\kelvin}/
\text{W}$ is the whole problem, and the only levers are area and
coefficient.

Add an aluminium fin array to the face: $16$ vertical fins, each $0.20\,
\text{m}$ deep, $0.030\,\text{m}$ tall, $2\,\si{\milli\meter}$ thick. Each
fin has perimeter $P = 2(0.20 + 0.002) = 0.404\,\text{m}$ and cross-section
$A_c = 0.20\times0.002 = 4.0\times10^{-4}\,\text{m}^2$, so with $k = 237$
the fin parameter is $m = \sqrt{hP/(kA_c)} = \sqrt{6\times0.404/(237\times
4.0\times10^{-4})} = \sqrt{25.6} = 5.06\,\text{m}^{-1}$ and $mL = 5.06\times
0.030 = 0.152$, a very stubby fin of efficiency $\eta_f = \tanh(0.152)/
0.152 = 0.992$, so the fins run almost isothermal with the base. The added
fin surface is $A_f = 16\times P L = 16\times0.404\times0.030 = 0.194\,
\text{m}^2$, and with the exposed base between fins ($\approx 0.02\,\text{m}
^2$) the total effective area is $A_t = 0.992\times0.194 + 0.02 = 0.212\,
\text{m}^2$, more than five times the bare face. The new resistance is
$1/(hA_t) = 1/(6\times0.212) = 0.786\,\si{\kelvin}/\text{W}$, so the rise
falls to $\Delta T = 25\times0.786 = 19.7\,\si{\kelvin}$ and the base sits
at $50\,\degC$, comfortably inside limits. The array cut the resistance by
the area ratio, over five to one, because the fins are short and highly
conductive so almost every square centimetre works. The diminishing return
appears if we push the fins taller: doubling the height to $0.060\,\text{m}$
raises $mL$ to $0.30$, still efficient, and doubles the area, but a natural-
convection array packed too tightly chokes its own buoyant flow between the
fins and the effective $h$ falls, so beyond a spacing near the boundary-
layer thickness (a few millimetres in air) more fins stop helping. The
fanless sink is sized by area against a fixed low coefficient, the exact
regime the fin effectiveness $\sqrt{kP/(hA_c)}$ was built to reward.

\subsection{A double-pane window with a low-emissivity coating}

The window exercise showed radiation carrying half a still-air gap's
conductance; here we work the design response, a low-emissivity coating, and
show why it is the single largest lever in modern glazing. We take a
double-pane unit, two $4\,\si{\milli\meter}$ glass panes across a
$16\,\si{\milli\meter}$ argon-filled gap, indoor $21\,\degC$, outdoor
$-5\,\degC$, and compare an ordinary uncoated unit against one whose inner
pane carries a low-emissivity coating of $\epsilon = 0.10$ facing the gap.
The gap conductance is the sum of the gas path and the radiation path in
parallel. Argon in a $16\,\si{\milli\meter}$ gap conducts at
$k_{\text{gas}}/L = 0.017/0.016 = 1.06\,\text{W}/(\text{m}^2\cdot\si{
\kelvin})$ (argon's lower conductivity than air is itself part of the
design). For the uncoated unit both surfaces have $\epsilon = 0.84$, giving
an effective emissivity $1/(1/0.84 + 1/0.84 - 1) = 0.72$ and a linearised
radiation coefficient at the gap mean temperature near $278\,\si{\kelvin}$,
\[
h_{\text{rad}} = \epsilon_{\text{eff}}\,4\sigma T_m^3 = 0.72\times4\times
5.67\times10^{-8}\times278^3 = 3.51\,\text{W}/(\text{m}^2\cdot\si{\kelvin}),
\]
so the gap conductance is $1.06 + 3.51 = 4.57$ and its resistance is
$0.219\,\si{\kelvin}\cdot\text{m}^2/\text{W}$. The coated unit drops the
effective emissivity to $1/(1/0.10 + 1/0.84 - 1) = 0.0989$, cutting the
radiation coefficient to $0.0989/0.72\times3.51 = 0.482\,\text{W}/(\text{m}
^2\cdot\si{\kelvin})$, so the gap conductance falls to $1.06 + 0.48 = 1.54$
and its resistance rises to $0.649\,\si{\kelvin}\cdot\text{m}^2/\text{W}$,
tripling the gap's insulating value. Adding the two glass panes (negligible)
and the indoor and outdoor films ($1/8 = 0.125$ and $1/25 = 0.040$), the
uncoated unit totals $R = 0.040 + 0.219 + 0.125 = 0.384$, $U = 2.60\,
\text{W}/(\text{m}^2\cdot\si{\kelvin})$, while the coated unit totals $R =
0.040 + 0.649 + 0.125 = 0.814$, $U = 1.23\,\text{W}/(\text{m}^2\cdot\si{
\kelvin})$. The coating alone more than halves the U-value, because it
attacks the radiation path that carried three-quarters of the gap
conductance, and it does so with an invisibly thin metal-oxide film rather
than more glass or a wider gap. This is why the low-emissivity coating and
the argon fill, not the pane count, define a modern insulating unit, and
why the exercise's warning against ignoring gap radiation is the very
opening the coating exploits.

\subsection{A pipe outside the lumped range by the one-term solution}

The lumped-capacitance model of section 5.2 fails once the Biot number
climbs past about $0.1$, and the one-term series solution of the mastery
box is the working tool for the range beyond it. We take a solid nylon rod
of diameter $D = 60\,\si{\milli\meter}$, initially at $T_i = 20\,\degC$
throughout, suddenly plunged into a $90\,\degC$ water bath with a
convection coefficient $h = 600\,\text{W}/(\text{m}^2\cdot\si{\kelvin})$,
and ask how long the centreline takes to reach $60\,\degC$. Take nylon
$k = 0.25\,\text{W}/(\text{m}\cdot\si{\kelvin})$, $\rho = 1140\,\text{kg/m}
^3$, $c_p = 1670\,\text{J}/(\text{kg}\cdot\si{\kelvin})$, so the diffusivity
is $\alpha = k/(\rho c_p) = 0.25/(1140\times1670) = 1.31\times10^{-7}\,
\text{m}^2/\text{s}$.

The Biot number for the long cylinder uses the radius as the length scale,
$\text{Bi} = h R/k = 600\times0.030/0.25 = 72$, far above the
lumped-capacitance ceiling, so the body carries a steep internal gradient
and the lumped model is disallowed. We use the one-term solution for the
centre of an infinite cylinder,
\[
\frac{T_0 - T_\infty}{T_i - T_\infty} = C_1\,e^{-\lambda_1^2\,\text{Fo}},
\]
where for a cylinder the eigenvalue solves $\lambda_1 J_1(\lambda_1)/
J_0(\lambda_1) = \text{Bi}$ and the coefficient is $C_1 = (2/\lambda_1)
J_1(\lambda_1)/[J_0^2(\lambda_1) + J_1^2(\lambda_1)]$. For a large Biot
number the eigenvalue saturates at the first zero of $J_0$, so
$\lambda_1 \to 2.405$; at $\text{Bi} = 72$ the tabulated values are
$\lambda_1 = 2.381$ and $C_1 = 1.599$. The target centre ratio is
$(60 - 90)/(20 - 90) = (-30)/(-70) = 0.429$, so
\[
e^{-\lambda_1^2\,\text{Fo}} = \frac{0.429}{1.599} = 0.268, \quad
\lambda_1^2\,\text{Fo} = -\ln(0.268) = 1.317,
\]
and with $\lambda_1^2 = 5.67$, $\text{Fo} = 1.317/5.67 = 0.232$, safely
above the $\text{Fo} \gtrsim 0.2$ threshold where the one-term truncation
is accurate (Figure~\ref{fig:vol04ch05:heisler-oneterm}). The physical
time follows from the Fourier number,
\[
t = \frac{\text{Fo}\,R^2}{\alpha} = \frac{0.232\times(0.030)^2}{1.31\times
10^{-7}} = \frac{2.09\times10^{-4}}{1.31\times10^{-7}} = 1594\,\text{s},
\]
about twenty-seven minutes. A lumped estimate would have used $L_c = R/2 =
0.015\,\text{m}$ and a time constant $\tau = \rho c_p L_c/h = 1140\times
1670\times0.015/600 = 47.6\,\text{s}$, predicting the centre at
$60\,\degC$ in $\tau\ln(70/30) = 47.6\times0.847 = 40\,\text{s}$, forty
times too fast, because it ignores the internal resistance that dominates
at $\text{Bi} = 72$. The one-term solution is not a refinement of the
lumped answer; the two live in different regimes, and reading the Biot
number first is what tells the engineer which tool the problem is in.

% Figure: the one-term transient chart for the symmetric plane wall,
% dimensionless centre excess temperature against Fourier number on a
% semilog plot for several Biot numbers. Computed from the leading term
% theta0 = C1 exp(-lambda1^2 Fo) with tabulated eigenvalues, so each
% curve is a straight line on the semilog axes past Fo ~ 0.2. The
% coordinates are precomputed constants, not evaluated in TeX.
\begin{figure}[htbp]
\centering
\begin{tikzpicture}
\begin{semilogyaxis}[
  width=12cm, height=8cm,
  xlabel={Fourier number $\text{Fo} = \alpha t/L^2$},
  ylabel={centre excess $\theta_0/\theta_i = (T_0 - T_\infty)/(T_i - T_\infty)$},
  xmin=0, xmax=3, ymin=0.02, ymax=1.2,
  axis lines=left,
  tick label style={font=\scriptsize},
  label style={font=\small},
  legend style={font=\scriptsize, at={(0.98,0.98)}, anchor=north east},
  legend cell align=left,
  legend entries={$\mathrm{Bi}=0.1$,$\mathrm{Bi}=1$,$\mathrm{Bi}=10$,$\mathrm{Bi}$ large},
  clip=false,
]
% Bi = 0.1: lambda1 = 0.3111, C1 = 1.0161 ; slope = -0.0968
\addplot[engblue, very thick, domain=0.2:3, samples=30]
  {1.0161*exp(-0.09678*x)};
% Bi = 1.0: lambda1 = 0.8603, C1 = 1.1191 ; lambda^2 = 0.7401
\addplot[engred, very thick, domain=0.2:3, samples=30]
  {1.1191*exp(-0.7401*x)};
% Bi = 10: lambda1 = 1.4289, C1 = 1.2620 ; lambda^2 = 2.0418
\addplot[engamber, very thick, domain=0.2:3, samples=30]
  {1.2620*exp(-2.0418*x)};
% Bi -> infinity: lambda1 = pi/2 = 1.5708, C1 = 1.2732 ; lambda^2 = 2.4674
\addplot[enggray, very thick, dashed, domain=0.2:3, samples=30]
  {1.2732*exp(-2.4674*x)};
% mark the one-term validity boundary
\draw[enggray, dotted] (axis cs:0.2,0.02) -- (axis cs:0.2,1.2);
\node[font=\scriptsize, enggray, anchor=south, rotate=90]
  at (axis cs:0.2,0.09) {one-term valid, $\mathrm{Fo}\gtrsim 0.2$};
\end{semilogyaxis}
\end{tikzpicture}
\caption[One-term transient chart for the plane wall]{The centre
temperature of a symmetric plane wall suddenly exposed to a convecting
fluid, plotted as the dimensionless centre excess against Fourier number
for several Biot numbers. Past a Fourier number of about $0.2$ the higher
diffusion modes have died and the response is the single leading term
$\theta_0/\theta_i = C_1 \exp(-\lambda_1^2\,\text{Fo})$, a straight line on
these semilogarithmic axes whose slope is set by the smallest eigenvalue
$\lambda_1$ of $\lambda_1 \tan\lambda_1 = \text{Bi}$. A low Biot number
gives a shallow slope and a slow, nearly uniform cooling, the
lumped-capacitance limit; a high Biot number gives a steep slope as the
surface convection strips heat faster than conduction can feed it from the
centre, and the eigenvalue saturates at $\pi/2$ for the perfectly
convected face. The chart is the working replacement for the full
eigenfunction series when the Biot number puts a body outside the
lumped-capacitance range.}
\label{fig:vol04ch05:heisler-oneterm}
\end{figure}


\subsection{Peak temperature in a nuclear fuel pin}

The busbar example showed a good conductor generating heat with a
negligible internal rise; the same Poisson problem in a poor conductor,
across the gap and clad that surround a reactor fuel pellet, sets the
design limit of the whole reactor. We take a cylindrical uranium-dioxide
fuel pellet of radius $R_f = 4.1\,\si{\milli\meter}$ generating a
volumetric heat rate $\dot{q}_v = 3.0\times10^8\,\text{W}/\text{m}^3$,
surrounded by a helium-filled gap of conductance $h_{\text{gap}} =
6000\,\text{W}/(\text{m}^2\cdot\si{\kelvin})$, a zirconium-alloy cladding
of inner radius $R_{ci} = 4.2\,\si{\milli\meter}$ and outer radius
$R_{co} = 4.75\,\si{\milli\meter}$ ($k_c = 17\,\text{W}/(\text{m}\cdot
\si{\kelvin})$), cooled on the outside by water at $T_w = 300\,\degC$ with
$h_w = 3.0\times10^4\,\text{W}/(\text{m}^2\cdot\si{\kelvin})$. Take the
fuel conductivity $k_f = 3.0\,\text{W}/(\text{m}\cdot\si{\kelvin})$, low
and falling with burn-up. We want the centreline temperature and its
margin against the $\sim 2800\,\degC$ melting point of the oxide.

The linear heat rate carried by the pin is the generation times the pellet
area, $q' = \dot{q}_v \pi R_f^2 = 3.0\times10^8\times\pi\times(0.0041)^2 =
1.58\times10^4\,\text{W}/\text{m}$, about $16\,\text{kW}$ per metre, a
representative value for a power reactor. We march outward from the coolant
through the series of resistances, each carrying this same linear rate.
The convective film on the clad drops
\[
\Delta T_w = \frac{q'}{h_w\,2\pi R_{co}} = \frac{1.58\times10^4}{3.0\times
10^4\times2\pi\times0.00475} = 17.6\,\si{\kelvin},
\]
so the clad outer surface sits at $317.6\,\degC$. The clad conduction drops
\[
\Delta T_c = \frac{q'\ln(R_{co}/R_{ci})}{2\pi k_c} = \frac{1.58\times10^4
\times\ln(4.75/4.2)}{2\pi\times17} = \frac{1.58\times10^4\times0.1230}{
106.8} = 18.2\,\si{\kelvin},
\]
lifting the clad inner surface to $335.8\,\degC$. The gap, a thin
low-conductance layer, drops
\[
\Delta T_{\text{gap}} = \frac{q'}{h_{\text{gap}}\,2\pi R_{ci}} =
\frac{1.58\times10^4}{6000\times2\pi\times0.0042} = 99.8\,\si{\kelvin},
\]
so the pellet surface runs at $435.6\,\degC$, the gap alone carrying nearly
$100\,\si{\kelvin}$ despite its thinness because its conductance is poor.
The last and largest drop is the conduction across the generating pellet
itself, which from the parabolic solution of the busbar example is
\[
T_0 - T_{s,f} = \frac{\dot{q}_v R_f^2}{4 k_f} = \frac{q'}{4\pi k_f} =
\frac{1.58\times10^4}{4\pi\times3.0} = 419\,\si{\kelvin}.
\]
The centreline therefore sits at $T_0 = 435.6 + 419 = 855\,\degC$. The
striking feature is where the temperature lives: of the $555\,\si{\kelvin}$
rise from coolant to centre, the pellet conduction carries $419$ and the
gap carries $100$, while the clad and the film together carry only $36$,
because the poor conductor and the poor-contact gap are the controlling
resistances and the good conductors are rounding terms, the exact inverse
of the busbar. The $855\,\degC$ centreline holds a large margin against
melting at nominal power, but the same calculation at an overpower
transient, with $k_f$ degraded by burn-up and the gap conductance changed
by pellet swelling, is what the fuel-performance codes track pin by pin,
because the centreline melt margin is the limit that caps reactor power.

\subsection{Freeze protection by electric heat tracing}

A pipe carrying a fluid that must not freeze or congeal is held above
ambient by an electric heat-tracing cable clipped to its wall under the
insulation, and sizing the trace is a steady balance between the trace
power in and the insulated heat loss out. We take a $50\,\si{\milli\meter}$
outer-diameter process line that must be kept at $T_p = 10\,\degC$ against
a design ambient of $T_a = -25\,\degC$, lagged with $40\,\si{\milli\meter}$
of insulation ($k = 0.035\,\text{W}/(\text{m}\cdot\si{\kelvin})$) and an
outer film $h_o = 20\,\text{W}/(\text{m}^2\cdot\si{\kelvin})$ against a
winter wind. We size the trace power per metre and choose a cable rating.

The heat the insulated pipe loses per metre is the pipe-to-ambient
difference over the series of the insulation and outer-film resistances.
The insulation runs from the pipe outer radius $r_i = 0.025\,\text{m}$ to
$r_o = 0.025 + 0.040 = 0.065\,\text{m}$, giving
\[
R_{\text{ins}} = \frac{\ln(r_o/r_i)}{2\pi k} = \frac{\ln(0.065/0.025)}{
2\pi\times0.035} = \frac{0.956}{0.220} = 4.35\,\si{\kelvin}\cdot\text{m}/
\text{W},
\]
and the outer film on the enlarged surface adds
\[
R_o = \frac{1}{h_o\,2\pi r_o} = \frac{1}{20\times2\pi\times0.065} =
0.122\,\si{\kelvin}\cdot\text{m}/\text{W},
\]
so the total is $R = 4.47\,\si{\kelvin}\cdot\text{m}/\text{W}$ and the loss
at the design condition is
\[
\frac{\dot{Q}}{L} = \frac{T_p - T_a}{R} = \frac{10 - (-25)}{4.47} =
\frac{35}{4.47} = 7.83\,\text{W}/\text{m}.
\]
The trace must replace exactly this loss to hold the setpoint, so a
$10\,\text{W}/\text{m}$ self-regulating cable, the smallest standard rating
above the requirement, carries the duty with a $28\,\%$ margin for cold
spots at supports and valves. Two design checks close the sizing. The
first is that the trace, sitting against the pipe wall, must not overheat
the fluid or itself: at $10\,\text{W}/\text{m}$ against $R = 4.47$ the
pipe wall settles only $\dot{Q}R = 10\times4.47 = 44.7\,\si{\kelvin}$
above ambient if the trace runs continuously with no fluid draw, so the
wall reaches $-25 + 44.7 = 20\,\degC$, safely below any degradation limit,
which is why self-regulating cable that cuts its own power as it warms is
preferred over constant-wattage cable that would keep driving. The second
is the sensitivity to the insulation: halving the lagging to
$20\,\si{\milli\meter}$ raises $r_o$ to $0.045\,\text{m}$, drops
$R_{\text{ins}}$ to $\ln(0.045/0.025)/(2\pi\times0.035) = 2.67$ and the
loss to $35/(2.67 + 0.177) = 12.3\,\text{W}/\text{m}$, which would force
the next cable rating up and raise the standing electricity cost for the
life of the line. The trace power and the insulation thickness trade
directly, and the economic optimum lags the pipe well enough that the
smallest standard cable suffices.

\subsection{Reading a surface temperature with a pyrometer}

A radiation thermometer infers a surface temperature from the thermal
radiation it collects, and the reading it returns is only as good as the
emissivity it assumes, a systematic error the Stefan-Boltzmann law lets us
size exactly. We take a single-band pyrometer that has been calibrated
against a blackbody, so it reports the temperature $T_{\text{app}}$ of a
blackbody that would emit the flux it measures. Pointed at a real surface
of emissivity $\epsilon$ and true temperature $T$, in a narrow wavelength
band the collected flux is $\epsilon$ times the blackbody flux at $T$, so
the instrument solves $\sigma_\lambda T_{\text{app}}^n = \epsilon\,
\sigma_\lambda T^n$ for the apparent temperature, where the effective
power $n$ is large in a short-wavelength band and near $4$ in a broadband
instrument. We use the broadband case for a clean estimate: the instrument
sees a flux $\epsilon\sigma T^4$ and reports $T_{\text{app}} =
\epsilon^{1/4} T$.

Take a hot steel billet whose true temperature is $T = 1000\,\degC =
1273\,\si{\kelvin}$, viewed with the pyrometer set for a blackbody
($\epsilon = 1$). Oxidised steel has an emissivity near $\epsilon = 0.80$,
so the instrument reads
\[
T_{\text{app}} = \epsilon^{1/4} T = 0.80^{1/4}\times1273 = 0.946\times
1273 = 1204\,\si{\kelvin} = 931\,\degC,
\]
$69\,\si{\kelvin}$ low, a large enough error to send a heat-treatment
process out of specification. Bright rolled steel is worse: at
$\epsilon = 0.30$ the reading is $0.30^{1/4}\times1273 = 0.740\times1273 =
942\,\si{\kelvin} = 669\,\degC$, a reading $331\,\si{\kelvin}$ below truth,
which would let an operator badly overheat the metal while trusting a cool
number. The size of the error tracks the quarter-power of emissivity, so a
surface whose emissivity is uncertain by a factor of two is uncertain in
apparent temperature by a factor $2^{1/4} = 1.19$ on the absolute scale.
The engineering responses read off the same relation: set the pyrometer's
emissivity dial to the surface's known value, prefer a two-colour
(ratio) pyrometer that divides two bands and cancels the emissivity to
first order, or contrive a near-blackbody viewing cavity such as a drilled
hole whose depth-to-diameter ratio drives its effective emissivity toward
unity. A thermocouple, by contrast, reads the true temperature of whatever
it touches and needs no emissivity, which is why the two instruments are
used together, the thermocouple to anchor the emissivity setting the
pyrometer then uses to scan a whole surface without contact.

\subsection{Vignette: a flat-plate solar-thermal collector}

A solar collector is the radiation and convection sections working against
each other on a single plate: sunlight is absorbed, and the same plate
loses heat to the cover and the sky by the very mechanisms the chapter has
developed. Sizing one is a combined-resistance balance with a source term,
and the device makes the abstract loss coefficient of section 5.5 into a
number a homeowner pays for. We take a glazed flat-plate collector of
absorber area $A = 2.0\,\text{m}^2$ receiving solar irradiance $G =
800\,\text{W}/\text{m}^2$, with a selective absorber coating of solar
absorptance $\tau\alpha = 0.85$ (the product of the glass transmittance
and the absorber absorptance) and an overall loss coefficient $U_L =
6.0\,\text{W}/(\text{m}^2\cdot\si{\kelvin})$, carrying water that enters at
$T_{\text{in}} = 40\,\degC$ on a $10\,\degC$ day.

The absorbed solar power is $\tau\alpha\,G\,A = 0.85\times800\times2.0 =
1360\,\text{W}$. The collector loses heat from its plate, running near the
inlet temperature, to ambient across the combined resistance the loss
coefficient represents,
\[
\dot{Q}_{\text{loss}} = U_L A\,(T_{\text{plate}} - T_a) \approx 6.0\times
2.0\times(40 - 10) = 360\,\text{W},
\]
so the useful heat delivered to the water is the absorbed power less the
loss,
\[
\dot{Q}_{\text{use}} = \tau\alpha\,G\,A - U_L A\,(T_{\text{in}} - T_a) =
1360 - 360 = 1000\,\text{W}.
\]
The instantaneous efficiency is the useful heat over the incident solar,
\[
\eta = \frac{\dot{Q}_{\text{use}}}{G A} = \frac{1000}{800\times2.0} =
0.625,
\]
which is exactly the Hottel-Whillier line $\eta = \eta_0 - U_L
(T_{\text{in}} - T_a)/G$ of Figure~\ref{fig:vol04ch05:collector-efficiency}
evaluated at the reduced temperature $(T_{\text{in}} - T_a)/G = 30/800 =
0.0375\,\si{\kelvin}\cdot\text{m}^2/\text{W}$, with the optical intercept
$\eta_0 = \tau\alpha = 0.85$ and slope $U_L = 6.0$. The plot makes the
design tension plain: as the collector runs hotter, to deliver hot water
rather than warm, the reduced temperature climbs, the loss term grows, and
the efficiency falls along the line until at the zero-efficiency intercept
$(T_{\text{in}} - T_a)/G = \eta_0/U_L = 0.85/6.0 = 0.142$ the plate loses
as fast as it absorbs and delivers nothing.

% Figure: the efficiency curve of a flat-plate solar-thermal collector,
% instantaneous efficiency against the reduced temperature difference
% (T_in - T_ambient)/G. The straight Hottel-Whillier line eta = eta0 -
% U_L (dT/G) with an optical intercept and a loss slope, three collector
% grades drawn as three slopes. Computed from the linear model with fixed
% constants, not measured data.
\begin{figure}[htbp]
\centering
\begin{tikzpicture}
\begin{axis}[
  width=12cm, height=7.5cm,
  xlabel={reduced temperature $(T_{\text{in}} - T_a)/G$ ($\si{\kelvin}\cdot\text{m}^2/\text{W}$)},
  ylabel={collector efficiency $\eta$},
  xmin=0, xmax=0.10, ymin=0, ymax=0.85,
  axis lines=left,
  tick label style={font=\scriptsize},
  label style={font=\small},
  legend style={font=\scriptsize, at={(0.02,0.02)}, anchor=south west},
  legend cell align=left,
  legend entries={unglazed ($U_L=20$),single-glazed ($U_L=6$),evacuated tube ($U_L=2$)},
  clip=false,
]
% unglazed: eta0 = 0.80, U_L = 20 -> zero at 0.040
\addplot[engamber, very thick, domain=0:0.040, samples=2] {0.80 - 20*x};
% single glazed: eta0 = 0.75, U_L = 6 -> zero at 0.125 (clipped at xmax)
\addplot[engblue, very thick, domain=0:0.10, samples=2] {0.75 - 6*x};
% evacuated tube: eta0 = 0.70, U_L = 2 -> zero at 0.35 (far right)
\addplot[engred, very thick, domain=0:0.10, samples=2] {0.70 - 2*x};
% operating window markers
\draw[enggray, dotted] (axis cs:0.05,0) -- (axis cs:0.05,0.85);
\node[font=\scriptsize, enggray, anchor=south west, rotate=90]
  at (axis cs:0.05,0.10) {domestic hot water};
\end{axis}
\end{tikzpicture}
\caption[Flat-plate solar-collector efficiency curve]{The instantaneous
efficiency of three grades of solar-thermal collector, plotted against the
reduced temperature difference $(T_{\text{in}} - T_a)/G$ that folds the
inlet-to-ambient gap and the incident irradiance into one variable. Each
line is the Hottel-Whillier form $\eta = \eta_0 - U_L (T_{\text{in}} -
T_a)/G$: the intercept $\eta_0$ is the optical efficiency, the fraction of
sunlight the glazing and the absorber capture at zero temperature
difference, and the slope $U_L$ is the collector's overall loss
coefficient, the same combined-resistance quantity that governs a wall.
An unglazed collector starts high but loses steeply, so it works only for
near-ambient duty such as pool heating; a single glazing cuts the loss
slope by trapping the convective and radiative losses behind a cover; an
evacuated tube removes the gas conduction and convection almost entirely
and holds useful efficiency out to the high reduced temperatures that
domestic hot water and process heat demand. The design reads straight off
the crossing points: below each line's zero-efficiency intercept the
collector delivers heat, above it the losses exceed the capture and the
collector runs backward.}
\label{fig:vol04ch05:collector-efficiency}
\end{figure}


The loss coefficient $U_L$ is the whole game, and it is a combined
resistance of the same kind the chapter has built repeatedly: the top loss
runs from the hot absorber through the air gap to the glass cover by
convection and radiation in parallel, then from the cover to the sky by
convection and radiation again, all in series, while the back and edges
lose through their insulation. A bare unglazed absorber has a $U_L$ near
$20$ and its efficiency line plunges, useful only for pool heating near
ambient; a single glazing traps the top convection and radiation and cuts
$U_L$ to about $6$, the value used here; an evacuated tube removes the gas
conduction and convection almost entirely and pushes $U_L$ toward $2$,
holding useful efficiency out to the high reduced temperatures that space
heating and process heat demand. The temperature rise of the water follows
from the useful heat and the flow: at a modest $0.02\,\text{kg/s}$ the rise
is $\dot{Q}_{\text{use}}/(\dot{m}c_p) = 1000/(0.02\times4180) =
12.0\,\si{\kelvin}$, so the water leaves near $52\,\degC$, and a slower
flow would raise the outlet temperature but also the mean plate
temperature and so the loss, the same effectiveness trade the heat
exchanger showed. The collector is the chapter's machinery assembled into
a device whose datasheet is one straight line whose intercept is optics
and whose slope is a loss coefficient.

\subsection{Vignette: the heat pipe as a super-conductor of heat}

A heat pipe carries heat along its length at an effective conductivity
tens to hundreds of times that of solid copper, and it does so with no
moving parts by running a closed evaporation-and-condensation cycle inside
a sealed tube (Figure~\ref{fig:vol04ch05:heat-pipe}). The device is the
condensation and boiling regimes of section 5.3 wired back to back: heat
entering the evaporator end boils a working fluid off a wick, the vapour
flows down the low-resistance core to a cooler condenser end where it gives
up its latent heat, and capillary action in the wick returns the condensate
to the evaporator. The transport step moves latent heat rather than
sensible heat, so the tube runs at nearly one temperature along its whole
length, set by the working fluid's saturation curve.

% Figure: schematic of a heat pipe. A sealed tube with an evaporator end,
% an adiabatic transport section, and a condenser end; the wick lines the
% wall, vapour flows down the core, liquid returns in the wick. Arrows for
% vapour (core) and liquid (wick). Schematic, not to scale.
\begin{figure}[htbp]
\centering
\begin{tikzpicture}[every node/.style={font=\small}]
% outer tube
\draw[engblack, very thick] (0,0) rectangle (10,2);
% wick lining top and bottom
\draw[enggray, fill=enggray!20] (0,0) rectangle (10,0.35);
\draw[enggray, fill=enggray!20] (0,1.65) rectangle (10,2);
% vapour core arrow (evaporator -> condenser), rightward
\draw[engred, line width=2.5pt, ->] (1.2,1.0) -- (8.8,1.0);
\node[font=\scriptsize, engred, anchor=south] at (5.0,1.05) {vapour};
% liquid return in wick (condenser -> evaporator), leftward
\draw[engblue, line width=1.5pt, ->] (8.8,0.55) -- (1.2,0.55);
\node[font=\scriptsize, engblue, anchor=north] at (5.0,0.5) {liquid (wick)};
% heat in at evaporator (left)
\foreach \y in {0.3,0.7,1.1,1.5} {
  \draw[engamber, ->] (-0.9,\y) -- (-0.05,\y);
}
\node[font=\scriptsize, engamber, anchor=east] at (-0.9,0.9) {$\dot{Q}_{\text{in}}$};
% heat out at condenser (right)
\foreach \y in {0.3,0.7,1.1,1.5} {
  \draw[engblue, ->] (10.05,\y) -- (10.9,\y);
}
\node[font=\scriptsize, engblue, anchor=west] at (10.9,0.9) {$\dot{Q}_{\text{out}}$};
% region labels
\node[font=\scriptsize, anchor=north] at (1.6,-0.15) {evaporator};
\node[font=\scriptsize, anchor=north] at (5.0,-0.15) {adiabatic transport};
\node[font=\scriptsize, anchor=north] at (8.4,-0.15) {condenser};
\end{tikzpicture}
\caption[Schematic of a heat pipe]{A heat pipe carries heat by a closed
evaporation-condensation cycle inside a sealed tube. Heat entering the
evaporator end boils the working fluid off the wick; the vapour flows down
the low-resistance core to the cooler condenser end, gives up its latent
heat, and condenses; capillary action in the wick pumps the liquid back to
the evaporator against the small pressure difference. Because the transport
step is a phase change rather than sensible conduction, a heat pipe carries
heat along its length at an effective conductivity far above any solid
metal, at a nearly uniform temperature set by the working fluid's
saturation curve. The device fails when the heat load outruns the wick's
capillary pumping (the capillary limit) or boils the wick dry, so a heat
pipe has a rated power above which its low resistance collapses.}
\label{fig:vol04ch05:heat-pipe}
\end{figure}


The size of the effect follows from a resistance comparison. Take a copper
rod $200\,\si{\milli\meter}$ long and $8\,\si{\milli\meter}$ in diameter
carrying $\dot{Q} = 30\,\text{W}$ from one end to the other. Its area is
$A = \pi(0.004)^2 = 5.03\times10^{-5}\,\text{m}^2$, so its axial resistance
is $R_{\text{Cu}} = L/(kA) = 0.200/(400\times5.03\times10^{-5}) = 9.94\,
\si{\kelvin}/\text{W}$, and the $30\,\text{W}$ would drop nearly
$300\,\si{\kelvin}$ end to end, which no electronics package would survive.
A water heat pipe of the same length and diameter carries the same
$30\,\text{W}$ against an end-to-end drop of only a few kelvin, because the
vapour transport adds almost no resistance and the small drops that remain
are the evaporator and condenser films. Its effective conductivity, read
from $k_{\text{eff}} = \dot{Q}L/(A\,\Delta T)$ with $\Delta T = 4\,\si{
\kelvin}$, is $30\times0.200/(5.03\times10^{-5}\times4) = 2.98\times10^4\,
\text{W}/(\text{m}\cdot\si{\kelvin})$, seventy-five times copper. The gain
is why heat pipes spread the load from a processor to a remote fin bank in
a laptop, move heat off orbiting instruments to a radiator panel, and drain
heat from permafrost-embedded pipeline supports to keep the ground frozen.

The device carries two limits an engineer must respect. The capillary limit
is the heat load at which the wick can no longer pump the returning liquid
fast enough against the vapour flow and the gravitational head; above it the
evaporator dries and its resistance jumps, the same burnout the boiling
curve of section 5.3 describes, so a heat pipe has a rated power beyond which
its low resistance collapses without warning. The temperature limit is set
by the working fluid: water pipes work from a little above freezing to about
$200\,\degC$, ammonia serves the cold end for spacecraft, and liquid-metal
pipes (sodium, potassium) carry the high-temperature end for reactor and
furnace duty, each fluid chosen so the operating temperature sits on the
steep part of its saturation curve where the latent-heat transport is
strong. The heat pipe is the chapter's phase-change convection turned into
a passive conductor, and its datasheet is a maximum power and a temperature
window rather than a single conductivity.

\subsection{Vignette: the wind-chill correction and a bare hand}

Wind-chill is Newton's law of cooling read as a felt temperature, and
working the numbers shows both why a breeze feels so much colder than still
air and why the effect saturates. Exposed skin loses heat to the air at
$\dot{q} = h(T_s - T_a)$, and the coefficient $h$ climbs sharply with wind
speed as forced convection replaces free convection. We take a bare hand at
skin temperature $T_s = 30\,\degC$ in air at $T_a = -5\,\degC$ and follow
the heat loss as the wind rises. In still air the free-convection
coefficient over a hand-sized surface is about $h_0 = 8\,\text{W}/(\text{m}
^2\cdot\si{\kelvin})$, so the loss is $8\times35 = 280\,\text{W}/\text{m}^2$.
A $5\,\text{m/s}$ breeze over a $0.10\,\text{m}$ characteristic length gives
a Reynolds number $\text{Re} = uL/\nu = 5\times0.10/1.3\times10^{-5} =
3.8\times10^4$, and the flat-plate laminar correlation $\text{Nu} = 0.664\,
\text{Re}^{1/2}\text{Pr}^{1/3}$ returns $\text{Nu} = 0.664\times196\times
0.887 = 115$, so $h = \text{Nu}\,k/L = 115\times0.024/0.10 = 27.6\,\text{W}
/(\text{m}^2\cdot\si{\kelvin})$, more than triple the still-air value. The
loss triples to $27.6\times35 = 966\,\text{W}/\text{m}^2$.

The felt temperature is the still-air temperature that would produce the
same loss, so setting $h_0(T_s - T_{\text{felt}}) = h(T_s - T_a)$ gives
\[
T_{\text{felt}} = T_s - \frac{h}{h_0}(T_s - T_a) = 30 - \frac{27.6}{8}
\times35 = 30 - 121 = -91\,\degC,
\]
a felt temperature far below the air temperature, which overstates the
published wind-chill because the skin does not in fact hold $30\,\degC$
once the loss climbs; blood flow drops and the skin cools, cutting the
driving difference. The saturation the standard wind-chill tables show
comes from two effects the crude balance misses: the convection coefficient
rises only as $u^{1/2}$ in the laminar range, so doubling the wind past a
brisk breeze adds less than the first breeze did, and the skin temperature
falls to defend the body, shrinking $T_s - T_a$. The engineering lesson
carries beyond frostbite: any surface losing heat to a moving fluid has a
loss set by the coefficient far more than by the temperature difference, so
a windbreak that cuts $h$ protects a pipe, a person, or a battery as much as
raising the ambient temperature would, and the same $u^{1/2}$ scaling means
the first reduction in exposure buys the most.

\subsection{A surface driven by a constant flux}

The step-temperature transient of section 5.2 fixed the surface temperature
and watched the flux fall; the complementary problem fixes the flux and
watches the surface temperature climb, and it is the right model for a
surface heated by an absorbed radiation load, a resistive film, or a friction
band, where the heat delivered is set and the surface temperature is the
unknown. We take a semi-infinite solid heated from $t = 0$ by a steady flux
$q'' = 1.0\times10^4\,\text{W}/\text{m}^2$ and ask how hot the surface runs
after five minutes for a metal against an insulator, and what property decides
the difference.

The heat equation with a constant-flux boundary
$-k\,\partial T/\partial x|_0 = q''$ has the exact surface solution
\[
\Delta T_s(t) = \frac{2 q''}{k}\sqrt{\frac{\alpha t}{\pi}} =
\frac{2 q''}{\sqrt{\pi}}\,\frac{\sqrt{t}}{\sqrt{k\rho c_p}},
\]
where the second form collects the three properties into the
\engterm{thermal effusivity} $e = \sqrt{k\rho c_p}$, the single group that
governs how a surface responds to an applied load. Take copper, $k = 400$,
$\rho = 8960$, $c_p = 385$, so $e = \sqrt{400\times8960\times385} =
\sqrt{1.38\times10^9} = 3.7\times10^4\,\text{J}/(\text{m}^2\cdot\si{\kelvin}
\cdot\text{s}^{1/2})$, and a rigid closed-cell foam, $k = 0.03$, $\rho = 40$,
$c_p = 1400$, so $e = \sqrt{0.03\times40\times1400} = \sqrt{1680} = 41\,
\text{J}/(\text{m}^2\cdot\si{\kelvin}\cdot\text{s}^{1/2})$, nearly three
orders below copper. At $t = 300\,\text{s}$ the common factor is
$2 q''\sqrt{t}/\sqrt{\pi} = 2\times10^4\times\sqrt{300}/1.772 = 1.955\times
10^5$, so copper rises $1.955\times10^5/3.7\times10^4 = 5.3\,\si{\kelvin}$
while the foam rises $1.955\times10^5/41 = 4.8\times10^3\,\si{\kelvin}$,
a figure the foam would never reach because it chars first. The physical
lesson the effusivity carries is that the surface temperature under a fixed
load depends on all three properties together, not on conductivity alone: the
foam conducts poorly and so cannot carry the heat inward, but it also stores
little per unit volume, so the thin surface layer it does heat rockets up in
temperature. The curves of Figure~\ref{fig:vol04ch05:constant-flux-surface}
show the two responses side by side, both climbing as $\sqrt{t}$ and separated
by the effusivity ratio. The same group explains why a metal doorknob and a
wooden one in the same cold room feel so different to the touch: the finger
delivers a roughly fixed flux, and the high-effusivity metal barely warms
under it while the low-effusivity wood warms fast at its surface, so the metal
draws heat from the skin and feels cold while the wood feels neutral. The
effusivity, not the conductivity, is the property of touch and of any
short-time surface transient, and it is why a heat spreader is specified by
$\sqrt{k\rho c_p}$ as much as by $k$.

% Figure: surface temperature rise of a semi-infinite solid heated by a
% constant surface flux, which climbs as the square root of time, contrasted
% with the falling surface flux of the step-temperature case. The curve is
% the exact result Delta T_s(t) = (2 q'' / k) sqrt(alpha t / pi), plotted
% here for a metal and an insulator to show the material sets the rate.
\begin{figure}[htbp]
\centering
\begin{tikzpicture}
\begin{axis}[
  width=12cm, height=7cm,
  xlabel={time $t$ (s)},
  ylabel={surface temperature rise $\Delta T_s$ (K)},
  xmin=0, xmax=300,
  ymin=0, ymax=260,
  axis lines=left,
  tick label style={font=\scriptsize},
  label style={font=\small},
  legend style={font=\scriptsize, at={(0.02,0.98)}, anchor=north west},
  legend cell align=left,
  legend entries={
    insulator ($\sqrt{k\rho c_p}=600$),
    metal ($\sqrt{k\rho c_p}=3.7\times10^{4}$)},
  clip=false,
]
% Delta T_s = (2 q''/sqrt(pi)) * sqrt(t) / sqrt(k rho c_p) = C * sqrt(t)
% Take q'' = 1e4 W/m^2. For the insulator sqrt(k rho c_p)=600 -> C = 2*1e4/(sqrt(pi)*600)=18.8
% For the metal sqrt(k rho c_p)=3.7e4 -> C = 2*1e4/(sqrt(pi)*3.7e4)=0.305
\addplot[engred, very thick, domain=0:300, samples=60] {18.8*sqrt(x)};
\addplot[engblue, very thick, domain=0:300, samples=60] {0.305*sqrt(x)};
\node[font=\scriptsize, engred, anchor=west] at (axis cs:150,205) {$\propto \sqrt{t}$};
\end{axis}
\end{tikzpicture}
\caption[Surface temperature under a constant applied flux]{The surface
temperature of a semi-infinite solid driven by a steady applied flux
$q''$ climbs without bound as the square root of time,
$\Delta T_s(t) = (2 q''/k)\sqrt{\alpha t/\pi}$, the mirror of the
step-temperature case whose surface flux instead falls as
$t^{-1/2}$. The rate is set by the thermal effusivity
$\sqrt{k\rho c_p}$: a metal with a large effusivity barely warms under a
flux that drives an insulator hundreds of kelvin hot in the same time,
which is why a hand rests briefly on cool metal but not on cool foam
carrying the same absorbed sunlight, and why a heat spreader is chosen
for its effusivity as much as its conductivity.}
\label{fig:vol04ch05:constant-flux-surface}
\end{figure}


\begin{mastery}
The effusivity that governed the constant-flux surface also fixes the
temperature two touching bodies share the instant they meet, and the
derivation is the cleanest place to see why touch reads effusivity rather
than temperature. Take two semi-infinite solids, one at $T_A$ with effusivity
$e_A = \sqrt{k_A\rho_A c_{p,A}}$ and one at $T_B$ with $e_B$, brought into
perfect contact along a plane at $t = 0$. Each body responds to the sudden
contact as a semi-infinite solid whose surface is held at the still-unknown
interface temperature $T_i$, so the flux each draws from or feeds to the
interface is the step-response surface flux of section 5.2,
$q''(t) = e\,(T_{\text{surface}} - T_{\text{far}})/\sqrt{\pi t}$, with the
effusivity as the coefficient. Continuity of heat at the interface requires
the flux body $A$ delivers to equal the flux body $B$ receives,
\[
\frac{e_A (T_A - T_i)}{\sqrt{\pi t}} = \frac{e_B (T_i - T_B)}{\sqrt{\pi t}},
\]
and the $\sqrt{\pi t}$ cancels, so the interface temperature is constant in
time and set by the effusivity-weighted average
\[
T_i = \frac{e_A T_A + e_B T_B}{e_A + e_B}.
\]
The interface leaps at once to a fixed value and stays there, which is why a
touch is felt instantly and does not drift. The value sits nearer the
temperature of the higher-effusivity body: a finger at $33\,\degC$
($e \approx 1500$ for skin) touching cold copper ($e \approx 3.7\times10^4$)
at $10\,\degC$ drives the interface to $(1500\times33 + 3.7\times10^4\times10)
/(1500 + 3.7\times10^4) = 10.9\,\degC$, almost the copper temperature, so the
skin surface plunges and the metal feels cold. The same finger on cold wood
($e \approx 400$) reaches $(1500\times33 + 400\times10)/(1500 + 400) =
28.2\,\degC$, barely cooled, so the wood feels neutral at the identical
temperature. The perceived warmth or cold of a surface is therefore read from
the effusivity ratio $e_{\text{skin}}/e_{\text{surface}}$, not from the
temperature difference, which is why a tile floor and a rug at the same room
temperature feel so different underfoot, why a metal ice tray feels colder
than a plastic one from the same freezer, and why burn severity from a brief
contact scales with the surface's effusivity as much as its temperature. The
result also caps how fast a spreader can pull heat from a hot spot in the
first instant: the interface it forms with the die is fixed by the two
effusivities before conduction has time to develop, so a high-effusivity
spreader wins the transient as well as the steady state.
\end{mastery}

\subsection{Response time of a temperature probe}

A temperature sensor reports the past, not the present, because it must first
come to the temperature of what it measures, and the lag is a lumped-
capacitance transient the instrument engineer must size against the rate the
process changes. We take a bare thermocouple bead, modelled as a sphere of
diameter $d = 1.5\,\si{\milli\meter}$ with the properties of its alloy
($\rho = 8500\,\text{kg/m}^3$, $c_p = 400\,\text{J}/(\text{kg}\cdot\si{
\kelvin})$, $k = 25\,\text{W}/(\text{m}\cdot\si{\kelvin})$), plunged into a
gas stream at a convection coefficient $h = 120\,\text{W}/(\text{m}^2\cdot
\si{\kelvin})$, and find its response time and the error it makes reading a
ramping temperature.

The bead is small and metallic, so we check the Biot number first with the
sphere length scale $L_c = d/6 = 2.5\times10^{-4}\,\text{m}$:
$\text{Bi} = h L_c/k = 120\times2.5\times10^{-4}/25 = 1.2\times10^{-3}$, far
below $0.1$, so the bead is isothermal and the lumped model holds. The time
constant is
\[
\tau = \frac{\rho V c_p}{h A} = \frac{\rho L_c c_p}{h} = \frac{8500\times
2.5\times10^{-4}\times400}{120} = 7.1\,\text{s}.
\]
The bead follows a step in gas temperature as $1 - e^{-t/\tau}$, reaching
$63\,\%$ of the step in $7.1\,\text{s}$ and $99\,\%$ only after about five
time constants, $35\,\text{s}$, which is slow for a control loop that must
catch a fast excursion. The more revealing case is a steady ramp: if the gas
temperature climbs linearly at rate $R$, the first-order sensor settles into a
steady lag, tracking a fixed amount $R\tau$ behind the true value, because the
solution of $\tau\,dT_s/dt + T_s = Rt$ is $T_s = R(t - \tau) + $ (a decaying
transient). A bead with $\tau = 7.1\,\text{s}$ reading a gas warming at
$R = 5\,\si{\kelvin}/\text{s}$ therefore lags by $R\tau = 35.5\,\si{\kelvin}$,
a large and easily missed systematic error. The engineering responses read
off the time constant: shrink the bead, since $\tau$ falls with $L_c$ and so
with diameter, raise the convection coefficient by exposing the bead to the
flow rather than shielding it, or correct the reading in software by adding
$\tau\,dT_s/dt$ back to the measured value, the standard first-order lead
compensation. A sheathed or grounded probe adds the sheath's mass and a
contact resistance in series, lengthening $\tau$ into the tens of seconds,
which is the trade of ruggedness against speed every probe choice makes.

\subsection{An insulating air gap on the edge of convection}

An enclosed gas layer insulates only while it stays still, and the design
question for a window cavity or a wall gap is how wide it can be made before
the trapped gas begins to convect and carries heat across faster than the
extra width blocks it. We take a vertical air gap between two glass panes,
the warm pane at $T_1 = 20\,\degC$ and the cool pane at $T_2 = 0\,\degC$, and
find the gap width at which convection sets in and what it does to the
conductance.

The gap Rayleigh number built on the gap width $L$ is $\text{Ra}_L =
g\beta(T_1 - T_2)L^3/(\nu\alpha)$. For air at the mean gap temperature
$10\,\degC = 283\,\si{\kelvin}$ take $\beta = 1/283 = 3.53\times10^{-3}\,
\si{\kelvin}^{-1}$, $\nu = 1.4\times10^{-5}\,\text{m}^2/\text{s}$, $\alpha =
2.0\times10^{-5}\,\text{m}^2/\text{s}$, and $k = 0.025\,\text{W}/(\text{m}
\cdot\si{\kelvin})$. Convection in a tall vertical cavity sets in near a
critical Rayleigh number of about $\text{Ra}_c \approx 2000$, so the onset
width solves
\[
L_c^3 = \frac{\text{Ra}_c\,\nu\alpha}{g\beta\,\Delta T} = \frac{2000\times
1.4\times10^{-5}\times2.0\times10^{-5}}{9.81\times3.53\times10^{-3}\times20},
\]
giving $L_c^3 = 5.60\times10^{-7}/0.693 = 8.08\times10^{-7}\,\text{m}^3$, so
$L_c = 9.3\times10^{-3}\,\text{m}$, about nine millimetres. Below this width
the gap conducts as a still layer with $\text{Nu} = 1$ and a conductance
$k/L$ that falls as the gap widens, the design rewarding a wider gap. At the
onset the reward stops: past $L_c$ the Nusselt number rises as roughly
$\text{Ra}^{1/4}$ (Figure~\ref{fig:vol04ch05:enclosed-cavity}), and since
$\text{Ra} \propto L^3$ the convecting conductance $\text{Nu}\,k/L \propto
L^{3/4}/L = L^{-1/4}$ falls only weakly, so widening the gap past onset buys
almost nothing. Compare a $9\,\si{\milli\meter}$ gap, still conducting at
$k/L = 0.025/0.009 = 2.78\,\text{W}/(\text{m}^2\cdot\si{\kelvin})$, against a
$20\,\si{\milli\meter}$ gap: its Rayleigh number is
$2000\times(20/9.3)^3 = 2000\times9.9 = 1.98\times10^4$, so
$\text{Nu} = 0.197\times(1.98\times10^4)^{1/4} = 0.197\times11.9 = 2.34$ and
the conductance is $2.34\times0.025/0.020 = 2.93\,\text{W}/(\text{m}^2\cdot
\si{\kelvin})$, slightly worse than the narrower still gap despite carrying
more than twice the width. The design rule the calculation gives is to hold
the gap at or just below the onset width, near a centimetre for air, which is
why sealed double-glazed units cluster at gap widths of twelve to sixteen
millimetres rather than the wide cavity intuition would pick, and why filling
the gap with argon, whose lower $\alpha$ and higher $\nu$ raise $\text{Ra}_c$
onset to a wider gap, lets the unit run a slightly wider still layer.

% Figure: Nusselt number for natural convection in a vertical enclosed air
% gap versus the gap Rayleigh number. Below a critical Rayleigh number the
% gap conducts (Nu = 1, still layer); above it convection cells form and
% Nu climbs as a power law. Log-log axes, schematic of the two regimes.
\begin{figure}[htbp]
\centering
\begin{tikzpicture}
\begin{axis}[
  width=12cm, height=7cm,
  xmode=log, ymode=log,
  xlabel={gap Rayleigh number $\text{Ra}_L$},
  ylabel={gap Nusselt number $\text{Nu}$},
  xmin=10, xmax=1e6,
  ymin=0.8, ymax=20,
  axis lines=left,
  tick label style={font=\scriptsize},
  label style={font=\small},
  legend style={font=\scriptsize, at={(0.02,0.98)}, anchor=north west},
  legend cell align=left,
  legend entries={conduction floor $\text{Nu}=1$,convection branch $0.197\,\text{Ra}^{1/4}$},
  clip=false,
]
% conduction plateau Nu = 1 up to Ra ~ 2000
\addplot[engblue, very thick, domain=10:2000, samples=20] {1};
% convection branch Nu = 0.197 Ra^{1/4} (a representative vertical-cavity fit)
% at Ra=2000 -> 0.197*2000^0.25 = 0.197*6.69=1.32 ; small mismatch at knee, acceptable schematic
\addplot[engred, very thick, domain=2000:1e6, samples=40] {0.197*(x^0.25)};
\draw[enggray, dotted] (axis cs:2000,0.8) -- (axis cs:2000,20);
\node[font=\scriptsize, enggray, anchor=south, rotate=90] at (axis cs:2000,1.3) {onset $\text{Ra}_c\approx 2000$};
\node[font=\scriptsize, engblue, anchor=south] at (axis cs:60,1.05) {still-air layer};
\node[font=\scriptsize, engred, anchor=south east] at (axis cs:6e5,10) {convection cells};
\end{axis}
\end{tikzpicture}
\caption[Nusselt number of an enclosed air gap]{An enclosed air gap does
not always convect. Below a critical gap Rayleigh number near $2000$ the
buoyant force cannot overcome viscous damping, no cells form, and the gap
transfers heat by pure conduction with a Nusselt number of one, the still
layer that gives a window air gap and a foam pore their insulating value.
Above the onset, convection rolls appear and the Nusselt number climbs as
roughly $\text{Ra}^{1/4}$, so a gap widened past the onset stops earning
its width because the new convection carries the extra heat the added
conduction resistance was meant to block. The design reading is that an
insulating gas gap is kept narrow enough to sit on the conduction floor,
which is why sealed-unit windows fix the gap near the width that keeps
$\text{Ra}$ below onset rather than making it as wide as possible.}
\label{fig:vol04ch05:enclosed-cavity}
\end{figure}


\subsection{A fin with a convecting tip}

The fin analysis of section 5.5 took the tip as insulated, which is exact only
when the tip loses no heat, and a real fin convects from its end face too. We
work the convecting-tip correction and show why the simple corrected-length
trick recovers it almost exactly, on a fin where the tip term actually
matters. Take an aluminium pin fin ($k = 237$), a cylinder of diameter
$d = 6\,\si{\milli\meter}$ and length $L = 40\,\si{\milli\meter}$, base at
$\theta_b = 60\,\si{\kelvin}$ above ambient, in air with $h = 40\,\text{W}/
(\text{m}^2\cdot\si{\kelvin})$.

The fin parameter uses the pin's perimeter $P = \pi d = 0.01885\,\text{m}$
and cross-section $A_c = \pi d^2/4 = 2.827\times10^{-5}\,\text{m}^2$,
\[
m = \sqrt{\frac{hP}{kA_c}} = \sqrt{\frac{40\times0.01885}{237\times2.827
\times10^{-5}}} = \sqrt{112.5} = 10.6\,\text{m}^{-1}, \quad mL = 10.6\times
0.040 = 0.424.
\]
With an insulated tip the heat rate is $\dot{Q}_{\text{ins}} = \sqrt{hPkA_c}
\,\theta_b\tanh(mL)$. The convecting-tip solution replaces $\tanh(mL)$ by the
ratio
\[
\frac{\tanh(mL) + (h/mk)}{1 + (h/mk)\tanh(mL)},
\]
where the extra term $h/(mk) = 40/(10.6\times237) = 0.0159$ is the tip
convection measured against the fin's own conductance. Here $\tanh(0.424) =
0.400$, so the tip ratio is $(0.400 + 0.0159)/(1 + 0.0159\times0.400) =
0.4159/1.00636 = 0.4133$, against the insulated $0.400$, a $3.3\,\%$ uplift in
heat rate from the tip face. The corrected-length shortcut adds half the tip
thickness to the length, $L_c = L + d/4 = 0.040 + 0.0015 = 0.0415\,\text{m}$,
because a pin's tip area $\pi d^2/4$ spread over its perimeter $\pi d$ adds a
length $d/4$; then $mL_c = 10.6\times0.0415 = 0.440$ and $\tanh(0.440) =
0.414$, matching the exact convecting-tip factor to three figures. The lesson
the numbers carry is that the tip correction is small on a slender fin, where
the tip area is a sliver of the total, and grows on a short fat fin, where the
end face is a real fraction of the surface; the corrected-length rule captures
it with no extra machinery, and the insulated-tip formula of the mastery box
is the right default precisely because most useful fins are slender enough
that the tip term is a few percent.

\subsection{Sizing a radiant floor panel}

A radiant floor heats a room from a warm surface rather than a warm air
stream, and sizing it is a combined convection-and-radiation balance on a
horizontal plate that must deliver a set power without running too hot to walk
on. We take a room whose floor is held at a surface temperature $T_f$ and must
deliver $\dot{Q}/A = 80\,\text{W}/\text{m}^2$ to the space, the air and the
surrounding surfaces both near $T_r = 20\,\degC$, and find the floor
temperature and how the load splits between the two modes.

The floor loses heat upward by free convection from a warm horizontal plate
facing up and by radiation to the room surfaces, the two in parallel. The
radiation to a large room is the small-surface limit,
$\dot{q}_{\text{rad}} = \epsilon\sigma(T_f^4 - T_r^4)$, and with a painted
floor $\epsilon = 0.9$ the linearised coefficient near the working
temperature is $h_{\text{rad}} = 4\epsilon\sigma T_m^3$. The free convection
from a warm floor uses the horizontal-plate correlation $\text{Nu} =
0.15\,\text{Ra}^{1/3}$, whose cube-root Rayleigh dependence makes the
coefficient nearly independent of plate size and roughly
$h_{\text{conv}} \approx 1.5\,\Delta T^{1/3}$ for a warm floor in air. We
guess a floor rise of $\Delta T = 8\,\si{\kelvin}$, so $T_f = 28\,\degC$ and
the mean radiating temperature is $T_m \approx 297\,\si{\kelvin}$. The
radiation coefficient is $h_{\text{rad}} = 4\times0.9\times5.67\times10^{-8}
\times297^3 = 5.35\,\text{W}/(\text{m}^2\cdot\si{\kelvin})$, and the
convection coefficient is $h_{\text{conv}} = 1.5\times8^{1/3} = 1.5\times2.0 =
3.0\,\text{W}/(\text{m}^2\cdot\si{\kelvin})$. The combined coefficient is
$8.35\,\text{W}/(\text{m}^2\cdot\si{\kelvin})$ and the delivered flux at
$\Delta T = 8$ is $8.35\times8 = 66.8\,\text{W}/\text{m}^2$, short of the
$80$ target, so the floor must run warmer. Raise the guess to $\Delta T =
9.5\,\si{\kelvin}$, $T_f = 29.5\,\degC$: the radiation coefficient barely
moves to $5.44$, the convection coefficient rises to $1.5\times9.5^{1/3} =
1.5\times2.12 = 3.18$, the combined $8.62$, and the flux is $8.62\times9.5 =
81.9\,\text{W}/\text{m}^2$, just over. The floor settles near $T_f =
29.4\,\degC$, marginally above the roughly $29\,\degC$ that
codes cap an occupied floor at, so the design sits at the comfort limit and
would need a lower supply temperature or a larger heated area to bring the
surface back under the cap, and the load splits with radiation carrying
$5.4/8.6 = 63\,\%$ and convection $37\,\%$. The split is the reason radiant
floors feel warm at a low surface temperature: most of the delivery is
radiant, warming the occupants and the room surfaces directly rather than the
air, so the air can be held a degree or two cooler for the same comfort, and
the low surface temperature keeps the floor gentle underfoot while still
meeting the heating load. The design lever is area, not temperature: a floor
that cannot deliver the room's peak load at its temperature cap needs more
heated area or a supplementary emitter, since pushing the surface hotter is
barred by comfort long before it is barred by physics.

\subsection{Boil-off from a cryogenic tank through its insulation}

The multilayer insulation of the spacecraft case blocked radiation into a
warm bus; the same stack blocks radiation \emph{into} a cold tank, and the
heat that leaks in boils off the stored cryogen at a rate the tank is bought
and sized against. We take a spherical liquid-nitrogen Dewar of inner
diameter $0.60\,\text{m}$ wrapped in a multilayer blanket of measured
effective emissivity $\epsilon^* = 0.03$, its inner surface at the nitrogen
boiling point $T_c = 77\,\si{\kelvin}$ and its outer surface near room
temperature $T_h = 295\,\si{\kelvin}$, and find the heat leak and the daily
boil-off fraction.

The blanket sees the warm outer wall on one side and the cold inner wall on
the other, so the radiative leak through the wrapped sphere of area $A =
\pi D^2 = \pi\times0.60^2 = 1.13\,\text{m}^2$ is
\[
\dot{Q} = \epsilon^* A \sigma (T_h^4 - T_c^4) = 0.03\times1.13\times5.67
\times10^{-8}\times(295^4 - 77^4).
\]
With $295^4 = 7.57\times10^9$ and $77^4 = 3.51\times10^7$, the bracket is
$7.54\times10^9$, so $\dot{Q} = 0.03\times1.13\times5.67\times10^{-8}\times
7.54\times10^9 = 14.5\,\text{W}$. That heat all goes to boiling nitrogen,
whose latent heat of vaporisation is $h_{fg} = 1.99\times10^5\,\text{J}/
\text{kg}$, so the boil-off rate is
\[
\dot{m} = \frac{\dot{Q}}{h_{fg}} = \frac{14.5}{1.99\times10^5} = 7.3\times
10^{-5}\,\text{kg/s},
\]
which over a day is $7.3\times10^{-5}\times86400 = 6.3\,\text{kg}$. Against a
full charge of liquid nitrogen ($\rho_l = 807\,\text{kg/m}^3$) filling the
$0.113\,\text{m}^3$ sphere, a mass of $91\,\text{kg}$, the daily loss is
$6.3/91 = 6.9\,\%$ per day, a representative figure for a well-insulated
laboratory Dewar. The number is the tank's headline specification, and the
blanket earns it: a bare double wall at $\epsilon = 0.2$ over the same
geometry would leak $0.2/0.03 = 6.7$ times as much, near $100\,\text{W}$,
boiling off almost half the charge a day. The ideal-shield reading of
Figure~\ref{fig:vol04ch05:cryostat-boiloff} shows why the blanket stops at a
few dozen layers rather than hundreds: the leak falls as $1/(N+1)$, so the
first ten shields win the bulk of the reduction and the spacer conduction sets
the floor the measured $\epsilon^*$ already folds in. The same balance sizes
the boil-off of liquefied-natural-gas ships and the hold time of a
liquid-helium magnet cryostat, where the far smaller latent heat of helium
makes even a milliwatt leak a serious loss and drives the use of a vapour-
cooled shield that intercepts the leak with the escaping cold gas before it
reaches the liquid.

% Figure: cryogenic tank boil-off. Left axis, the radiative heat leak
% through a multilayer blanket falls as 1/(N+1) with the number of ideal
% shields; right, the same leak converted to a daily boil-off fraction of
% a stored cryogen. Plotted as heat leak (W) versus number of shields on
% log-linear axes, with the bare-gap and ten-layer points marked.
\begin{figure}[htbp]
\centering
\begin{tikzpicture}
\begin{axis}[
  width=12cm, height=7cm,
  ymode=log,
  xlabel={number of radiation shields $N$},
  ylabel={radiative heat leak $\dot{Q}$ (W)},
  xmin=0, xmax=30,
  ymin=0.2, ymax=30,
  axis lines=left,
  tick label style={font=\scriptsize},
  label style={font=\small},
  legend style={font=\scriptsize, at={(0.98,0.98)}, anchor=north east},
  legend cell align=left,
  legend entries={ideal stack $\dot{Q}_0/(N+1)$},
  clip=false,
]
% Bare-gap leak Q0 = 20 W (schematic). Ideal stack: Q0/(N+1).
\addplot[engblue, very thick, domain=0:30, samples=31] {20/(x+1)};
\addplot[only marks, engred, mark=*, mark size=2pt] coordinates {(0,20) (10,1.82)};
\node[font=\scriptsize, engred, anchor=south west] at (axis cs:0,20) {bare gap};
\node[font=\scriptsize, engred, anchor=south west] at (axis cs:10,1.82) {10 shields};
\end{axis}
\end{tikzpicture}
\caption[Heat leak through a multilayer cryogenic blanket]{An ideal stack
of $N$ floating radiation shields cuts the bare-gap heat leak by a factor
$1/(N+1)$, so ten shields drop the leak roughly elevenfold and the curve
flattens: each further shield returns less than the one before, and a real
blanket stops improving once the spacer netting that separates the shields
conducts across them faster than the added shield blocks. This is why a
practical cryogenic Dewar quotes a measured effective conductivity for the
whole blanket rather than an ideal shield count, and why the last decade of
heat leak is fought with better spacers and fewer penetrations rather than
with more layers. The leak fixes the boil-off: the escaping heat divided by
the latent heat of the cryogen is the mass lost per unit time, the number a
tank is bought and sized against.}
\label{fig:vol04ch05:cryostat-boiloff}
\end{figure}


\subsection{Vignette: the thermos, the cool roof, and the passive daytime
radiator}

Three everyday surfaces run cold in the sun or hold heat against a cold room
by the same two levers the radiation section supplies, the emissivity in a
band and the ratio $\alpha_s/\epsilon$ across bands, and reading them together
shows how far a purely radiative design can be pushed. The ordinary vacuum
flask of the chapter's estimation runs its inner and outer walls at low
emissivity so the vacuum gap carries almost no radiation; that is a
same-band choice, cutting $\epsilon$ everywhere to throttle the only path a
vacuum leaves open, and it is why the flask holds coffee hot for a day. The
cool roof is a cross-band choice: it wears a finish with low solar
absorptance and high infrared emittance, so it reflects the incoming sunlight
it cannot use and radiates strongly in its own thermal band, holding a summer
roof tens of kelvin below the black roof of the estimation exercise, and the
selective-surface balance of the earlier example fixes the gap at the ratio
$\alpha_s/\epsilon$.

The third surface pushes the cross-band idea to its limit and reaches a
result that first looks impossible: a passive radiator that sits in full
daylight and cools below the air temperature with no power input. The trick is
to make the surface reflect almost all the sunlight, so its absorbed solar
load is small, while emitting strongly in the $8$ to $13\,\mu\text{m}$
atmospheric window, the band in which the sky is nearly transparent and the
effective radiating sink is the cold of deep space rather than the warm air.
A surface with solar absorptance $\alpha_s \approx 0.05$ under
$G_s = 900\,\text{W}/\text{m}^2$ absorbs only $0.05\times900 = 45\,\text{W}/
\text{m}^2$, while a high emittance in the window band lets it radiate to a
sky whose effective temperature in that band can sit $30$ to $40\,\si{\kelvin}$
below the air. The net radiative cooling is the window emission to the cold
sky less the absorbed sunlight and less the parasitic convection and radiation
from the warm air, and when the window emission wins, the surface settles
below ambient. A representative balance holds a few tens of watts per square
metre of net cooling at ambient temperature, or a surface several kelvin below
air with no net cooling delivered, which is the same $\alpha_s/\epsilon$ trade
the selective-surface example drew, now split further into an in-window and
out-of-window emittance so the radiator emits where the sky is open and stays
dark where it is not. The engineering reading is that radiation is not only a
loss to be blocked but a resource to be aimed: the thermos aims low emissivity
at a vacuum, the cool roof aims a band split at the sun, and the daytime
radiator aims a sharper band split at the one window in the atmosphere,
turning the sky into a cold sink a surface can dump heat into in broad
daylight.

\subsection{A transient probe: the thermal penetration time}

An engineer choosing where to place a temperature sensor, or judging how
fast a surface treatment will be felt at depth, needs the time for a
thermal disturbance to reach a given depth, and the semi-infinite result
gives it directly. The penetration depth grows as $\delta \sim \sqrt{\alpha
t}$, so the time to reach depth $d$ scales as $t \sim d^2/\alpha$, and the
factor in front is fixed by how much of the surface change one calls
``felt.'' We take the depth at which the temperature has moved to one
percent of the surface step, a common threshold, which the error-function
profile places at $\text{erf}(\eta) = 0.99$, or $\eta = d/(2\sqrt{\alpha t})
= 1.82$, so $t = d^2/(4\times1.82^2\,\alpha) = d^2/(13.2\,\alpha)$. Compare
three materials at a depth $d = 10\,\si{\milli\meter}$. Copper, with
$\alpha = 1.1\times10^{-4}\,\text{m}^2/\text{s}$, feels the change in
$t = (0.010)^2/(13.2\times1.1\times10^{-4}) = 0.069\,\text{s}$, faster than
a heartbeat. Glass, at $\alpha = 5\times10^{-7}$, takes
$(0.010)^2/(13.2\times5\times10^{-7}) = 15.2\,\text{s}$. A brick wall, at
$\alpha = 5\times10^{-7}$ over $d = 0.10\,\text{m}$, takes
$(0.10)^2/(13.2\times5\times10^{-7}) = 1.52\times10^3\,\text{s}$, about
twenty-five minutes, and a full masonry wall of $0.30\,\text{m}$ takes nine
times that, over three hours, which is the thermal-mass lag that keeps a
stone building cool through a hot afternoon. The quadratic dependence on
depth is the whole story: a sensor buried twice as deep reports the surface
event four times later, so a fast measurement demands a shallow, small
probe, and a slow deliberate lag, the phase shift a thick wall gives, demands
depth and low diffusivity together. The same $d^2/\alpha$ scaling tells a
cook that a thick roast must rest, its centre still climbing on stored
surface heat long after it leaves the oven, and tells a foundry that a thick
casting cools from the outside in over a time that grows with the square of
its section, the reason heavy sections trap the last liquid and the shrinkage
that goes with it.

\begin{mastery}
The steady, sinusoidally driven surface is the periodic cousin of the step
response, and it turns the diffusion equation into a single complex
wavenumber that governs both the amplitude decay and the phase lag a wall
imposes on a daily or annual temperature swing. Drive the surface of a
semi-infinite solid with $T(0,t) = \bar{T} + \Delta T\cos(\omega t)$ and
seek a solution $T(x,t) = \bar{T} + \operatorname{Re}[\Theta(x)e^{i\omega
t}]$. Substituting into $\partial T/\partial t = \alpha\,\partial^2 T/
\partial x^2$ collapses to $i\omega\Theta = \alpha\Theta''$, whose decaying
solution is $\Theta(x) = \Delta T\exp(-\sqrt{i\omega/\alpha}\,x)$. The
square root of $i$ is $(1+i)/\sqrt{2}$, so writing the diurnal penetration
depth $\delta = \sqrt{2\alpha/\omega}$ gives
\[
T(x,t) = \bar{T} + \Delta T\,e^{-x/\delta}\cos\!\left(\omega t -
\frac{x}{\delta}\right).
\]
The amplitude falls as $e^{-x/\delta}$ and the phase lags by $x/\delta$
radians, the two effects tied to the one length $\delta$. For soil with
$\alpha = 5\times10^{-7}\,\text{m}^2/\text{s}$ the daily period gives
$\delta = \sqrt{2\times5\times10^{-7}/(2\pi/86400)} = 0.117\,\text{m}$, so
the diurnal swing is down to $e^{-1} = 0.37$ of its surface value at
$0.12\,\text{m}$ and to a percent by half a metre, while the annual swing,
whose period is $365$ times longer, penetrates $\sqrt{365} = 19$ times
deeper, past three metres. The two depths explain the design of both the
cellar and the ground-source heat pump: a cellar a metre down never feels
the day but tracks the season, and a ground loop laid a few metres down
sees a nearly constant temperature all year, the annual wave having decayed
away. The phase lag $x/\delta$ also runs to a full period and beyond at
depth, so the deep soil is warmest in autumn and coldest in spring, half a
season out of step with the surface, a shift the same complex wavenumber
predicts exactly.
\end{mastery}

\begin{principle}[Heat transfer is resistance in series and parallel]
The working tool for any engineering heat-transfer problem is the
combined-resistance network: convection on each side of any surface,
conduction through each layer, radiation across any gap, with the
resistances combined in series for serial paths and in parallel for
parallel paths. The total heat-flow rate is the total temperature
difference divided by the total resistance. The discipline is to
identify which resistance dominates and to redesign that one; a
calculation that adds together resistances of vastly different
magnitudes (say, copper-wall conduction in series with still-air
convection) is signalling that the small resistance is
inconsequential and the design effort belongs at the large one.
\end{principle}

\begin{archetype}[Transport]
A transport phenomenon is the movement of a conserved quantity (mass,
momentum, energy, charge, species) under a driving gradient. Each
mode (conduction, convection, radiation for heat; diffusion,
advection for mass; momentum flux for fluid flow) has a flux
proportional to a gradient, with a transport coefficient (thermal
conductivity, viscosity, diffusivity) setting the proportionality.
The dimensionless groups of the chapter (Reynolds, Prandtl, Nusselt,
Grashof, Rayleigh) parameterise the relations between modes; the
mass-transfer analogue (Schmidt, Sherwood) in chapter~7 follows the
same template. The transport archetype recurs in Vol~III (momentum
in fluid flow), Vol~IV (heat and mass in this volume), and Vol~V
(species in reacting systems).
\end{archetype}

\begin{estimation}
\textbf{Estimate the rate of heat loss from a $1\,\text{L}$ thermos
flask of $90\,\degC$ coffee in a $20\,\degC$ kitchen, and the time
for the coffee to cool to $60\,\degC$.}

\smallskip
\textit{Estimate, before reading on.}

\smallskip
The thermos consists of an inner stainless-steel wall, an evacuated
gap with reflective silvered surfaces, and an outer wall. The
dominant heat-transfer modes:

\begin{itemize}[leftmargin=*]
  \item Radiation across the vacuum gap. Both surfaces are silvered
    with $\epsilon \approx 0.05$. Surface area $A \approx 0.025\,\text{m}^2$
    (a cylinder $0.3\,\text{m}$ tall, $0.08\,\text{m}$ diameter).
    Using the parallel-plate formula with $T_1 = 360\,\si{\kelvin}$ and
    $T_2 = 295\,\si{\kelvin}$ and effective emissivity $1/(1/\epsilon_1
    + 1/\epsilon_2 - 1) = 1/(20 + 20 - 1) = 0.0256$:
    \[
    \dot{Q}_{\text{rad}} = 0.0256 \times 0.025 \times 5.67 \times 10^{-8}
    \times (360^4 - 295^4) \approx 0.7\,\text{W}.
    \]
  \item Conduction through the cap and the support struts. These add
    perhaps another $1$ to $2\,\text{W}$.
\end{itemize}

Total heat-loss rate: $\sim 2\,\text{W}$. At the coffee mass
$1\,\text{kg}$ and $c = 4.18\,\text{kJ}/(\text{kg}\cdot\si{\kelvin})$,
the lumped-capacitance time constant is
\[
\tau = \frac{m c}{(\dot{Q}/\Delta T)} \approx \frac{4180}{2/70} \approx
1.5 \times 10^5\,\text{s} \approx 40\,\text{h}.
\]
The cooling from $90\,\degC$ to $60\,\degC$ corresponds to $\Delta T$
falling from $70$ to $40\,\si{\kelvin}$, a factor of $0.57$; the time
is $\tau \ln(70/40) \approx 1.5 \times 10^5 \times 0.56 \approx 8.4
\times 10^4\,\text{s} = 23\,\text{hours}$.

The published cooling figure for a good consumer-grade vacuum thermos
is roughly $10\,\si{\kelvin}/\text{hour}$ of cooling initially, dropping
as the temperature difference falls. Our estimate at $\sim 2.4\,\si{
\kelvin}/\text{hour}$ initially is a factor of four too low, suggesting
the actual heat-loss rate is closer to $5$ to $8\,\text{W}$. The
discrepancy is in the cap and the lid seal, which are larger losses
than our estimate. The vacuum-gap radiation calculation is good; the
working error is in the auxiliary paths.

The postmortem: the order of magnitude of the heat-loss rate
(a few watts) and the time scale (hours, not minutes or days) are
correctly estimated. The factor-of-three discrepancy is the level
expected from a working estimate of a heat-loss calculation with a
geometrically complex object.
\end{estimation}

\section{An end-to-end case study: cooling a desktop
processor}\pathtag{standard}

The chapter's resistance machinery pays off when an engineer must clear
a real part for service from a power number and a temperature limit. We
take a desktop processor and carry it from the dissipated power through
the junction-to-case-to-sink-to-ambient resistance network, the fin-array
sizing, the airflow the fan must supply, and a junction-temperature
margin against the rated limit, ending in a verdict an engineer would
sign. The processor is a good vehicle because every move of the chapter
appears in it once: conduction through a package and a thermal-interface
layer, spreading in a base, fin-array convection under forced air, and a
combined-resistance verdict.

\subsection{Specification and the temperature budget}

The processor dissipates $\dot{Q} = 60\,\text{W}$ at its design
operating point, spread over a $12\,\si{\milli\meter} \times 12\,\si{
\milli\meter}$ integrated heat spreader. The silicon vendor rates the
maximum junction temperature at $T_{j,\max} = 105\,\degC$. The chassis
draws air from a room at $25\,\degC$, and the air warms as it passes
through the case, so the air actually reaching the heat sink is at
$T_a = 35\,\degC$, a $10\,\si{\kelvin}$ case rise we carry as the
ambient for the sink. The whole budget is the temperature rise the
network is allowed to produce:
\[
T_{j,\max} - T_a = 105 - 35 = 70\,\si{\kelvin},
\]
and at $60\,\text{W}$ the total junction-to-ambient resistance the
design may not exceed is
\[
R_{ja,\text{allow}} = \frac{T_{j,\max} - T_a}{\dot{Q}} = \frac{70}{60}
= 1.17\,\si{\kelvin}/\text{W}.
\]
Every resistance from the junction outward must sum to less than this.

\subsection{The resistance network}

The path from the junction to the moving air is a series chain of four
resistances, drawn as the ladder of
Figure~\ref{fig:vol04ch05:heatsink-network}.

% Figure: the junction-to-ambient thermal-resistance ladder for the CPU
% heat-sink case study. A series chain of four resistances (junction-to-
% case, thermal interface, base spreading, fin-array convection) carries
% the dissipated power from the junction node to the ambient node, with
% each resistance labelled by its value and each node by its temperature.
\begin{figure}[htbp]
\centering
\begin{tikzpicture}[
  every node/.style={font=\small},
  res/.style={draw, engblue, very thick, minimum width=18mm, minimum height=8mm},
  term/.style={circle, draw, engblack, fill=enggray!15, minimum size=7mm, inner sep=1pt, font=\scriptsize},
]
% nodes (temperatures) along the chain
\node[term] (Tj) at (0,0) {$T_j$};
\node[res, right=7mm of Tj] (Rjc) {$R_{jc}$};
\node[res, right=7mm of Rjc] (Rtim) {$R_{\text{TIM}}$};
\node[res, right=7mm of Rtim] (Rsp) {$R_{\text{sp}}$};
\node[res, right=7mm of Rsp] (Rsa) {$R_{sa}$};
\node[term, right=7mm of Rsa] (Ta) {$T_a$};
% connecting wires
\draw[engblue, very thick] (Tj) -- (Rjc);
\draw[engblue, very thick] (Rjc) -- (Rtim);
\draw[engblue, very thick] (Rtim) -- (Rsp);
\draw[engblue, very thick] (Rsp) -- (Rsa);
\draw[engblue, very thick] (Rsa) -- (Ta);
% resistance values under each block
\node[font=\scriptsize, enggray, below=2mm of Rjc] {$0.30$};
\node[font=\scriptsize, enggray, below=2mm of Rtim] {$0.16$};
\node[font=\scriptsize, enggray, below=2mm of Rsp] {$0.10$};
\node[font=\scriptsize, enggray, below=2mm of Rsa] {$0.45$};
\node[font=\scriptsize, enggray, below=8mm of Rtim] {(all in $\si{\kelvin}/\text{W}$)};
% temperature labels above the terminal nodes
\node[font=\scriptsize, engred, above=2mm of Tj] {$95.8\,\degC$};
\node[font=\scriptsize, engred, above=2mm of Ta] {$35\,\degC$};
% the driving heat flow
\draw[engred, ->, thick] (-0.9,0.9) -- (-0.1,0.9);
\node[font=\scriptsize, engred, anchor=west] at (-0.9,1.25) {$\dot{Q}=60\,\text{W}$};
\end{tikzpicture}
\caption[Junction-to-ambient resistance ladder for the heat-sink case
study]{The junction-to-ambient thermal-resistance ladder for the
processor cooling case study, drawn as a series chain. The dissipated
power $\dot{Q}$ enters at the junction node $T_j$ and crosses four
resistances in turn: the junction-to-case resistance $R_{jc}$ fixed by
the package, the thermal-interface resistance $R_{\text{TIM}}$ of the
grease layer, the spreading resistance $R_{\text{sp}}$ of the heat-sink
base, and the sink-to-ambient resistance $R_{sa}$ of the fin array under
forced air. The temperature falls by $\dot{Q}R_i$ across each block, so
the junction sits above ambient by $\dot{Q}\sum_i R_i$. The single
largest block is the fin-array convection, which is where added fin area
or higher airflow buys the most margin.}
\label{fig:vol04ch05:heatsink-network}
\end{figure}


\begin{itemize}[leftmargin=*]
  \item \emph{Junction to case}, $R_{jc} = 0.30\,\si{\kelvin}/\text{W}$,
    a package property the vendor publishes. It is fixed once the part
    is chosen.
  \item \emph{Thermal-interface material}, $R_{\text{TIM}}$. A grease
    layer of conductivity $k_{\text{TIM}} = 4\,\text{W}/(\text{m}\cdot
    \si{\kelvin})$ and thickness $50\,\si{\micro\meter}$ over the
    $12\,\si{\milli\meter} \times 12\,\si{\milli\meter}$ spreader gives
    $R_{\text{TIM}} = t/(k A) = 50 \times 10^{-6}/(4 \times 0.012^2) =
    0.087\,\si{\kelvin}/\text{W}$. We carry $0.16\,\si{\kelvin}/\text{W}$
    to allow for two grease lines (spreader to base) and contact
    imperfection.
  \item \emph{Base spreading}, $R_{\text{sp}} = 0.10\,\si{\kelvin}/
    \text{W}$. The $12\,\si{\milli\meter}$ heat source must spread out
    to feed the full fin base; the spreading resistance for a small
    source on a larger aluminium base of this size is of this order.
  \item \emph{Sink to ambient}, $R_{sa}$, the fin-array convection,
    which is the term the designer controls and the one we size next.
\end{itemize}

The fixed part of the chain is $R_{jc} + R_{\text{TIM}} + R_{\text{sp}}
= 0.30 + 0.16 + 0.10 = 0.56\,\si{\kelvin}/\text{W}$, so the fin array
must deliver
\[
R_{sa} \leq R_{ja,\text{allow}} - 0.56 = 1.17 - 0.56 = 0.61\,\si{
\kelvin}/\text{W}.
\]

\subsection{Sizing the fin array}

We size an aluminium ($k = 237\,\text{W}/(\text{m}\cdot\si{\kelvin})$)
extruded sink with a $60\,\si{\milli\meter} \times 60\,\si{\milli\meter}$
base carrying $N$ rectangular fins, each $40\,\si{\milli\meter}$ tall,
$1.5\,\si{\milli\meter}$ thick, $60\,\si{\milli\meter}$ deep, with a fan
driving $3\,\text{m/s}$ of air through the channels. For air at the film
temperature take $k_{\text{air}} = 0.028\,\text{W}/(\text{m}\cdot\si{
\kelvin})$, $\nu = 1.7 \times 10^{-5}\,\text{m}^2/\text{s}$, $\text{Pr}
= 0.71$. Treating each inter-fin channel as flow over the fin walls of
length $L = 0.060\,\text{m}$,
\[
\text{Re}_L = \frac{u L}{\nu} = \frac{3 \times 0.060}{1.7 \times 10^{-5}}
= 1.06 \times 10^4,
\]
laminar over a flat plate, so the averaged correlation
$\text{Nu}_L = 0.664\,\text{Re}_L^{1/2}\,\text{Pr}^{1/3}$ gives
$\text{Nu}_L = 0.664 \times 103 \times 0.892 = 61$, and
\[
h = \frac{\text{Nu}_L\,k_{\text{air}}}{L} = \frac{61 \times 0.028}{0.060}
= 28.5\,\text{W}/(\text{m}^2\cdot\si{\kelvin}).
\]
Each fin has perimeter $P = 2(0.060 + 0.0015) = 0.123\,\text{m}$ and
cross-section $A_c = 0.060 \times 0.0015 = 9.0 \times 10^{-5}\,\text{m}^2$,
so the fin parameter is
\[
m = \sqrt{\frac{h P}{k A_c}} = \sqrt{\frac{28.5 \times 0.123}{237 \times
9.0 \times 10^{-5}}} = \sqrt{164} = 12.8\,\text{m}^{-1},
\]
and $mL = 12.8 \times 0.040 = 0.51$, a stubby effective fin with
efficiency $\eta_f = \tanh(mL)/(mL) = 0.470/0.51 = 0.92$. The heat-
transfer area per fin is the fin surface $P L = 0.123 \times 0.040 =
4.92 \times 10^{-3}\,\text{m}^2$, of which the efficiency makes
$\eta_f \times 4.92 \times 10^{-3} = 4.53 \times 10^{-3}\,\text{m}^2$
effective. With $N = 18$ fins plus the exposed base between them
(roughly $0.6 \times 0.060 \times 0.060 = 2.2 \times 10^{-3}\,\text{m}^2$
of unfinned base), the total effective area is
\[
A_{\text{eff}} = 18 \times 4.53 \times 10^{-3} + 2.2 \times 10^{-3} =
0.0837\,\text{m}^2,
\]
and the sink-to-ambient resistance is
\[
R_{sa} = \frac{1}{h A_{\text{eff}}} = \frac{1}{28.5 \times 0.0837} =
0.42\,\si{\kelvin}/\text{W},
\]
comfortably below the $0.61\,\si{\kelvin}/\text{W}$ budget. We carry
$0.45\,\si{\kelvin}/\text{W}$ as the value after derating for the
non-ideal channel entry and the bypass air that escapes around the fin
stack.

\subsection{Airflow and fan check}

The air carries the heat away, and its temperature rise sets the
$10\,\si{\kelvin}$ case ambient we assumed. The volumetric flow through
the fin channels is roughly the face velocity times the free channel
area. With $18$ fins of $1.5\,\si{\milli\meter}$ across a $60\,\si{
\milli\meter}$ width, the blocked width is $27\,\si{\milli\meter}$ and
the free width is $33\,\si{\milli\meter}$, so the channel face area is
$0.033 \times 0.040 = 1.32 \times 10^{-3}\,\text{m}^2$ and the flow is
$\dot{V} = u A = 3 \times 1.32 \times 10^{-3} = 3.96 \times 10^{-3}\,
\text{m}^3/\text{s}$, about $8\,\text{cfm}$. The air-temperature rise
across the sink is
\[
\Delta T_{\text{air}} = \frac{\dot{Q}}{\dot{m} c_p} = \frac{\dot{Q}}{
\rho \dot{V} c_p} = \frac{60}{1.16 \times 3.96 \times 10^{-3} \times
1007} = 12.9\,\si{\kelvin}.
\]
The air leaves about $13\,\si{\kelvin}$ warmer than it entered. The
$35\,\degC$ sink-ambient we assumed already folds in a $10\,\si{\kelvin}$
case rise; the fin-bundle correlation uses the inlet air, and the
$13\,\si{\kelvin}$ through-rise confirms the assumption is the right
size. A small axial fan rated for $8$ to $12\,\text{cfm}$ against the
modest pressure drop of a $40\,\si{\milli\meter}$ fin stack supplies
this flow with margin.

\subsection{Junction temperature and verdict}

The junction temperature is the ambient plus the dissipated power times
the summed resistance:
\[
T_j = T_a + \dot{Q}\,(R_{jc} + R_{\text{TIM}} + R_{\text{sp}} + R_{sa})
= 35 + 60 \times (0.30 + 0.16 + 0.10 + 0.45).
\]
The summed resistance is $1.01\,\si{\kelvin}/\text{W}$, so
\[
T_j = 35 + 60 \times 1.01 = 35 + 60.6 = 95.6\,\degC.
\]
Against the rated $105\,\degC$ the design has a margin of $9.4\,\si{
\kelvin}$, drawn as the budget stack of
Figure~\ref{fig:vol04ch05:junction-margin}.

% Figure: the junction-temperature budget for the heat-sink case study,
% drawn as a stacked bar. Each segment is one resistance's temperature
% rise above ambient; the stack top is the junction temperature, compared
% against the rated limit drawn as a horizontal line. The gap between the
% stack top and the limit is the design margin.
\begin{figure}[htbp]
\centering
\begin{tikzpicture}
\begin{axis}[
  width=10cm, height=8cm,
  ybar stacked,
  bar width=34pt,
  ymin=0, ymax=112,
  xmin=0, xmax=2,
  ylabel={temperature ($\degC$)},
  xtick=\empty,
  axis lines=left,
  legend pos=outer north east,
  legend cell align=left,
  tick label style={font=\scriptsize},
  label style={font=\small},
  legend style={font=\scriptsize},
  clip=false,
]
% segments stack from ambient (35 C) upward; values are dot Q * R_i
\addplot[fill=enggray!30, draw=enggray] coordinates {(1,35)};
\addlegendentry{ambient $35\,\degC$}
\addplot[fill=engblue!70, draw=engblue] coordinates {(1,18)};
\addlegendentry{$R_{jc}$ rise $18.0$}
\addplot[fill=engblue!40, draw=engblue] coordinates {(1,9.6)};
\addlegendentry{$R_{\text{TIM}}$ rise $9.6$}
\addplot[fill=englime!60!black, draw=englime!50!black] coordinates {(1,6)};
\addlegendentry{$R_{\text{sp}}$ rise $6.0$}
\addplot[fill=engamber!70, draw=engamber] coordinates {(1,27)};
\addlegendentry{$R_{sa}$ rise $27.0$}
% the rated junction limit
\draw[engred, very thick, dashed] (axis cs:0.45,105) -- (axis cs:1.55,105);
\node[font=\scriptsize, engred, anchor=south] at (axis cs:1.0,105.5) {rated limit $105\,\degC$};
% the achieved junction temperature line and label
\draw[engblack, thick] (axis cs:0.45,95.6) -- (axis cs:1.55,95.6);
\node[font=\scriptsize, engblack, anchor=north] at (axis cs:1.0,95.0) {junction $95.6\,\degC$};
% the margin bracket
\draw[engred, <->] (axis cs:1.7,95.6) -- (axis cs:1.7,105);
\node[font=\scriptsize, engred, anchor=west] at (axis cs:1.72,100.3) {margin $9.4\,\si{\kelvin}$};
\end{axis}
\end{tikzpicture}
\caption[Junction-temperature budget for the heat-sink case
study]{The junction-temperature budget for the processor cooling case
study, stacked from ambient upward. Each segment is the temperature rise
$\dot{Q}R_i$ across one resistance in the ladder of
Figure~\ref{fig:vol04ch05:heatsink-network}. The stack reaches a
junction temperature of $95.6\,\degC$ against a rated limit of
$105\,\degC$, leaving a margin of $9.4\,\si{\kelvin}$. The fin-array
convection segment is the tallest, so the cheapest way to widen the
margin is to lift the airflow or the fin area rather than to chase the
package or interface terms. A design that landed the stack top above the
dashed limit would fail the check and force a larger sink, a faster fan,
or a lower clock.}
\label{fig:vol04ch05:junction-margin}
\end{figure}


The verdict is that the
sink passes at the design power with a thin but real margin. The stack
also shows where the margin lives: the fin-array convection is the
single largest segment at $27\,\si{\kelvin}$ of rise, so a designer who
needed more headroom would raise the airflow or add fins before touching
the package or the interface. A worst-case audit would test the same
network at a hotter room ($30\,\degC$ rather than $25$), at the fan's
end-of-life reduced flow, and at a boosted power state; any of these can
erase the $9.4\,\si{\kelvin}$, which is why production designs target a
larger margin than this baseline and why the throttling logic that drops
the clock when the junction nears the limit is the last line of defence
rather than the first.

\section{A second end-to-end case study: a house retrofit and its
payback}\pathtag{standard}

The processor case carried a single part from a power number to a
temperature margin. A building inverts the problem: the temperatures are
fixed by comfort and climate, and the engineering question is how much
energy the envelope leaks over a season and whether a retrofit pays for
itself. We take a detached two-storey house in a cold-winter climate and
carry it from the component conductances through the seasonal degree-day
energy, the thermal-bridge correction the clear-field calculation misses,
the upgrade that targets the worst components, and a simple payback against
the retrofit cost, ending in a recommendation a homeowner could act on.
Every tool of the combined-resistance and U-value sections appears once:
series resistance through layered walls, parallel paths through framing,
the infiltration term, and the conversion of a steady conductance into an
annual energy through the degree-day integral.

\subsection{The envelope and its conductances}

The house has the envelope below, each component specified by its area and
its clear-field U-value computed from the layer build-up by the series-
resistance method of section 5.5. The indoor set point is $21\,\degC$.

\begin{center}
\begin{tabular}{lccc}
\toprule
component & area $A$ (m$^2$) & $U$ (W/m$^2$K) & $UA$ (W/K) \\
\midrule
walls & $180$ & $0.45$ & $81.0$ \\
windows & $32$ & $2.8$ & $89.6$ \\
roof/ceiling & $95$ & $0.20$ & $19.0$ \\
ground floor & $95$ & $0.35$ & $33.3$ \\
\bottomrule
\end{tabular}
\end{center}

The walls are a timber-framed cavity build (brick outer leaf, partial-fill
mineral wool, plasterboard inner lining) whose clear-field $U = 0.45$ comes
from summing the layer resistances and the two air films, as in the
composite-wall exercise. The windows are an older double-glazed unit at
$U = 2.8$. The roof is well insulated; the suspended ground floor is
moderately so. The clear-field transmission conductance is the sum,
\[
(UA)_{\text{fabric}} = 81.0 + 89.6 + 19.0 + 33.3 = 222.9\,\text{W}/\si{
\kelvin}.
\]

\subsection{The thermal-bridge correction}

The clear-field sum understates the loss, because the timber studs, the
floor and ceiling junctions, the window reveals, and the wall corners are
parallel high-conductivity paths that short-circuit the insulation
(Figure~\ref{fig:vol04ch05:thermal-bridge}). Building practice captures
these as linear thermal transmittances $\psi$ (W per metre of junction per
kelvin) summed over the junction lengths, plus a fraction added for the
repeating framing. For this house take a repeating-bridge uplift of
$15\,\%$ on the wall conductance (the framing fraction of a timber wall)
and a linear-bridge allowance of $\sum \psi L = 18\,\text{W}/\si{\kelvin}$
over the perimeter junctions and window reveals. The wall uplift adds
$0.15 \times 81.0 = 12.2\,\text{W}/\si{\kelvin}$, so the corrected fabric
conductance is
\[
(UA)_{\text{fabric}}^{\text{corr}} = 222.9 + 12.2 + 18.0 =
253.1\,\text{W}/\si{\kelvin}.
\]
The bridges add $30\,\text{W}/\si{\kelvin}$, an extra $13\,\%$ over the
clear-field figure, the systematic gap between a textbook U-value sum and
a measured whole-house loss that the project's wall measurement is built to
expose.

% Figure: a thermal bridge through a framed wall. The insulated cavity is a
% high-resistance path; the timber (or steel) stud is a parallel low-
% resistance path that short-circuits the insulation, so the heat-flow lines
% crowd through the stud and the inside surface runs cold over it. Schematic
% wall cross-section with crowding flux lines and an inset isotherm note.
\begin{figure}[htbp]
\centering
\begin{tikzpicture}[every node/.style={font=\small}]
% wall body: insulated cavity
\fill[engblue!8] (0,0) rectangle (6,3.2);
% two studs (low-resistance bridges)
\fill[engamber!30] (1.4,0) rectangle (1.8,3.2);
\fill[engamber!30] (4.2,0) rectangle (4.6,3.2);
% inner and outer face lines
\draw[very thick, engred] (0,0) -- (0,3.2);
\draw[very thick, engblue] (6,0) -- (6,3.2);
\node[font=\scriptsize, engred, anchor=south, rotate=90] at (-0.2,1.6) {warm inside};
\node[font=\scriptsize, engblue, anchor=north, rotate=90] at (6.2,1.6) {cold outside};
% flux lines through insulation: straight, sparse
\foreach \y in {0.5,1.1,1.7,2.3,2.9} {
  \draw[enggray, ->] (0.1,\y) -- (1.3,\y);
}
% flux lines crowding through the studs: denser, converging
\foreach \y in {0.4,0.8,1.2,1.6,2.0,2.4,2.8} {
  \draw[engred, thick, ->] (0.1,\y) .. controls (1.0,\y) and (1.4,1.6) .. (1.6,1.6);
}
\draw[engred, thick, ->] (1.6,1.6) -- (6,1.6);
% labels
\node[font=\scriptsize, anchor=south] at (3.0,3.25) {insulated cavity (high $R$)};
\node[font=\scriptsize, engamber, anchor=north] at (1.6,-0.1) {stud};
\node[font=\scriptsize, engamber, anchor=north] at (4.4,-0.1) {stud};
\node[font=\scriptsize, engred, anchor=west] at (6.4,2.4)
  {cold stripe on};
\node[font=\scriptsize, engred, anchor=west] at (6.4,2.05)
  {inside surface};
\end{tikzpicture}
\caption[Thermal bridge through a framed wall]{A thermal bridge through a
framed wall. The insulated cavity is the high-resistance path the design
intends; the framing studs are a parallel low-resistance path that
short-circuits the insulation. Heat-flow lines crowd through the studs, so
the outward flux is locally high and the inside surface runs cold in a
stripe over each member, the pattern an infrared camera reads as the
cold-spot diagnosis of the failure-analysis exercise. The bridge is
quantified by a linear or point transmittance added to the clear-field
$UA$, and a continuous layer of exterior insulation over the framing is the
standard cure because it covers the bridge rather than interrupting it.}
\label{fig:vol04ch05:thermal-bridge}
\end{figure}


\subsection{Infiltration and the total loss coefficient}

Air leakage carries heat out with the exfiltrating air. The heated volume
is $V = 320\,\text{m}^3$ and the infiltration rate is $0.6$ air-changes
per hour. The infiltration conductance is
\[
(UA)_{\text{inf}} = \dot{V}\rho c_p = \frac{0.6 \times 320}{3600} \times
1.2 \times 1005 = 64.3\,\text{W}/\si{\kelvin}.
\]
The whole-house loss coefficient is the sum of fabric and infiltration,
\[
(UA)_{\text{tot}} = 253.1 + 64.3 = 317\,\text{W}/\si{\kelvin}.
\]
At the design outdoor condition of $-12\,\degC$, the design heat-loss rate
is $(UA)_{\text{tot}}\,\Delta T = 317 \times 33 = 10.5\,\text{kW}$, which
sizes the heating plant.

\subsection{Seasonal energy by the degree-day method}

The heating energy over a season is the loss coefficient times the
time-integral of the indoor-outdoor temperature difference, and that
integral is the \engterm{heating degree-days}. A degree-day accrues one
degree-kelvin of deficit held for one day; summing the daily deficits of
outdoor temperature below a base temperature over the season gives the
degree-day total. The base is the outdoor temperature at which the
house needs no heat, the indoor set point less the rise that internal
gains and solar gain supply, here taken as a base of $15.5\,\degC$ (a
$5.5\,\si{\kelvin}$ free-gain credit below the $21\,\degC$ set point). The
site has $H = 2800$ heating degree-days (kelvin-days) per season against
that base. The seasonal heating energy is
\[
E = (UA)_{\text{tot}} \times H \times \frac{86400\,\text{s}}{\text{day}}
= 317 \times 2800 \times 86400\,\text{J} = 7.67 \times 10^{10}\,\text{J},
\]
which converted to the working unit is $E = 7.67\times10^{10}/(3.6\times
10^9) = 21.3\,\text{MWh}$ of delivered heat per season. Charged to the
components in proportion to their conductances (with the bridges and
infiltration distributed), the walls and windows carry the largest blocks,
the roof the smallest, the breakdown plotted in the before-upgrade bars of
Figure~\ref{fig:vol04ch05:degree-day-energy}.

% Figure: annual heating energy by envelope component for the building
% case study, shown before and after the insulation upgrade. Stacked-style
% grouped bars in MWh per heating season, with the windows and walls the
% largest movers. Numbers are the case-study results in the text.
\begin{figure}[htbp]
\centering
\begin{tikzpicture}
\begin{axis}[
  width=12cm, height=7.5cm,
  ybar,
  bar width=10pt,
  enlarge x limits=0.18,
  xlabel={envelope component},
  ylabel={annual heating energy (MWh/season)},
  ymin=0, ymax=7,
  symbolic x coords={walls, windows, roof, floor, infiltration},
  xtick=data,
  x tick label style={font=\scriptsize, rotate=20, anchor=east},
  tick label style={font=\scriptsize},
  label style={font=\small},
  legend style={font=\scriptsize, at={(0.98,0.98)}, anchor=north east},
  legend cell align=left,
  nodes near coords,
  every node near coord/.append style={font=\tiny},
]
\addplot[draw=engred, fill=engred!25]
  coordinates {(walls,2.6) (windows,3.3) (roof,0.9) (floor,1.8) (infiltration,1.8)};
\addplot[draw=engblue, fill=engblue!25]
  coordinates {(walls,1.7) (windows,1.3) (roof,0.9) (floor,1.8) (infiltration,1.2)};
\legend{before upgrade, after upgrade}
\end{axis}
\end{tikzpicture}
\caption[Annual heating energy by component, before and after upgrade]{Annual
heating energy charged to each envelope component for the building case
study, computed from the component conductances and the site degree-days.
The windows and the walls carry the largest pre-upgrade loads and move the
most after the glazing and wall-insulation upgrade; the roof is already
well insulated and barely changes, and the slab floor is left untouched in
this scope. The upgrade trims the season total from about $10.4$ to
$6.9\,$MWh, a saving the text turns into a simple payback against the
retrofit cost.}
\label{fig:vol04ch05:degree-day-energy}
\end{figure}


\subsection{The retrofit and its payback}

The breakdown points the money at the windows and the walls, so the
retrofit replaces the windows with a triple-glazed low-emissivity unit at
$U = 1.1$ and adds $100\,\si{\milli\meter}$ of external wall insulation
($k = 0.035\,\text{W}/(\text{m}\cdot\si{\kelvin})$) that also covers the
framing bridges. The external layer adds $R = 0.100/0.035 = 2.86\,\si{
\kelvin}\cdot\text{m}^2/\text{W}$ in series with the existing wall, taking
the wall from $U = 0.45$ ($R = 2.22$) to $U = 1/(2.22 + 2.86) = 0.197\,
\text{W}/(\text{m}^2\cdot\si{\kelvin})$, and because the new layer is
continuous over the studs, the repeating-bridge uplift collapses from
$15\,\%$ to about $3\,\%$. The post-retrofit conductances are

\begin{center}
\begin{tabular}{lccc}
\toprule
component & $U_{\text{new}}$ (W/m$^2$K) & $UA$ (W/K) & with bridges (W/K) \\
\midrule
walls & $0.197$ & $35.5$ & $36.5$ \\
windows & $1.1$ & $35.2$ & $35.2$ \\
roof/ceiling & $0.20$ & $19.0$ & $19.0$ \\
ground floor & $0.35$ & $33.3$ & $33.3$ \\
\bottomrule
\end{tabular}
\end{center}

The linear-bridge allowance falls to $10\,\text{W}/\si{\kelvin}$ as the
external insulation wraps the perimeter junctions, and air-sealing during
the work cuts infiltration from $0.6$ to $0.45$ air-changes per hour, so
$(UA)_{\text{inf}}$ drops to $48.2\,\text{W}/\si{\kelvin}$. The new total
is
\[
(UA)_{\text{tot}}' = 36.5 + 35.2 + 19.0 + 33.3 + 10.0 + 48.2 =
182.2\,\text{W}/\si{\kelvin},
\]
a $43\,\%$ cut from $317$. The post-retrofit seasonal energy scales with
the loss coefficient, $E' = 21.3 \times 182.2/317 = 12.3\,\text{MWh}$, the
after-upgrade bars of Figure~\ref{fig:vol04ch05:degree-day-energy}, a
saving of $9.0\,\text{MWh}$ per season. The two bar groups in the figure
are drawn for a focused scope that upgrades windows, walls, and
infiltration and leaves the roof and floor, so the figure's component
totals differ from the full-energy numbers carried here; both tell the
same story, that the windows and walls move and the roof barely does.

At a delivered-heat cost of \$0.12 per kilowatt-hour and a boiler seasonal
efficiency of $0.90$, the saving is worth $9000/0.90 \times 0.12 =
\$1200$ per season in fuel. Against a retrofit cost of \$18{,}000 for the
windows and the external wall insulation, the simple payback is
\[
\text{payback} = \frac{\$18{,}000}{\$1200/\text{season}} = 15\,\text{
seasons}.
\]
The fabric measures pay back over the slow side of a building-improvement
horizon, which is the honest verdict: the comfort gain and the carbon
reduction are immediate, the financial payback is long, and the window
replacement carries most of the cost while the wall insulation carries
most of the energy saving. A homeowner working to a budget would stage the
work, doing the air-sealing and the wall insulation first for the larger
energy return per dollar and deferring the windows, and the degree-day
calculation is what lets that comparison be made on numbers rather than on
intuition.

\begin{principle}[A steady conductance becomes annual energy through the
degree-day integral]
A building's heat loss is set at any instant by its loss coefficient
$(UA)_{\text{tot}}$, the sum of every fabric path, its thermal bridges,
and the infiltration term. The seasonal energy is that coefficient times
the heating degree-days, the time-integral of the indoor-outdoor deficit
above the building's free-gain base temperature. The design discipline is
to rank the components by their share of $(UA)_{\text{tot}}$, to correct
the clear-field fabric sum upward for the bridges that short-circuit the
insulation, and to weigh any upgrade by the energy it saves against its
cost over the building's life, not by its U-value alone.
\end{principle}

\section{A third end-to-end case study: thermal control of a small
spacecraft}\pathtag{standard}

The processor case rejected heat into moving air, and the house leaked it
through layered walls. A spacecraft has neither: in orbit there is no
fluid to convect into and no ground to conduct into, so radiation is the
only path out and the engineering reduces to a fourth-power balance
between absorbed loads and a radiator. We take a small satellite bus and
carry it from the electronics and solar loads through the radiator sizing,
the multilayer insulation that isolates the bus from the environment, the
hot-case and cold-case equilibrium temperatures, and the survival-heater
budget, ending in a verdict the thermal engineer would sign. Every tool of
the radiation section appears: the Stefan-Boltzmann law sizes the
radiator, the small-body-in-enclosure formula sets the bus-to-space
exchange, the surface-resistance algebra fixes the multilayer stack, and
the combined-load balance closes the temperature.

\subsection{Specification and the two design cases}

The avionics box dissipates between $80\,\text{W}$ in a low-power safe
mode and $140\,\text{W}$ at full operation. The electronics must stay
between a lower survival limit of $-10\,\degC$ and an upper operating
limit of $+40\,\degC$. The bus is wrapped in multilayer insulation on all
faces except a dedicated radiator panel that looks at deep space (taken as
a $4\,\si{\kelvin}$ sink, so its $T^4$ is negligible against any panel
temperature). Two attitudes bracket the mission. In the \emph{hot case}
the radiator runs warm at the upper limit, the electronics dissipate the
full $140\,\text{W}$, and the panel absorbs a solar and albedo leak of
$90\,\text{W}$ from grazing sunlight and from the planet below. In the
\emph{cold case} the electronics idle at $80\,\text{W}$, the radiator
sees no sun, and the bus tends to fall below its lower limit unless
heaters make up the difference (Figure~\ref{fig:vol04ch05:spacecraft-balance}).
The radiator is sized for the hot case and the heaters for the cold case,
the two conditions every flight attitude falls between.

% Figure: the spacecraft thermal-balance heat budget, drawn as two stacked
% bars (cold case and hot case) of the radiator heat budget. Each bar
% stacks the loads the radiator must reject; a dashed line marks the
% radiator's rejection capability at the survival temperature. The cold
% case shows where heater power must fill the deficit.
\begin{figure}[htbp]
\centering
\begin{tikzpicture}
\begin{axis}[
  width=11cm, height=8cm,
  ybar stacked,
  bar width=42pt,
  ymin=0, ymax=320,
  ylabel={power (W)},
  symbolic x coords={cold case, hot case},
  xtick=data,
  enlarge x limits=0.5,
  axis lines=left,
  legend pos=outer north east,
  legend cell align=left,
  tick label style={font=\scriptsize},
  label style={font=\small},
  legend style={font=\scriptsize},
  clip=false,
]
% cold case: low electronics dissipation, no solar absorbed on radiator,
% heater fills to survival. hot case: full electronics + solar leak.
\addplot[fill=engblue!70, draw=engblue] coordinates {(cold case,80) (hot case,140)};
\addlegendentry{electronics dissipation}
\addplot[fill=engamber!70, draw=engamber] coordinates {(cold case,0) (hot case,90)};
\addlegendentry{absorbed solar/albedo}
\addplot[fill=englime!60!black, draw=englime!50!black] coordinates {(cold case,70) (hot case,0)};
\addlegendentry{heater power}
% the radiator rejection capability at the control band
\draw[engred, very thick, dashed] (axis cs:cold case,150) -- (axis cs:hot case,150);
\node[font=\scriptsize, engred, anchor=south] at (axis cs:hot case,151) {radiator reject $\approx 230$ at hot limit};
\node[font=\scriptsize, engred, anchor=south] at (axis cs:cold case,151) {balance line $150$};
\end{axis}
\end{tikzpicture}
\caption[Spacecraft thermal balance: cold and hot cases]{The radiator
heat budget for the spacecraft thermal-control case study, drawn for the
two design extremes. In the cold case the electronics run at reduced
power and the radiator sees no sun, so survival heaters supply the
deficit that keeps the avionics above their lower limit. In the hot case
the electronics run at full power and the radiator absorbs a solar and
albedo leak, and the radiator is sized so that its fourth-power rejection
at the upper temperature limit clears the summed load. The balance line
is the steady load the equilibrium calculation matches; the radiator is
sized for the hot case and the heater for the cold case, the two
conditions that bracket every flight attitude.}
\label{fig:vol04ch05:spacecraft-balance}
\end{figure}


\subsection{Sizing the radiator for the hot case}

A radiator of area $A_r$, emissivity $\epsilon_r = 0.85$ (a typical
white-paint or silvered-Teflon finish), and temperature $T_r$ rejects to
space, by the small-body-in-enclosure formula with a $4\,\si{\kelvin}$
sink,
\[
\dot{Q}_{\text{rej}} = \epsilon_r A_r \sigma (T_r^4 - T_{\text{space}}^4)
\approx \epsilon_r A_r \sigma T_r^4.
\]
In the hot case the radiator must reject the full electronics load plus
the absorbed environmental leak, $\dot{Q}_{\text{rej}} = 140 + 90 =
230\,\text{W}$, while running no hotter than the upper limit, $T_r =
40\,\degC = 313\,\si{\kelvin}$. The required area is
\[
A_r = \frac{\dot{Q}_{\text{rej}}}{\epsilon_r \sigma T_r^4} =
\frac{230}{0.85 \times 5.67 \times 10^{-8} \times 313^4}.
\]
With $313^4 = 9.60 \times 10^9$, the denominator is $0.85 \times 5.67
\times 10^{-8} \times 9.60 \times 10^9 = 463\,\text{W}/\text{m}^2$, so
$A_r = 230/463 = 0.497\,\text{m}^2$. We carry a $0.55\,\text{m}^2$ panel,
about ten percent larger, to hold margin against finish degradation: the
silvered surface's solar absorptance climbs over years of ultraviolet and
atomic-oxygen exposure, which raises the absorbed leak and forces the
panel warmer for the same duty. The absorbed leak itself follows from the
panel's solar absorptance $\alpha_s$ and the incident flux. Taking
$\alpha_s = 0.20$ for the white finish and a worst-case grazing solar plus
albedo input of about $450\,\text{W}/\text{m}^2$ on the panel, the
absorbed solar is $\alpha_s \times 450 \times 0.55 = 49.5\,\text{W}$, and
a further $40\,\text{W}$ of infrared from the warm planet below brings the
leak near the $90\,\text{W}$ carried, which is the value the budget used.

\subsection{The multilayer insulation around the bus}

The bus faces that are not the radiator must be near-perfect insulators,
or the cold-case heater budget explodes. Multilayer insulation does this
by stacking thin reflective shields that each add a radiative resistance
in series, the principle of the radiation-shield worked example carried to
many layers (Figure~\ref{fig:vol04ch05:mli-layers}). An ideal stack of
$N$ floating shields cuts the bare-gap flux by $1/(N+1)$, but a flight
blanket is degraded by the spacer netting that conducts across the shields
and by seams and penetrations, so the design quotes a measured
\engterm{effective emissivity} $\epsilon^* \approx 0.02$ to $0.05$ for the
whole blanket rather than the ideal layer count.

% Figure: the effective radiative conductance of a multilayer-insulation
% stack against the number of floating shields, showing the 1/(N+1)
% reduction. The ideal-shield curve falls as 1/(N+1); real MLI is worse
% because supports and edges short-circuit the shields, drawn as a second
% curve that flattens.
\begin{figure}[htbp]
\centering
\begin{tikzpicture}
\begin{axis}[
  width=11cm, height=7.5cm,
  xlabel={number of floating shields $N$},
  ylabel={heat flux relative to bare gap},
  xmin=0, xmax=20,
  ymin=0, ymax=1.05,
  axis lines=left,
  legend pos=north east,
  legend cell align=left,
  tick label style={font=\scriptsize},
  label style={font=\small},
  legend style={font=\scriptsize},
  domain=0:20,
  samples=80,
]
% ideal: flux ratio = 1/(N+1)
\addplot[engblue, very thick] {1/(x+1)};
\addlegendentry{ideal shields, $1/(N{+}1)$}
% real MLI: degraded by conduction through spacers and edge leaks,
% modelled as a floor that the stack cannot beat
\addplot[engred, very thick, dashed] {1/(x+1) + 0.02};
\addlegendentry{real MLI with parasitic paths}
% mark a working stack
\addplot[only marks, mark=*, englime!50!black] coordinates {(10,0.111)};
\node[font=\scriptsize, engblack, anchor=west] at (axis cs:10.5,0.16)
  {ten-shield stack: $\approx 9\%$ of bare};
\end{axis}
\end{tikzpicture}
\caption[Multilayer-insulation effective conductance against shield
count]{The radiative heat flux across a stack of $N$ floating
low-emissivity shields, relative to the flux across the same bare gap.
Ideal shields each add a radiative resistance in series, so $N$ shields
cut the flux by the factor $1/(N+1)$, the curve a ten-shield stack rides
to about nine percent of the bare value. Real multilayer insulation
falls short of the ideal because the spacer netting conducts heat across
the shields and the seams and penetrations leak, so the stack reaches a
floor that more layers cannot beat, drawn as the dashed curve. The
spacecraft case carries a flight-representative effective emissivity for
the wrapped stack rather than the ideal count, which is why the design
quotes a measured effective conductance instead of multiplying the
single-shield result by the layer number.}
\label{fig:vol04ch05:mli-layers}
\end{figure}

 We take $\epsilon^* =
0.03$ over a wrapped bus area of $A_b = 2.0\,\text{m}^2$ at a bus skin
temperature near $T_b = 290\,\si{\kelvin}$. The parasitic heat leaking out
through the blanket in the cold case is
\[
\dot{Q}_{\text{MLI}} = \epsilon^* A_b \sigma (T_b^4 - T_{\text{space}}^4)
\approx 0.03 \times 2.0 \times 5.67 \times 10^{-8} \times 290^4.
\]
With $290^4 = 7.07 \times 10^9$, this is $0.03 \times 2.0 \times 5.67
\times 10^{-8} \times 7.07 \times 10^9 = 24.1\,\text{W}$. A bare bus with
$\epsilon = 0.85$ over the same area would leak $0.85/0.03 = 28$ times as
much, near $680\,\text{W}$, which no practical heater could replace. The
blanket turns an impossible heater budget into a manageable one, and the
remaining leak is the term the cold-case heater must cover.

\subsection{Hot-case and cold-case equilibrium}

The bus temperature settles where the loads balance the radiator
rejection. In the hot case the radiator was sized to sit at $313\,
\si{\kelvin}$ rejecting $230\,\text{W}$, so the equilibrium is the upper
limit by construction, with the $0.55\,\text{m}^2$ panel running slightly
cooler than the limit at the design loads and reaching the limit only as
the finish degrades. The cold case is the binding one. With the
electronics at $80\,\text{W}$, no absorbed sun, and the radiator now
running cold, the panel would settle where $\epsilon_r A_r \sigma T_r^4 =
80 + \dot{Q}_{\text{MLI}}$ less any heater, and without a heater that
balance lands far below the survival limit. Setting the radiator to the
lower limit $T_r = -10\,\degC = 263\,\si{\kelvin}$, its rejection there is
\[
\dot{Q}_{\text{rej,cold}} = 0.85 \times 0.55 \times 5.67 \times 10^{-8}
\times 263^4 = 0.85 \times 0.55 \times 5.67 \times 10^{-8} \times 4.79
\times 10^9 = 127\,\text{W}.
\]
The radiator that was right-sized to dump $230\,\text{W}$ when warm still
insists on rejecting $127\,\text{W}$ when cold, because the fourth-power
law does not switch off. Against an electronics supply of only $80\,\text{W}$
plus a $24\,\text{W}$ blanket leak the bus would over-reject by $127 -
(80 - 24) = 71\,\text{W}$ and freeze, so the cold-case deficit the heater
must supply is about $70\,\text{W}$.

\subsection{Heater budget and verdict}

The survival-heater budget is the cold-case deficit, near $70\,\text{W}$,
which the thermal design supplies with a thermostatically switched
resistance heater on the radiator or the avionics plate. The verdict is
that a $0.55\,\text{m}^2$ radiator and a $70\,\text{W}$ heater hold the
electronics inside their $-10$ to $+40\,\degC$ band across both extremes.
The budget also shows the structural tension of every passive radiator: it
is sized for the hot case and therefore oversized for the cold case, and
the heater pays the standing penalty of that oversize for the life of the
mission. A designer who wanted to cut the heater power would fit a
louvre or a variable-emissivity coating that throttles the radiator's
effective emissivity down in the cold case, trading mechanism complexity
for heater energy, the same hot-case-versus-cold-case compromise that sets
every spacecraft thermal design. The general treatment of radiative
exchange and the surface-resistance network this case rests on is
developed at length in the standard heat-transfer references
\cite{text:incropera2007}.

\begin{principle}[In a vacuum, radiation is the only budget and it is set
by the worst two cases]
A body that can neither convect nor conduct to its surroundings reaches
steady state where its absorbed loads equal its radiated output, a
balance governed entirely by the fourth-power law and the surface
emissivities. The design is fixed by two bracketing conditions: a hot
case of maximum dissipation and maximum absorbed environment, which sizes
the radiator at the upper temperature limit, and a cold case of minimum
dissipation and no absorbed sun, which sizes the survival heater at the
lower limit. Multilayer insulation isolates every face that is not a
radiator, quoted by a measured effective emissivity rather than an ideal
layer count, and the standing heater penalty is the price of a radiator
sized for the hot case but still rejecting in the cold one.
\end{principle}

\section{Failure: thermal runaway and where it actually happens}\pathtag{core}

\engterm{Thermal runaway} is a positive-feedback instability in which
a small temperature rise causes the system to generate heat faster
than it can be removed, leading to an accelerating temperature
increase. The mechanism appears across engineering domains, with the
common form: a heat-generation rate that rises faster with temperature
than the heat-removal rate.

The mathematical condition for thermal runaway in a system with
internal generation $\dot{q}_g(T)$ and removal $\dot{q}_r(T)$ is
\[
\dot{q}_g(T) > \dot{q}_r(T) \quad \text{and} \quad
\frac{d\dot{q}_g}{dT} > \frac{d\dot{q}_r}{dT}.
\]
The first condition makes the temperature rise; the second makes the
rise self-accelerating. Engineering design seeks $\dot{q}_g(T) \leq
\dot{q}_r(T)$ at all operating conditions or, where the first
condition is violated only transiently, ensures the second is not.

% Figure: thermal-runaway heat-balance diagram. Heat-generation curve rises
% exponentially (Arrhenius) with temperature; heat-removal line is linear in
% (T - T_coolant). Intersections are operating points: a stable low point, an
% unstable middle point (tangency / ignition), and the runaway branch above.
% Two removal lines shown for different cooling capacities.
\begin{figure}[htbp]
\centering
\begin{tikzpicture}
\begin{axis}[
  width=12cm, height=7.5cm,
  xlabel={temperature $T$},
  ylabel={heat rate $\dot{q}$},
  xmin=0, xmax=10, ymin=0, ymax=10,
  axis lines=left,
  legend pos=north west,
  legend cell align=left,
  tick label style={font=\scriptsize},
  label style={font=\small},
  legend style={font=\scriptsize},
  xtick=\empty, ytick=\empty,
  samples=120, domain=0:10,
  clip=false,
]
% generation: Arrhenius-like sigmoid (heat release rises then saturates as
% reactant depletes). Use a logistic-shaped curve so it stays bounded.
\addplot[engred, very thick] {9.2/(1+exp(-1.5*(x-5.2)))};
\addlegendentry{generation $\dot{q}_g(T)$}
% removal line 1: adequate cooling, steep slope -> single low stable point
\addplot[engblue, very thick] {1.15*(x-1.0)};
\addlegendentry{removal $\dot{q}_r(T)$, high cooling}
% removal line 2: marginal cooling, shallow slope -> three intersections
\addplot[engamber, very thick, dashed] {0.62*(x-0.8)};
\addlegendentry{removal $\dot{q}_r(T)$, low cooling}
% mark intersections on the marginal-cooling case (approximate, illustrative)
\fill[engamber] (axis cs:1.55,0.46) circle (2.0pt);
\node[font=\scriptsize, engamber, anchor=north west] at (axis cs:1.55,0.5) {stable};
\fill[engamber] (axis cs:4.5,2.29) circle (2.0pt);
\node[font=\scriptsize, engamber, anchor=south east] at (axis cs:4.5,2.29) {unstable (ignition)};
\fill[engamber] (axis cs:8.4,4.71) circle (2.0pt);
\node[font=\scriptsize, engamber, anchor=south east] at (axis cs:8.4,4.71) {runaway branch};
\end{axis}
\end{tikzpicture}
\caption[Thermal-runaway heat-balance diagram]{The heat-balance diagram
that diagnoses thermal runaway. The generation curve $\dot{q}_g(T)$ rises
steeply with temperature (Arrhenius kinetics in a reactor, leakage
current in a chip, side reactions in a cell) and saturates as the driving
reactant depletes. The removal rate $\dot{q}_r(T)$ is close to linear in
the temperature above coolant. Where the two cross is a steady operating
point: it is stable if the removal curve cuts the generation curve from
below ($d\dot{q}_r/dT > d\dot{q}_g/dT$) and unstable otherwise. With ample
cooling the steep removal line meets the generation curve at a single low
stable point. With marginal cooling the shallow line produces a low
stable point, an unstable ignition point, and a high runaway branch; a
disturbance past the ignition point drives the system up to the runaway
branch. The design rule is to keep the removal line steep enough that no
unstable intersection sits inside the operating envelope.}
\label{fig:vol04ch05:runaway-balance}
\end{figure}


Figure~\ref{fig:vol04ch05:runaway-balance} is the diagram that decides
the question. The generation curve climbs steeply with temperature; the
removal rate is close to linear. Where they cross is a steady operating
point, stable if the removal curve cuts from below and unstable
otherwise. Ample cooling gives a single low stable point; marginal
cooling admits an unstable ignition point past which the system climbs
to a high runaway branch.

\begin{mastery}
The Semenov analysis makes the runaway condition quantitative for a
well-stirred exothermic reactor. Heat generation follows Arrhenius
kinetics, $\dot{q}_g = V \Delta H_r A_0 \, c \, e^{-E_a/(RT)}$, and heat
removal is $\dot{q}_r = U A_s (T - T_c)$ to a coolant at $T_c$. Plotting
both against $T$, a stable steady state needs the removal line to be
steeper than the generation curve at the intersection,
$d\dot{q}_r/dT > d\dot{q}_g/dT$. Setting the two rates and their slopes
equal at the tangency gives the critical condition. Expanding the
Arrhenius term about the coolant temperature and writing
$\theta = E_a(T - T_c)/(RT_c^2)$ collapses the result to the Semenov
number $\psi = \dot{q}_g(T_c) E_a/(U A_s R T_c^2)$, with ignition at
$\psi = e^{-1} \approx 0.368$. The lesson the dimensionless form carries
is that doubling the cooling coefficient $U A_s$ buys the same margin as
halving the reactant concentration or the batch size, which is why
inherent-safety arguments favour smaller inventories over larger chillers.
\end{mastery}

% Figure: three heat-removal lines against one Arrhenius generation curve,
% showing the tangency (ignition) case between an amply cooled subcritical
% case and a supercritical case with no low intersection. Generation is an
% exponential proxy for the Arrhenius rate; removal lines are straight.
\begin{figure}[htbp]
\centering
\begin{tikzpicture}
\begin{axis}[
  width=12cm, height=7.5cm,
  xlabel={temperature $T$ (arb. units above coolant)},
  ylabel={heat rate (arb. units)},
  xmin=0, xmax=10, ymin=0, ymax=10,
  axis lines=left,
  tick label style={font=\scriptsize},
  label style={font=\small},
  legend style={font=\scriptsize, at={(0.02,0.98)}, anchor=north west},
  legend cell align=left,
  legend entries={generation (Arrhenius),ample cooling,critical (tangent),marginal cooling},
  clip=false,
]
% generation: exponential proxy, clamped by domain
\addplot[engred, very thick, domain=0:9.2, samples=60] {0.35*exp(0.36*x)};
% ample cooling: steep removal line, high slope, cuts generation once low
\addplot[engblue, thick, domain=0:10, samples=2] {1.55*x};
% critical: tangent line -- slope chosen to graze the generation curve
\addplot[engamber, thick, domain=0:10, samples=2] {0.9*x + 0.55};
% marginal: shallow line, only high intersection
\addplot[enggray, thick, dashed, domain=0:10, samples=2] {0.55*x + 0.9};
% mark the stable low point on the ample-cooling case (approx)
\addplot[only marks, mark=*, engblack, mark size=1.6pt] coordinates {(0.24,0.37)};
\node[font=\scriptsize, engblack, anchor=west] at (axis cs:0.5,0.9) {stable};
\node[font=\scriptsize, engamber, anchor=south west] at (axis cs:4.6,5.1) {ignition};
\end{axis}
\end{tikzpicture}
\caption[Ignition as a tangency of removal to generation]{The thermal-runaway
condition read off the heat-balance diagram. The generation rate climbs
steeply with temperature (an Arrhenius or leakage-current curve); the
removal rate is a straight line whose slope is the cooling conductance.
Ample cooling (steep line) cuts the generation curve once at a low, stable
operating point. As the cooling weakens the line tilts down until it just
grazes the generation curve: that tangency is the ignition condition, the
last configuration with any low steady state. Weaker still (dashed) the
line clears the low crossing entirely and the only balance lies far up the
curve, so the system has no cool operating point and climbs to the runaway
branch. The design margin against runaway is the gap between the operating
cooling line and this tangent.}
\label{fig:vol04ch05:runaway-stability}
\end{figure}


Figure~\ref{fig:vol04ch05:runaway-stability} sharpens the balance diagram
into the criterion that decides whether a design has any safe operating
point at all. As the cooling weakens, the straight removal line tilts down
until it just grazes the steep generation curve, and that tangency is the
ignition condition: the last configuration with a low, cool steady state.
Any weaker cooling clears the low crossing entirely, and the only balance
left lies far up the generation curve, on the runaway branch. The design
margin against runaway is the angular gap between the operating cooling line
and this tangent, and it is worth putting a number to it.

\engterm{An ignition margin for a stored-material drum}. We take a drum of
self-heating material whose generation rate we model over the range that
matters as $\dot{q}_g(T) = \dot{q}_0\,e^{(T - T_a)/T^*}$ with $\dot{q}_0 =
2.0\,\text{W}$ at the ambient $T_a = 300\,\si{\kelvin}$ and a characteristic
temperature $T^* = 12\,\si{\kelvin}$ over which the rate rises by a factor
$e$ (a proxy for the Arrhenius slope near ambient). The drum loses heat to
its surroundings as $\dot{q}_r(T) = \Lambda(T - T_a)$ with a cooling
conductance $\Lambda$ we are choosing. Ignition is the tangency where the
two rates and their slopes match at once,
\[
\dot{q}_0\,e^{(T-T_a)/T^*} = \Lambda(T - T_a), \qquad \frac{\dot{q}_0}{T^*}
e^{(T-T_a)/T^*} = \Lambda.
\]
Dividing the first by the second gives $T - T_a = T^* = 12\,\si{\kelvin}$ at
the tangency, so the drum sits only twelve kelvin above ambient at the
knife-edge, and back-substituting fixes the critical conductance
\[
\Lambda_c = \frac{\dot{q}_0}{T^*}e^{1} = \frac{2.0}{12}\times2.718 =
0.453\,\text{W}/\si{\kelvin}.
\]
A drum with cooling conductance above $0.453\,\text{W}/\si{\kelvin}$ has a
stable low operating point; one below it has none and runs away from
ambient. This is the dimensionless Semenov result of the mastery box in
concrete numbers: the critical group $\dot{q}_0 T^{*-1}\Lambda^{-1}$ crosses
$e^{-1}$ exactly at $\Lambda_c$. The lesson the tangency carries is that the
approach to ignition gives almost no warning, because the drum sits a mere
twelve kelvin warm right up to the edge and then leaves; a temperature
alarm set to trip at a modest rise buys little time, which is why bulk
storage of self-heating materials is limited by inventory and pile size,
the levers that lower $\dot{q}_0$, rather than by monitoring alone.

\begin{mastery}
The leakage-current runaway of the electronics subsection is the same
tangency, and casting it in the drum's algebra shows the semiconductor
version of the critical cooling conductance. A part's leakage power near an
operating point rises as $\dot{q}_g(T_j) = P_0\,2^{(T_j - T_0)/T_d}$, where
$T_d$ is the temperature interval over which the leakage doubles, eight to
eleven kelvin for a modern process. Writing it as an exponential,
$\dot{q}_g = P_0\exp[(T_j - T_0)\ln 2/T_d]$, its slope is $\dot{q}_g\ln
2/T_d$. The package removes heat as $\dot{q}_r = (T_j - T_a)/R_{ja}$, a line
of slope $1/R_{ja}$. Ignition is where the slopes match while the rates
match, so as in the drum the operating excess at tangency is the
characteristic interval $T_d/\ln 2$, and the critical junction-to-ambient
resistance is
\[
R_{ja,c} = \frac{T_d}{\ln 2}\,\frac{1}{\dot{q}_g^*},
\]
with $\dot{q}_g^*$ the leakage power at the tangency. A part dissipating a
tenth of a watt of leakage at the edge with $T_d = 10\,\si{\kelvin}$ tips
over once $R_{ja}$ exceeds $(10/0.693)/0.1 = 144\,\si{\kelvin}/\text{W}$,
which a bare package in still air can easily reach, while the same part in a
good heat sink at $R_{ja} = 20\,\si{\kelvin}/\text{W}$ sits far on the
stable side. The identical device is therefore stable or unstable depending
on its cooling, at the same power, the quantitative form of the
subsection's claim that a processor runs for years and then fails in minutes
once its fan clogs. Undervolting lowers $P_0$ and flattens the whole
generation curve, raising $R_{ja,c}$ and buying margin from the software
side, the complement to the fan and the sink on the hardware side.
\end{mastery}

\subsection{Lithium-ion batteries}

A lithium-ion battery generates internal heat by Ohmic resistance
during charge and discharge, and additionally by exothermic
electrochemical side reactions whose rates rise with temperature. The
exothermic side reactions begin near $80\,\degC$ for typical lithium-
ion chemistries, with rates that double every $10\,\si{\kelvin}$
(Arrhenius behaviour). Above $130\,\degC$, the solid electrolyte
interface (SEI) layer on the anode begins to decompose, releasing
additional heat and exposing fresh lithium to the electrolyte.
Above $200\,\degC$, the electrolyte itself decomposes and the
cathode releases oxygen, which sustains combustion of the
electrolyte and binder.

The thermal-runaway pattern: a thermal stimulus (overcharging,
external heating, mechanical damage that creates an internal short)
raises a cell's temperature to the side-reaction onset; the side
reactions add to the temperature rise; SEI decomposition adds further;
cathode oxygen release ignites the electrolyte; the cell vents flames
and propagates heat to neighbouring cells. A single cell's runaway
can propagate through a pack, with the entire pack consumed in
seconds to minutes.

Engineering controls for lithium-ion thermal runaway include cell-
level: separators that shut down at $\sim 130\,\degC$, pressure-
relief vents, current-interrupt devices; pack-level: thermal
isolation between cells to prevent propagation, active cooling
systems, battery-management systems that monitor each cell's
temperature and voltage; system-level: fire-resistant enclosures,
deluge systems, geographic separation between battery packs in
storage facilities.

\subsection{Reactive chemicals and process reactors}

A chemical reaction with $\Delta H_r < 0$ (exothermic) in a poorly-
cooled vessel can run away when the reaction rate (Arrhenius:
exponential in $T$) outstrips the cooling rate (typically linear in
$T - T_{\text{coolant}}$). The Semenov analysis (1928) gives the
critical condition for thermal runaway in a stirred tank:
\[
\frac{E_a}{RT_s^2} > \frac{1}{\Delta T_{\text{ad}}},
\]
where $E_a$ is the activation energy, $T_s$ the steady-state
temperature, and $\Delta T_{\text{ad}}$ the adiabatic temperature rise.

Industrial chemical-process disasters associated with thermal runaway
include the Seveso dioxin release (1976), the T2 Laboratories
explosion (2007), and the Bhopal MIC release (1984), among others.
The Bhopal case is in the registry as
\cite{acc:icmr-bhopal-2004}: a tank of methyl isocyanate
underwent runaway when water entered and triggered an exothermic
reaction; the cooling system that should have removed the reaction
heat had been switched off, and the runaway proceeded to a
catastrophic vapour release. Chapter~7 of this volume develops the
mass-transfer side of the same case; the heat-transfer side is the
inability of the cooling system to keep up with the reaction's
exothermic generation rate.

The engineering disciplines for reactor thermal-runaway prevention
are inherent safety (limit the inventory of reactive material and
its concentration so that a runaway is bounded), reaction calorimetry
(measure the heat generation rate as a function of temperature and
composition to design cooling capacity), and redundant cooling
(multiple independent paths to remove heat, so the loss of one is
not catastrophic).

\subsection{Nuclear reactors}

A nuclear reactor generates heat by fission of fuel nuclei. The
heat-generation rate at any instant depends on the neutron flux and
the fuel composition; the heat-removal rate depends on the coolant
flow, temperature, and phase. Several feedback mechanisms relate
the two: the fuel's temperature affects the resonance absorption of
neutrons (Doppler broadening, typically a negative feedback that
stabilises the reactor), the moderator's temperature and density
affect the neutron moderation (often a negative feedback in water-
moderated reactors; can be positive in graphite-moderated reactors
operating in unusual regimes), and the fuel's thermal expansion
affects the geometry.

A nuclear thermal runaway requires the positive feedback to dominate
the negative feedback. In a stable reactor, the feedbacks are
designed to be net-negative across the operating envelope. The
Chernobyl accident in 1986 is in the registry as
\cite{acc:iaea-insag1-1986, acc:iaea-insag7-1992}: a positive void
coefficient (water voiding in a graphite-moderated reactor reduces
neutron absorption faster than it reduces moderation, increasing
reactivity) combined with a control-rod design that briefly
increased reactivity on insertion caused a runaway power excursion
that destroyed the reactor.

The Three Mile Island accident in 1979 (registry:
\cite{acc:kemeny-tmi-1979}) is the heat-transfer failure mode rather
than a thermal-runaway one: a loss of coolant inventory through a
stuck-open relief valve uncovered the core, which then overheated
from decay heat (the residual radioactivity of fission products,
which decays from initial values around $7\,\%$ of full power immediately
after shutdown to $\sim 1\,\%$ at one hour and $\sim 0.5\,\%$ at one day).
Decay heat is not a thermal runaway (it does not accelerate with
temperature; it follows a known time profile) but it is a heat-
generation source that persists for days after shutdown and demands
its own cooling system independent of the main coolant loop.

\subsection{Electronic components}

Solid-state electronics fail thermally when the junction temperature
exceeds the device's rated maximum. The failure mode for individual
devices is electromigration acceleration, oxide breakdown, or
catastrophic destruction. For integrated circuits, the thermal-
runaway concern is at the package level: an overheated chip drops
in performance, but a chip that exceeds $\sim 125\,\degC$ at the
junction begins to draw additional leakage current (which scales
exponentially with temperature), generating more heat, raising the
temperature further, and so on. The runaway is bounded by the
device's thermal-protection circuits (typically a thermal shutdown
at $\sim 150\,\degC$); without protection, the device destroys
itself within seconds of crossing the threshold.

The leakage-current runaway is worth casting in the chapter's own
heat-balance form, because it shows the positive-feedback condition met by
a device that has no chemistry and no phase change, only a temperature-
dependent generation term. A transistor's static leakage current rises
roughly exponentially with junction temperature, doubling every $8$ to
$11\,\si{\kelvin}$ for the sub-threshold leakage of a modern process, so
the leakage power dissipated at fixed supply voltage is $\dot{q}_g(T_j) =
P_0\,2^{(T_j - T_0)/10}$ over the range that matters, a curve that climbs
steeply. The removal side is the package conducting to ambient through the
junction-to-ambient resistance of the case study, $\dot{q}_r(T_j) = (T_j -
T_a)/R_{ja}$, a straight line. The two plotted together are exactly the
generation-and-removal diagram of the failure section: a well-cooled part
with a small $R_{ja}$ gives a steep removal line that cuts the leakage
curve once, at a low stable junction temperature; a starved part with a
large $R_{ja}$ gives a shallow line that either grazes the leakage curve
at a tangency or misses it entirely, leaving no stable operating point
below the destruction temperature. The runaway condition
$d\dot{q}_g/dT_j > d\dot{q}_r/dT_j$ becomes, at the operating point,
$\ln(2)/10\cdot\dot{q}_g > 1/R_{ja}$, so the same device is stable in a
good heat sink and unstable in a poor one at the identical power, which is
why a processor that runs for years can fail in minutes once its fan
clogs or its thermal-interface layer dries out. The leakage curve also
steepens with the supply voltage, so undervolting a part flattens the
generation curve and buys thermal margin directly, the software-side
defence that complements the fan and the sink. The thermal-protection
shutdown is the engineered version of removing the generation term when the
removal term cannot keep up, the last of the three defences the failure
principle names.

\subsection{The diagnostic for thermal runaway}

The diagnostic for any thermal-runaway-prone system is the heat-
balance plot: plot $\dot{q}_g(T)$ and $\dot{q}_r(T)$ on the same
axes. The intersections of the two curves are the steady-state
operating points. A stable point has $d\dot{q}_g/dT < d\dot{q}_r/dT$
(the removal-rate curve crosses the generation-rate curve from
below); an unstable point has the opposite. A system with two stable
points and an unstable one between them can run away if perturbed
past the unstable point; the engineering design discipline is to
operate well below the unstable intersection, or to ensure no
unstable intersection exists in the operating envelope.

\begin{principle}[Thermal runaway is a positive-feedback instability]
A system is prone to thermal runaway when its heat-generation rate
rises faster with temperature than its heat-removal rate. The
condition is captured by the Semenov criterion in chemical reactors,
by feedback coefficients in nuclear reactors, by leakage-current
behaviour in semiconductors, and by side-reaction Arrhenius rates in
batteries. The engineering discipline against runaway is the
operating-point analysis on a heat-balance diagram, the design of
adequate cooling capacity at the system's worst-case operating
condition, and the provision of independent shutdown mechanisms that
remove the generation term if the removal term cannot keep up.
\end{principle}

\section{Project}\pathtag{core}

\begin{project}[Measure the U-value of a household wall over a winter
month]
The reader instruments a single exterior wall of a residential
building and measures the heat flux through it over a winter month,
extracting the wall's U-value and comparing to the value implied by
the construction.

\textbf{Track}. Household instrumentation plus analysis.
\textbf{Hazard class}: low.

\textbf{Measurements}.

\begin{itemize}[leftmargin=*]
  \item Wall heat flux: a heat-flux sensor (HFP01-style, $\sim
    \$ 500$ commercial, or a home-built calibrated insulation pad
    with thermocouples) mounted on the interior wall surface.
  \item Indoor air temperature and surface temperature: thermocouples
    or RTDs at the interior wall surface and in the indoor air a
    few centimetres away.
  \item Outdoor air temperature and surface temperature: similar
    instrumentation outside (sheltered from direct sun and rain).
  \item Time resolution: $30\,\text{min}$ averages over four weeks,
    capturing diurnal and weather variations.
\end{itemize}

\textbf{Analysis}.

\begin{itemize}[leftmargin=*]
  \item Compute the wall's average heat flux divided by the average
    indoor-to-outdoor air-temperature difference; this is the
    apparent U-value.
  \item Correct for transient effects: filter the data to retain only
    quasi-steady periods (overnight, when solar gain is absent and
    the wall has had time to settle from daytime warming).
  \item Compare the measured U-value to the manufacturer or code
    value for the wall's construction (assumed from drawings or
    visual inspection).
  \item Identify any discrepancy: thermal bridging through framing,
    moisture in insulation, air infiltration through the wall.
\end{itemize}

\textbf{Reflection ($1{,}500$ to $2{,}500$ words)}. The reader
addresses:

\begin{itemize}[leftmargin=*]
  \item How close is the measured U-value to the design value? What
    are the most likely sources of the discrepancy?
  \item What fraction of the building's total heat loss does this
    wall represent?
  \item Which envelope upgrade (more insulation, better windows, air-
    sealing) would have the largest payback?
\end{itemize}

\textbf{Deliverable}. The data log, the U-value calculation, the
comparison to design, and the reflection.

\textbf{Time}. Four weeks of passive monitoring, four to six hours of
data reduction, four to six hours of write-up.
\end{project}

\section{Exercises}

\subsection*{Calculation}\pathtag{core}

\begin{exercise}
A flat brick wall of thickness $25\,\si{\centi\meter}$ and area
$10\,\text{m}^2$ separates a $20\,\degC$ room from a $-5\,\degC$
exterior. Take $k = 0.7\,\text{W}/(\text{m}\cdot\si{\kelvin})$ for
brick. Find the steady-state heat-loss rate and the temperature at
the midplane of the wall.
\paragraph{Solution sketch.} Pure conduction, no films given. The
resistance is $R = L/(kA) = 0.25/(0.7 \times 10) = 0.0357\,\si{\kelvin}
/\text{W}$. The heat-loss rate is $\dot{Q} = \Delta T/R = 25/0.0357 =
700\,\text{W}$, equivalently $\dot{Q} = kA\,\Delta T/L = 0.7 \times 10
\times 25/0.25$. The temperature varies linearly through a single
constant-$k$ layer, so the midplane sits halfway between the faces,
$(20 + (-5))/2 = 7.5\,\degC$. The same answer comes from following the
heat flux $q = \dot{Q}/A = 70\,\text{W}/\text{m}^2$ inward from the
cold face: $T = -5 + q(L/2)/k = -5 + 70 \times 0.125/0.7 = 7.5\,\degC$.
\end{exercise}

\begin{exercise}
A composite wall has a $10\,\si{\centi\meter}$ brick outer layer,
$15\,\si{\centi\meter}$ of insulation ($k = 0.04\,\text{W}/(\text{m}
\cdot\si{\kelvin})$), and a $1.5\,\si{\centi\meter}$ gypsum-board
inner layer ($k = 0.17\,\text{W}/(\text{m}\cdot\si{\kelvin})$). Air
films give $h_o = 25$, $h_i = 8\,\text{W}/(\text{m}^2\cdot\si{
\kelvin})$. Find the U-value and the temperature at each material
interface for indoor $22\,\degC$, outdoor $-15\,\degC$.
\paragraph{Solution sketch.} Per unit area the resistances in series
are: indoor film $1/8 = 0.125$; gypsum $0.015/0.17 = 0.088$;
insulation $0.15/0.04 = 3.75$; brick $0.10/0.7 = 0.143$; outdoor film
$1/25 = 0.04$, all in $\si{\kelvin}\cdot\text{m}^2/\text{W}$. The total
is $R = 4.146$, so $U = 1/R = 0.241\,\text{W}/(\text{m}^2\cdot
\si{\kelvin})$. The heat flux is $q = U\,\Delta T = 0.241 \times 37 =
8.9\,\text{W}/\text{m}^2$. Walking inward to outward, each interface
temperature drops by $q R_i$: indoor air $22.0$; gypsum inner face
$22 - 8.9(0.125) = 20.9$; gypsum/insulation $20.9 - 8.9(0.088) = 20.1$;
insulation/brick $20.1 - 8.9(3.75) = -13.3$; brick outer face $-13.3 -
8.9(0.143) = -14.6$; outdoor air $-15.0\,\degC$. The insulation carries
almost all of the drop, as Figure~\ref{fig:vol04ch05:wall-temperature-profile}
shows. The \texttt{thermal\_resistance.py} module reproduces these
numbers.
\end{exercise}

\begin{exercise}
A steel pipe of inner diameter $40\,\si{\milli\meter}$, wall thickness
$3\,\si{\milli\meter}$, carries water at $80\,\degC$ at a velocity of
$1.5\,\text{m/s}$. The pipe is in $20\,\degC$ air with $h_o = 10\,
\text{W}/(\text{m}^2\cdot\si{\kelvin})$. Using the Dittus-Boelter
correlation for the inner convection, find the heat-loss rate per
metre.
\paragraph{Solution sketch.} For water near $80\,\degC$ take
$\rho = 972\,\text{kg/m}^3$, $\mu = 3.55 \times 10^{-4}\,\text{Pa}\cdot
\text{s}$, $k = 0.67\,\text{W}/(\text{m}\cdot\si{\kelvin})$,
$\text{Pr} = 2.2$. With $D_i = 0.040\,\text{m}$ and $u = 1.5\,\text{m/s}$,
$\text{Re}_D = \rho u D/\mu = 972 \times 1.5 \times 0.040/3.55\times
10^{-4} \approx 1.6 \times 10^5$, fully turbulent. Dittus-Boelter for
cooling, $\text{Nu} = 0.023\,\text{Re}^{0.8}\text{Pr}^{0.3} = 0.023
(1.6\times10^5)^{0.8}(2.2)^{0.3} \approx 460$, so $h_i = \text{Nu}\,k/D
= 460 \times 0.67/0.040 \approx 7700\,\text{W}/(\text{m}^2\cdot
\si{\kelvin})$. Per metre the resistances are: inner film $1/(h_i \pi
D_i) = 1/(7700 \times \pi \times 0.040) = 1.03 \times 10^{-3}$; steel
wall $\ln(0.046/0.040)/(2\pi \times 50) = 4.5 \times 10^{-4}$; outer
film $1/(h_o \pi D_o) = 1/(10 \times \pi \times 0.046) = 0.692$, all in
$\si{\kelvin}\cdot\text{m}/\text{W}$. The outer film dominates, so
$R_{\text{tot}} \approx 0.693$ and $\dot{Q}/L = (80 - 20)/0.693 \approx
87\,\text{W}/\text{m}$. The inner convection and the steel are
negligible; the design lever is the outside, which is what insulation
or a higher outer coefficient would address.
\end{exercise}

\begin{exercise}
A copper sphere of diameter $20\,\si{\milli\meter}$ at $90\,\degC$ is
dropped into $20\,\degC$ water with $h = 500\,\text{W}/(\text{m}^2
\cdot\si{\kelvin})$. Take $k_{Cu} = 400$, $\rho_{Cu} = 8960\,\text{kg/m}
^3$, $c_{Cu} = 385\,\text{J}/(\text{kg}\cdot\si{\kelvin})$. Determine
whether the lumped-capacitance approximation applies, and find the
time for the sphere to cool to $30\,\degC$.
\paragraph{Solution sketch.} The characteristic length is $L_c = V/A_s
= (D/6) = 0.020/6 = 3.33 \times 10^{-3}\,\text{m}$ for a sphere. The
Biot number is $\text{Bi} = h L_c/k = 500 \times 3.33\times10^{-3}/400
= 4.2 \times 10^{-3}$, far below $0.1$, so the lumped approximation is
excellent. The time constant is $\tau = \rho V c_p/(h A_s) = \rho L_c
c_p/h = 8960 \times 3.33\times10^{-3} \times 385/500 = 23.0\,\text{s}$.
The cooling follows $(T - T_\infty)/(T_0 - T_\infty) = e^{-t/\tau}$, so
$(30 - 20)/(90 - 20) = 0.143 = e^{-t/\tau}$, giving $t = \tau\ln(1/0.143)
= 23.0 \times 1.946 = 45\,\text{s}$.
\end{exercise}

\begin{exercise}
A heated horizontal plate of $0.5\,\text{m} \times 0.5\,\text{m}$ at
$80\,\degC$ faces upward in $20\,\degC$ air. Estimate the free-
convection heat-transfer coefficient and the heat-loss rate. Use the
correlation $\text{Nu}_L = 0.54\,\text{Ra}_L^{1/4}$.
\paragraph{Solution sketch.} For a horizontal plate the characteristic
length is $L_c = A/P = (0.5 \times 0.5)/(4 \times 0.5) = 0.125\,\text{m}$.
For air near the film temperature $50\,\degC$, take $\nu = 1.8\times
10^{-5}\,\text{m}^2/\text{s}$, $\alpha = 2.5\times10^{-5}\,\text{m}^2/
\text{s}$, $k = 0.028\,\text{W}/(\text{m}\cdot\si{\kelvin})$,
$\text{Pr} = 0.72$, $\beta = 1/323\,\si{\kelvin}^{-1}$. Then
$\text{Ra}_L = g\beta\,\Delta T\,L_c^3/(\nu\alpha) = 9.81 \times
(1/323) \times 60 \times 0.125^3/(1.8\times10^{-5}\times2.5\times
10^{-5}) \approx 7.9 \times 10^6$, inside the $0.54$-coefficient range.
So $\text{Nu}_L = 0.54(7.9\times10^6)^{1/4} = 0.54 \times 53 = 28.6$,
and $h = \text{Nu}\,k/L_c = 28.6 \times 0.028/0.125 = 6.4\,\text{W}/
(\text{m}^2\cdot\si{\kelvin})$. The heat-loss rate is $\dot{Q} = hA\,
\Delta T = 6.4 \times 0.25 \times 60 = 96\,\text{W}$ from the top face;
radiation from the warm plate would add a comparable amount and is the
subject of the next exercise's machinery.
\end{exercise}

\begin{exercise}
A $20\,\degC$ room has walls at $15\,\degC$ (a cold winter day). A
person at $32\,\degC$ skin temperature with $1.7\,\text{m}^2$ of
exposed body area and $\epsilon = 0.95$ stands in the room. Compute
the net radiative heat-loss rate.
\paragraph{Solution sketch.} The body is a small convex object in a
large enclosure, so $F = 1$ and $\dot{Q} = \epsilon A \sigma (T_s^4 -
T_{\text{wall}}^4)$ with absolute temperatures. Skin at $32\,\degC =
305\,\si{\kelvin}$, walls at $15\,\degC = 288\,\si{\kelvin}$:
$\dot{Q} = 0.95 \times 1.7 \times 5.67\times10^{-8} \times (305^4 -
288^4)$. With $305^4 = 8.65\times10^9$ and $288^4 = 6.88\times10^9$,
the difference is $1.77\times10^9$, so $\dot{Q} = 0.95 \times 1.7
\times 5.67\times10^{-8} \times 1.77\times10^9 \approx 162\,\text{W}$.
This is the dominant heat-loss mode for a still, lightly clothed person
and explains why a room with cold walls feels cold even at a comfortable
air temperature; \texttt{radiation\_exchange.py} returns the same value.
\end{exercise}

\begin{exercise}
A flat double-pane window of total area $2\,\text{m}^2$ has two
$4\,\si{\milli\meter}$ glass panes separated by a $12\,\si{\milli\meter}$
air gap. Compute the U-value, ignoring the frame and treating the
air gap as still air with negligible radiation between the panes.
Then compute the U-value including the radiation across the gap (both
panes at emissivity $0.9$).
\paragraph{Solution sketch.} Per unit area, take outdoor and indoor
films $h_o = 25$, $h_i = 8\,\text{W}/(\text{m}^2\cdot\si{\kelvin})$.
The glass panes contribute $2 \times 0.004/1.0 = 0.008$, negligible.
The still-air gap alone gives $R_{\text{gap}} = 0.012/0.025 = 0.48\,
\si{\kelvin}\cdot\text{m}^2/\text{W}$. The series total is $R = 1/25 +
0.008 + 0.48 + 1/8 = 0.653$, so $U \approx 1.53\,\text{W}/(\text{m}^2
\cdot\si{\kelvin})$. Now add radiation across the gap in parallel with
the gas conduction. With both panes near $283\,\si{\kelvin}$ and an
effective emissivity $1/(1/0.9 + 1/0.9 - 1) = 0.82$, the linearised
radiation coefficient is $h_{\text{rad}} = 0.82 \times 4 \times
5.67\times10^{-8} \times 283^3 \approx 4.2\,\text{W}/(\text{m}^2\cdot
\si{\kelvin})$. The gas-conduction coefficient across the gap is
$0.025/0.012 = 2.1$, so the combined gap coefficient is $2.1 + 4.2 =
6.3$, giving $R_{\text{gap}} = 0.16$ instead of $0.48$. The series
total falls to $R = 0.04 + 0.008 + 0.16 + 0.125 = 0.333$ and $U
\approx 3.0\,\text{W}/(\text{m}^2\cdot\si{\kelvin})$. Radiation roughly
doubles the gap conductance, which is why low-emissivity coatings (which
cut $h_{\text{rad}}$) are the main lever in modern glazing.
\end{exercise}

\begin{exercise}
A long fin of length $50\,\si{\milli\meter}$, cross-section $1\,\si{
\milli\meter}\times 10\,\si{\milli\meter}$, attached to a surface at
$80\,\degC$ and exposed to air at $20\,\degC$ with $h = 20\,\text{W}/
(\text{m}^2\cdot\si{\kelvin})$, made of aluminium ($k = 237\,\text{W}/
(\text{m}\cdot\si{\kelvin})$). Find the fin temperature at the tip,
the heat dissipated by the fin, and the fin efficiency.
\paragraph{Solution sketch.} The cross-section is $A_c = 0.001 \times
0.010 = 1.0\times10^{-5}\,\text{m}^2$ and the perimeter $P = 2(0.001 +
0.010) = 0.022\,\text{m}$. The fin parameter is $m = \sqrt{hP/(kA_c)} =
\sqrt{20 \times 0.022/(237 \times 1.0\times10^{-5})} = \sqrt{185.7} =
13.6\,\text{m}^{-1}$, so $mL = 13.6 \times 0.050 = 0.68$. With an
insulated tip, the tip excess is $\theta_b/\cosh(mL) = 60/\cosh(0.68) =
60/1.24 = 48.3\,\si{\kelvin}$, giving a tip temperature of $20 + 48.3 =
68.3\,\degC$. The heat rate is $\dot{Q} = \sqrt{hPkA_c}\,\theta_b
\tanh(mL) = \sqrt{20 \times 0.022 \times 237 \times 1.0\times10^{-5}}
\times 60 \times \tanh(0.68)$. The square root is $\sqrt{1.043\times
10^{-3}} = 0.0323$, and $\tanh(0.68) = 0.591$, so $\dot{Q} = 0.0323
\times 60 \times 0.591 = 1.15\,\text{W}$. The efficiency is
$\tanh(mL)/(mL) = 0.591/0.68 = 0.87$, a stubby effective fin.
\texttt{fin\_analysis.py} returns these values.
\end{exercise}

\begin{exercise}
A horizontal pipe of outer diameter $50\,\si{\milli\meter}$ carries
saturated steam at $150\,\degC$ ($T_s = 150\,\degC$). It is bare and
in $20\,\degC$ air. Compute the free-convection coefficient using the
Churchill-Chu correlation, the radiation coefficient (assume
$\epsilon = 0.85$), and the total heat-loss rate per metre.
\paragraph{Solution sketch.} The film temperature is $85\,\degC$; for
air take $\nu = 2.1\times10^{-5}\,\text{m}^2/\text{s}$, $\alpha =
3.0\times10^{-5}$, $k = 0.030\,\text{W}/(\text{m}\cdot\si{\kelvin})$,
$\text{Pr} = 0.70$, $\beta = 1/358\,\si{\kelvin}^{-1}$. With $D =
0.050\,\text{m}$ and $\Delta T = 130\,\si{\kelvin}$, $\text{Ra}_D =
g\beta\,\Delta T\,D^3/(\nu\alpha) = 9.81 \times (1/358) \times 130
\times 0.050^3/(2.1\times10^{-5}\times3.0\times10^{-5}) \approx 7.1
\times10^5$. Churchill-Chu gives $\text{Nu}_D \approx \{0.6 + 0.387
(7.1\times10^5)^{1/6}/[1+(0.559/0.70)^{9/16}]^{8/27}\}^2 \approx 13.5$,
so $h_{\text{conv}} = \text{Nu}\,k/D = 13.5 \times 0.030/0.050 = 8.1\,
\text{W}/(\text{m}^2\cdot\si{\kelvin})$. The radiation coefficient is
$h_{\text{rad}} = \epsilon\sigma(T_s^2+T_\infty^2)(T_s+T_\infty) = 0.85
\times 5.67\times10^{-8}(423^2+293^2)(423+293) \approx 9.1\,\text{W}/
(\text{m}^2\cdot\si{\kelvin})$. The two are comparable; the combined
coefficient is $17.2$, and per metre $\dot{Q}/L = (h_{\text{conv}} +
h_{\text{rad}})\pi D\,\Delta T = 17.2 \times \pi \times 0.050 \times
130 \approx 352\,\text{W}/\text{m}$. Radiation carries roughly half the
loss from a bare hot pipe, the standard warning against ignoring it.
\end{exercise}

\subsection*{Derivation}\pathtag{standard}

\begin{exercise}
Derive the steady-state radial heat-conduction resistance $R = \ln
(r_2/r_1)/(2\pi k L)$ for a cylindrical wall by direct integration of
Fourier's law in cylindrical coordinates.
\paragraph{Solution sketch.} At steady state with no generation, the
radial heat-flow rate $\dot{Q}$ is the same through every cylindrical
shell. Fourier's law on a shell of radius $r$ and area $A = 2\pi r L$
reads $\dot{Q} = -k(2\pi r L)\,dT/dr$. Separating, $dT = -\dot{Q}/(2\pi
k L)\,dr/r$, and integrating from $r_1$ (at $T_1$) to $r_2$ (at $T_2$),
$T_2 - T_1 = -\dot{Q}/(2\pi k L)\ln(r_2/r_1)$. Rearranging, $\dot{Q} =
2\pi k L(T_1 - T_2)/\ln(r_2/r_1) = (T_1 - T_2)/R$ with $R = \ln(r_2/r_1)
/(2\pi k L)$. The logarithm appears because the area grows linearly with
radius, unlike the plane wall where the area is constant.
\end{exercise}

\begin{exercise}
Derive the critical-insulation-radius formula $r_c = k/h$ from the
condition that adding insulation reduces (rather than increases) the
heat-loss rate.
\paragraph{Solution sketch.} For a cylinder with insulation outer radius
$r$, the two outer resistances per unit length are the insulation
$R_{\text{ins}} = \ln(r/r_i)/(2\pi k)$ and the outer film $R_o = 1/(2\pi
r h)$. Their sum is $R(r) = \ln(r/r_i)/(2\pi k) + 1/(2\pi r h)$. Heat
loss falls when $R$ rises, so set $dR/dr = 0$: $dR/dr = 1/(2\pi k r) -
1/(2\pi h r^2) = 0$ gives $r = k/h$. This stationary point is a minimum
of $R$ (a maximum of heat loss): for $r < k/h$ the $1/r^2$ film term
dominates and $R$ falls as $r$ grows, while for $r > k/h$ the conduction
term dominates and $R$ rises. So $r_c = k/h$ is the radius below which
adding insulation increases the loss; only once the outer radius exceeds
$r_c$ does insulation help. For a fat pipe $r_i > r_c$ already and
insulation always helps; for a thin wire the opposite can hold.
\end{exercise}

\begin{exercise}
Derive the lumped-capacitance time constant $\tau = \rho V c_p/(hA)$
from the energy balance $\rho V c_p\, dT/dt = -hA(T - T_\infty)$ and
the assumption of spatial uniformity.
\paragraph{Solution sketch.} Spatial uniformity means the body has one
temperature $T(t)$. Energy conservation sets the rate of change of
stored energy equal to the surface heat loss: $\rho V c_p\,dT/dt =
-hA(T - T_\infty)$. Let $\theta = T - T_\infty$, so $d\theta/dt =
-(hA/\rho V c_p)\theta$. This is a first-order linear equation with
solution $\theta(t) = \theta_0 e^{-t/\tau}$, where the time constant is
the reciprocal of the bracketed coefficient, $\tau = \rho V c_p/(hA)$.
The numerator $\rho V c_p$ is the body's thermal capacitance and the
denominator $hA$ its surface conductance, so $\tau$ is capacitance over
conductance, the thermal analogue of the $RC$ time constant of an
electrical circuit.
\end{exercise}

\begin{exercise}
Derive the linearised radiation coefficient $h_{\text{rad}} \approx
4 \epsilon \sigma T_{\text{avg}}^3$ from the Stefan-Boltzmann law
applied to small temperature differences.
\paragraph{Solution sketch.} The net radiative exchange for a small body
in an enclosure is $\dot{Q} = \epsilon\sigma A(T_1^4 - T_2^4)$. Factor
the difference of fourth powers: $T_1^4 - T_2^4 = (T_1^2 + T_2^2)(T_1 +
T_2)(T_1 - T_2)$. For $T_1 \approx T_2 \approx T_{\text{avg}}$, the first
two factors become $2T_{\text{avg}}^2$ and $2T_{\text{avg}}$, so $T_1^4
- T_2^4 \approx 4T_{\text{avg}}^3(T_1 - T_2)$. Substituting, $\dot{Q}
\approx 4\epsilon\sigma T_{\text{avg}}^3 A(T_1 - T_2) = h_{\text{rad}} A
(T_1 - T_2)$ with $h_{\text{rad}} = 4\epsilon\sigma T_{\text{avg}}^3$.
The same factorisation follows from a first-order Taylor expansion of
$T^4$ about $T_{\text{avg}}$. The linear form lets radiation be combined
with convection as a parallel resistance, valid when $|T_1 - T_2| \ll
T_{\text{avg}}$.
\end{exercise}

\begin{exercise}
Show that the Biot number $\text{Bi} = hL/k$ is the ratio of internal
conduction resistance to external convection resistance, and that
the lumped-capacitance approximation applies when $\text{Bi} \ll 1$.
\paragraph{Solution sketch.} Compare the two resistances in series
between the body's core and the ambient fluid. The internal conduction
resistance over the characteristic length is $R_{\text{cond}} = L/(k A)$,
and the surface convection resistance is $R_{\text{conv}} = 1/(h A)$.
Their ratio is $R_{\text{cond}}/R_{\text{conv}} = [L/(kA)]/[1/(hA)] =
hL/k = \text{Bi}$. So the Biot number is exactly the ratio of internal
to external resistance. When $\text{Bi} \ll 1$ the internal resistance is
small compared with the surface resistance, the temperature drop within
the body is negligible against the drop across the film, and the body is
nearly isothermal, which is the premise of the lumped-capacitance model.
When $\text{Bi}$ is large the dominant drop is inside the body, spatial
gradients matter, and the full heat equation is needed.
\end{exercise}

\subsection*{Estimation}\pathtag{core}

\begin{exercise}
Estimate the rate of heat loss from a person standing outdoors in
$0\,\degC$ still air in summer clothing. State the assumed body
surface area, clothing insulation value, and metabolic heat
generation.
\paragraph{Reference estimate.} Take a body surface area of $1.8\,
\text{m}^2$, light summer clothing at about $0.5\,\text{clo}$
($1\,\text{clo} = 0.155\,\text{m}^2\cdot\si{\kelvin}/\text{W}$, so
$R_{\text{cl}} = 0.078$), and a skin temperature near $33\,\degC$. The
outer surface loses heat by combined convection and radiation at roughly
$h \approx 12\,\text{W}/(\text{m}^2\cdot\si{\kelvin})$ in still air, so
$R_{\text{surf}} = 1/12 = 0.083\,\text{m}^2\cdot\si{\kelvin}/\text{W}$.
The series resistance per unit area is $0.078 + 0.083 = 0.161$, and the
heat-loss flux is $\Delta T/R = 33/0.161 = 205\,\text{W}/\text{m}^2$,
giving $\dot{Q} \approx 205 \times 1.8 \approx 370\,\text{W}$. Resting
metabolism supplies only about $100\,\text{W}$, so the person loses heat
roughly three to four times faster than the body generates it: the body
shivers and the person feels cold within minutes, consistent with the
real experience of standing outdoors at freezing in summer clothes.
\end{exercise}

\begin{exercise}
Estimate the temperature reached by a black-painted metal roof on a
summer day with $1000\,\text{W}/\text{m}^2$ of solar irradiance,
$30\,\degC$ ambient air, and natural-convection cooling. State the
emissivity and absorptivity assumed.
\paragraph{Reference estimate.} Black paint absorbs most of the solar
flux, so take solar absorptivity $\alpha_s = 0.95$ and thermal
emissivity $\epsilon = 0.9$. The absorbed flux is $0.95 \times 1000 =
950\,\text{W}/\text{m}^2$. At steady state this leaves by convection and
radiation. Convection takes $h_{\text{conv}}(T_s - T_\infty)$ with
$h_{\text{conv}} \approx 10\,\text{W}/(\text{m}^2\cdot\si{\kelvin})$ for
a warm horizontal surface in light air; radiation takes $\epsilon\sigma
(T_s^4 - T_{\text{sky}}^4)$, and with a sky near ambient we approximate
it as $h_{\text{rad}}(T_s - T_\infty)$ with $h_{\text{rad}} \approx 4
\times 0.9 \times 5.67\times10^{-8} \times 330^3 \approx 7.3$. The
combined coefficient is about $17\,\text{W}/(\text{m}^2\cdot
\si{\kelvin})$. Balancing, $950 = 17(T_s - 303)$, so $T_s - 303 \approx
56\,\si{\kelvin}$ and $T_s \approx 359\,\si{\kelvin} = 86\,\degC$. A
black roof in full sun runs $50$ to $60\,\si{\kelvin}$ above air, the
familiar reason dark roofs are hot enough to burn skin and why
high-reflectance ``cool roofs'' cut summer cooling loads.
\end{exercise}

\begin{exercise}
Estimate the time for a cup of tea ($250\,\text{mL}$, $80\,\degC$
initial) to cool to drinkable temperature ($55\,\degC$) on a counter
in a $20\,\degC$ kitchen. State the heat-transfer mechanisms
included.
\paragraph{Reference estimate.} An open mug of $250\,\text{mL}$ presents
a top surface of about $0.005\,\text{m}^2$ (a $40\,\si{\milli\meter}$
radius) plus side walls of roughly $0.015\,\text{m}^2$. Heat leaves by
evaporation from the surface, by convection and radiation from the
liquid and walls. Evaporation is the largest term for an open hot
drink: a wet surface at $80\,\degC$ loses on the order of $400$ to
$800\,\text{W}/\text{m}^2$, say $600$, so $\approx 3\,\text{W}$ from the
top. Convection and radiation from the top and sides at a combined
$h \approx 15\,\text{W}/(\text{m}^2\cdot\si{\kelvin})$ over $0.02\,
\text{m}^2$ with $\Delta T \approx 60$ add $15 \times 0.02 \times 60
\approx 18\,\text{W}$, dropping as the drink cools. Take an average
loss near $15\,\text{W}$. The thermal mass is $mc = 0.25 \times 4180 =
1045\,\text{J}/\si{\kelvin}$, so the time constant is $\tau = mc/(\dot{Q}
/\Delta T) \approx 1045/(15/60) = 4200\,\text{s}$. Cooling from $80$ to
$55\,\degC$ means the excess falls from $60$ to $35\,\si{\kelvin}$, a
factor of $0.58$, so $t = \tau\ln(60/35) \approx 4200 \times 0.54
\approx 2300\,\text{s} \approx 38\,\text{min}$. The order of magnitude,
tens of minutes, matches everyday experience; a lid that suppresses
evaporation roughly doubles the time.
\end{exercise}

\begin{exercise}
Estimate the U-value of a typical $100\,\si{\milli\meter}$ thick
poured-concrete basement wall (treating the soil outside as a
semi-infinite medium at $10\,\degC$ at the soil's deep-ground
temperature).
\paragraph{Reference estimate.} The concrete itself is a poor barrier:
$R_{\text{conc}} = 0.10/1.7 = 0.059\,\si{\kelvin}\cdot\text{m}^2/\text{W}$.
The dominant resistance is the soil path. A buried wall does not see a
thin film; heat spreads three-dimensionally into the ground, which acts
as a thick conducting medium. Treating the soil as a semi-infinite
medium of conductivity $k_{\text{soil}} \approx 1.5\,\text{W}/(\text{m}
\cdot\si{\kelvin})$, the effective path length to the deep-ground
temperature is on the order of metres rather than the wall thickness;
an effective soil resistance of $1$ to $2\,\si{\kelvin}\cdot\text{m}^2/
\text{W}$ is representative for a wall a metre or two below grade. Adding
an interior film $1/8 = 0.125$, the series total is about $1.7\,\si{
\kelvin}\cdot\text{m}^2/\text{W}$, giving $U \approx 0.6\,\text{W}/
(\text{m}^2\cdot\si{\kelvin})$. The soil coupling, not the concrete, sets
the value, which is why basement insulation is placed on the wall face
and why deep-ground temperature (near $10\,\degC$ year-round) makes
basements milder than above-grade rooms in both seasons.
\end{exercise}

\subsection*{Simulation}\pathtag{mastery}

\begin{exercise}
Implement a finite-difference solver for the 1D heat equation in a
slab. Compute the temperature evolution after a step change in one
surface temperature; compare to the semi-infinite-solid analytical
solution at short times.
\paragraph{Model-answer sketch.} Discretise the slab into $N$ nodes with
spacing $\Delta x$ and march the explicit forward-time centred-space
scheme $T_i^{n+1} = T_i^n + \text{Fo}_{\Delta}(T_{i+1}^n - 2T_i^n +
T_{i-1}^n)$, where the mesh Fourier number $\text{Fo}_{\Delta} = \alpha
\Delta t/\Delta x^2$ must stay at or below $1/2$ for stability; the
reference implementation \texttt{heat\_equation\_1d.py} uses
$\text{Fo}_{\Delta} = 0.4$ and raises an error otherwise. Fix one face
to the stepped temperature and hold the far face at the initial value.
At short times, before the disturbance reaches the far face, the slab
behaves as a semi-infinite solid and the node temperatures should match
$T(x,t) = T_s + (T_0 - T_s)\,\text{erf}(x/2\sqrt{\alpha t})$ to within
the spatial truncation error. The reference run on a concrete slab
($\alpha = 5\times10^{-7}\,\text{m}^2/\text{s}$) matches the error
function to two decimals at every node until the penetration depth
$\sqrt{\alpha t}$ approaches the slab thickness, where the finite far
boundary makes the two diverge. Halving $\Delta x$ should roughly halve
the discrepancy in the unaffected region, confirming first-to-second
order convergence.
\end{exercise}

\begin{exercise}
Implement a 2D finite-difference solver for steady conduction in a
rectangular domain with mixed boundary conditions (fixed temperature
on some boundaries, convection on others). Verify on a problem with
a known analytical solution.
\paragraph{Model-answer sketch.} On a rectangular grid the steady heat
equation $\nabla^2 T = 0$ discretises to the five-point Laplacian
$T_{i,j} = \tfrac14(T_{i+1,j} + T_{i-1,j} + T_{i,j+1} + T_{i,j-1})$ at
interior nodes. Fixed-temperature boundaries set $T$ directly; a
convective boundary uses the discretised flux balance $-k(T_{\text{int}}
- T_{\text{bndry}})/\Delta x = h(T_{\text{bndry}} - T_\infty)$, which
gives a ghost-node or one-sided update for the boundary value in terms of
the interior neighbour, the Biot number $h\Delta x/k$, and $T_\infty$.
Iterate with Gauss-Seidel or successive over-relaxation until the maximum
node change falls below a tolerance. Verify on a square with three faces
held at $0$ and one at $T_0$, whose analytical solution is the Fourier
series $T(x,y) = T_0\sum_n b_n \sin(n\pi x/L)\sinh(n\pi y/L)/\sinh(n\pi)$;
the numerical field should converge to this as the grid is refined, with
the maximum error falling roughly as $\Delta x^2$.
\end{exercise}

\begin{exercise}
Simulate the thermal response of a household-wall sandwich to a
$24\,\text{h}$ outdoor temperature swing of $\pm 10\,\si{\kelvin}$.
Identify the time delay and the amplitude reduction of the inner
surface temperature.
\paragraph{Model-answer sketch.} Drive the outer face of the wall
sandwich with $T_o(t) = \bar{T}_o + 10\sin(2\pi t/24\text{h})$ and march
the one-dimensional transient solver (the \texttt{heat\_equation\_1d.py}
scheme extended to layered conductivity and capacitance) until the inner
surface reaches a periodic steady state. The inner surface temperature
oscillates at the same $24\,\text{h}$ period but lagged and attenuated.
For a sinusoidal surface forcing on a diffusive medium the theory
predicts an amplitude that decays as $e^{-d/\delta}$ and a phase lag of
$d/\delta$ radians, with the diurnal penetration depth $\delta =
\sqrt{2\alpha/\omega}$ and $\omega = 2\pi/(24\text{h})$. For a typical
insulated wall the inner-surface swing should come out a small fraction
of the $\pm 10\,\si{\kelvin}$ outdoor swing, perhaps $\pm 1$ to $2\,
\si{\kelvin}$, with a lag of several hours, which is the thermal-mass
effect that keeps massive buildings cool through a hot afternoon. The
deliverable reports the measured amplitude ratio and lag and compares
them to the $e^{-d/\delta}$ and $d/\delta$ predictions.
\end{exercise}

\subsection*{Design}\pathtag{standard}

\begin{exercise}
Specify the insulation thickness for a domestic hot-water tank
($60\,\degC$, $200\,\text{L}$) such that the standby heat loss is
not more than $50\,\text{W}$ at $20\,\degC$ ambient. State the
insulation material and the tank geometry assumed.
\paragraph{Solution sketch.} Model the $200\,\text{L}$ tank as a
cylinder roughly $0.5\,\text{m}$ in diameter and $1.0\,\text{m}$ tall,
surface area $A \approx \pi D h + 2(\pi D^2/4) \approx 1.57 + 0.39
\approx 2.0\,\text{m}^2$. The driving difference is $60 - 20 = 40\,
\si{\kelvin}$, so the allowable overall conductance is $UA \leq 50/40 =
1.25\,\text{W}/\si{\kelvin}$, i.e. $U \leq 0.63\,\text{W}/(\text{m}^2
\cdot\si{\kelvin})$. The outer film and tank wall contribute little, so
the insulation must supply almost all of $R = 1/U = 1.6\,\si{\kelvin}
\cdot\text{m}^2/\text{W}$. With polyurethane foam at $k = 0.025\,
\text{W}/(\text{m}\cdot\si{\kelvin})$, the thickness is $t = kR = 0.025
\times 1.6 = 0.040\,\text{m}$, about $40\,\si{\milli\meter}$. Mineral
wool at $k = 0.04$ would need $\approx 64\,\si{\milli\meter}$. A
$40$ to $50\,\si{\milli\meter}$ foam jacket meets the standby-loss
target with margin; the answer should note that the cylindrical critical
radius is far exceeded, so the insulation unambiguously helps.
\end{exercise}

\begin{exercise}
Design a heat sink for a power semiconductor dissipating $100\,\text{W}$,
to limit the junction temperature to $100\,\degC$ in $40\,\degC$
ambient air. Specify the fin geometry, the airflow (forced or
natural), and the thermal-interface material.
\paragraph{Solution sketch.} The total allowable junction-to-ambient
resistance is $R_{\text{tot}} = (100 - 40)/100 = 0.6\,\si{\kelvin}/
\text{W}$. Budget this across the stack: junction-to-case $\approx 0.2$
(a device datasheet number), a thermal-interface layer at $\approx 0.05$
(use thermal grease or a phase-change pad, $R_{\text{TIM}} = t/(k A)$ kept
thin and high-$k$), leaving a heat-sink-to-ambient budget of $\approx
0.35\,\si{\kelvin}/\text{W}$. Natural convection cannot reach this in a
reasonable size (a bare $100\,\text{W}$ sink would need square metres),
so specify forced air. A finned aluminium sink, say a $60\,\si{
\milli\meter}\times 60\,\si{\milli\meter}$ base with $20$ to $30$ fins of
$30\,\si{\milli\meter}$ height at $5\,\text{m/s}$, delivers $0.2$ to
$0.4\,\si{\kelvin}/\text{W}$ from fin-array correlations, meeting the
budget. The design should check the fin parameter $mL$ stays near unity
so the fins run effective (Figure~\ref{fig:vol04ch05:fin-temperature})
and confirm the case temperature $100 - 100\times0.25 = 75\,\degC$
leaves margin to any case rating.
\end{exercise}

\begin{exercise}
Specify the thermal protection (insulation, fire-rated cladding) for
a $1000\,\text{L}$ propane storage tank such that, under a hydrocarbon
pool-fire exposure, the tank's working pressure does not exceed its
rated value within a $30\,\text{min}$ window. State the heat-flux
assumption.
\paragraph{Solution sketch.} A standard hydrocarbon pool-fire exposure
is taken at an incident flux near $100\,\text{kW}/\text{m}^2$ (the value
used in tank fire-test standards), radiating at a flame temperature
around $1000$ to $1100\,\degC$. The protection goal is to keep the tank
wall and contents below the temperature at which the relief valve cannot
hold working pressure, roughly limiting the wall to a few hundred
degrees over $30\,\text{min}$. The insulation must throttle the absorbed
flux: with an outer face driven toward flame temperature and an inner
face held near the contents, a layer of intumescent or ceramic-fibre
cladding of conductivity $k \approx 0.1\,\text{W}/(\text{m}\cdot
\si{\kelvin})$ and thickness $t$ passes $q = k\,\Delta T/t$. Sizing $t$
so the conducted flux times the exposed area, integrated over
$30\,\text{min}$, raises the tank contents and wall by less than the
allowable temperature rise (using the tank thermal mass $mc$) sets the
thickness; a $25$ to $50\,\si{\milli\meter}$ rated fireproofing layer is
typical. The answer should state the assumed pool-fire flux, the
allowable wall temperature, and the transient energy balance $mc\,
dT/dt = q_{\text{in}} A - q_{\text{relief}}$ that the insulation makes
survivable for the required window.
\end{exercise}

\subsection*{Diagnosis and reverse engineering}\pathtag{core}

\begin{exercise}
A residential wall is observed (by an infrared camera) to have
substantial cold spots at the corners and along the wall-to-floor
junction. Identify the heat-transfer mechanism and the construction
defect responsible.
\paragraph{Solution sketch.} Cold spots on the interior surface mark
places where the local heat flux outward is high, so the inner face runs
cold. The mechanism is thermal bridging: a high-conductivity path that
bypasses the insulation, drawn as the parallel $R_{\text{stud}}$ branch
in Figure~\ref{fig:vol04ch05:thermal-resistance-circuit}. At corners,
geometry concentrates the heat flow (two walls meeting present more
outside area than inside), and the wall-to-floor junction usually carries
a continuous framing or slab edge that conducts around the insulation.
The construction defects responsible are uninsulated framing members
(studs, plates, the slab edge), missing insulation at the geometric
corners, and sometimes air leakage washing the inner surface. The fix is
continuous exterior insulation that covers the framing and a thermal
break at the slab edge. The diagnosis confirms the bridge by comparing
the cold-spot surface temperature to the clear-field surface temperature:
the colder the stripe, the lower its local resistance.
\end{exercise}

\begin{exercise}
A datacenter has experienced a $5\,\degC$ rise in average rack inlet
temperature over the past six months. Identify the most probable
causes from among (a) increased computational load, (b) fouling of
the air-handler coils, (c) a chilled-water flow restriction, and
state the diagnostic measurement for each.
\paragraph{Solution sketch.} The rack inlet temperature rises when
generation outpaces removal, $\dot{Q}_{\text{gen}} = \dot{m}c_p\,\Delta
T_{\text{air}}$. (a) Increased computational load raises
$\dot{Q}_{\text{gen}}$ directly; diagnose by reading the facility power
draw (kW at the distribution units) and correlating it with the
temperature rise, since electrical power in equals heat out. (b) Fouled
air-handler coils add a resistance on the air side, so the same chilled
water removes less heat; diagnose by measuring the air-side temperature
drop across the coil and the air pressure drop, both degraded by
fouling, and by comparing supply-air temperature to chilled-water
temperature. (c) A chilled-water flow restriction cuts the water mass
flow, raising the water-side $\Delta T$ for the same duty; diagnose by
measuring the chilled-water flow rate and the supply-to-return water
temperature difference, which rises when flow falls. The disciplined
approach plots generation against removal and asks which side moved: a
power-draw rise points to (a), a coil-side resistance rise to (b), a
water-flow drop to (c).
\end{exercise}

\subsection*{Failure analysis}\pathtag{core}

\begin{exercise}
A lithium-ion battery pack in an electric vehicle catches fire
following a high-speed collision. Identify the thermal-runaway
sequence and the design features (cell-level and pack-level) that
should have limited propagation.
\paragraph{Solution sketch.} The collision creates an internal short in
one or more cells (crushed separator, breached can), dumping the cell's
stored energy as local heat. The cell temperature climbs through the
side-reaction onset near $80\,\degC$, where exothermic reactions add
heat with rates that roughly double every $10\,\si{\kelvin}$ (Arrhenius);
above $\sim 130\,\degC$ the SEI layer decomposes; above $\sim 200\,\degC$
the electrolyte breaks down and the cathode releases oxygen, which
sustains combustion. This is the positive-feedback runaway of the
chapter's failure section, where generation outruns removal. Cell-level
features that should limit it: a shutdown separator that opens the
circuit near $130\,\degC$, a current-interrupt device, and a pressure
vent to release gas before rupture. Pack-level features: thermal barriers
and spacing between cells so a single runaway cannot conduct or radiate
its neighbours over their onset temperature, a cooling system sized to
remove the worst-case heat, and a battery-management system that detects
the abnormal cell and isolates it. The failure analysis asks whether the
propagation path between cells had enough thermal resistance to break the
chain; if neighbouring cells reached onset, the inter-cell isolation was
inadequate.
\end{exercise}

\begin{exercise}
Re-read the Kemeny Commission's report on Three Mile Island
\cite{acc:kemeny-tmi-1979}. Identify the heat-transfer failure mode
(loss of coolant inventory, decay-heat removal failure) and state
how the design's defence-in-depth was breached.
\paragraph{Solution sketch.} The heat-transfer failure mode at Three
Mile Island is decay-heat removal failure following a loss of coolant,
not a thermal runaway. After the reactor scrammed, fission stopped but
decay heat persisted at several percent of full power, falling to about
$1\,\%$ at an hour and $0.5\,\%$ at a day, and this heat had to be
removed continuously. A pressure-operated relief valve stuck open,
draining coolant inventory; the operators, misreading the indications,
throttled the emergency injection that would have replaced it. The core
uncovered, and the residual decay heat with no liquid coolant overheated
and damaged the fuel. The defence-in-depth was breached at several
layers: the stuck valve defeated the pressure-control barrier, the
instrumentation did not make the loss of inventory plain, and the
operator action removed the makeup-water barrier that should have caught
the first failure. The lesson is that a shut-down reactor still demands
an independent decay-heat path for days; the registry entry
\cite{acc:kemeny-tmi-1979} records the sequence.
\end{exercise}

\begin{exercise}
Re-read the IAEA INSAG-7 report on Chernobyl
\cite{acc:iaea-insag7-1992}. Identify the positive-feedback mechanism
(positive void coefficient, scram control-rod insertion) that
produced the thermal-runaway power excursion.
\paragraph{Solution sketch.} The positive-feedback mechanism is the
positive void coefficient of the graphite-moderated, water-cooled RBMK
design. When the cooling water boiled to steam (voided), it absorbed
fewer neutrons while the graphite kept moderating, so reactivity rose,
which raised power, which made more steam: generation feeding back on
itself, the runaway condition $d\dot{q}_g/dT > d\dot{q}_r/dT$ of the
chapter. The excursion was triggered during a low-power test that left
the reactor in an unstable configuration; the attempted shutdown made it
worse, because the control rods carried graphite tips that briefly
displaced water and added reactivity on insertion before the absorber
followed. Power spiked far above rating in seconds, destroying the core.
The contrast with a negative-feedback reactor (where rising temperature
cuts reactivity through Doppler broadening and moderator density) is the
whole point: stable designs keep the net feedback negative across the
operating envelope. The registry entry \cite{acc:iaea-insag7-1992}
records the mechanism and the design and operational contributors.
\end{exercise}

\subsection*{Judgment}\pathtag{standard}

\begin{exercise}
A junior engineer argues that radiation can be neglected in any
heat-transfer calculation involving room-temperature surfaces. Write
a one-paragraph reply identifying when this assumption is and is not
acceptable.
\paragraph{Model-answer sketch.} The reply should reject the blanket
claim. At room temperature the linearised radiation coefficient is
$h_{\text{rad}} = 4\epsilon\sigma T_{\text{avg}}^3 \approx 5.5\,\text{W}/
(\text{m}^2\cdot\si{\kelvin})$ for a high-emissivity surface, which is
comparable to still-air free-convection coefficients of $2$ to $25$. So
radiation may be neglected only when the competing convection is much
larger, as in forced flow or boiling, or when the surface emissivity is
genuinely low (bright polished metal, $\epsilon \sim 0.05$). It must be
kept when convection is weak (still air, a person in a cold-walled room,
a low-velocity radiator, a clear-night surface radiating to the sky), or
the radiative term carries a comparable or larger share of the heat. The
discipline is to estimate $h_{\text{rad}}$ and compare it to $h_{\text{conv}}$
rather than assume either dominates.
\end{exercise}

\begin{exercise}
A facilities manager argues that adding insulation to a small-
diameter chilled-water pipe ($25\,\si{\milli\meter}$ outside diameter,
$5\,\degC$ contents, $20\,\degC$ ambient) will reduce heat gain. Write
a one-paragraph reply that addresses the critical-insulation-radius
issue and identifies whether the manager is correct in this case.
\paragraph{Model-answer sketch.} The reply checks the critical
insulation radius $r_c = k/h$ against the pipe outer radius of $12.5\,
\si{\milli\meter}$. For typical pipe insulation $k \approx 0.04\,\text{W}
/(\text{m}\cdot\si{\kelvin})$ and an indoor free-convection coefficient
$h \approx 8\,\text{W}/(\text{m}^2\cdot\si{\kelvin})$, $r_c = 0.04/8 =
5\,\si{\milli\meter}$, well below the bare outer radius. Since the outer
radius already exceeds $r_c$, adding insulation raises the total
resistance and cuts the heat gain, so the manager is correct here. The
reply should flag the caveat the manager glossed: on a much thinner line
or with a lower film coefficient, $r_c$ could exceed the bare radius and
the first increment of insulation would increase heat gain rather than
reduce it, which is why the check matters before generalising. A
chilled-water pipe behaves the same way as a hot pipe; only the sign of
$\Delta T$ differs, and condensation control is a second reason to
insulate it regardless.
\end{exercise}

\begin{exercise}
A senior colleague proposes to design a battery-pack thermal-runaway
mitigation by ``oversizing the cooling system.'' Write a one-paragraph
reply that argues for cell-level isolation as the complementary
defence, with reference to the propagation mechanism.
\paragraph{Model-answer sketch.} The reply should accept oversized
cooling as useful but insufficient on its own. Cooling addresses the
removal side of the heat balance under normal operation, but a runaway
begins when a single cell shorts and generates heat far faster than any
practical pack-level cooling can extract from that one cell; the
generation term spikes locally and outruns removal. The threat is then
propagation: the runaway cell conducts and radiates heat to its
neighbours, driving them over their side-reaction onset and cascading
through the pack. Cell-level isolation (thermal barriers, spacing,
intumescent layers, vent paths directed away from neighbours) attacks
the propagation path directly by adding thermal resistance between cells,
so one cell's failure cannot ignite the next. The two defences are
complementary: cooling holds the steady operating point well below the
unstable intersection, and isolation contains the failure when a cell
crosses it anyway. Relying on cooling alone leaves the propagation path
open.
\end{exercise}
